\section{The Galois link III: assembly}\label{sec:pf:linkN}

\providecommand{\linkNup}{M^{\mathrm{up}}}
\providecommand{\linkNlow}{M^{\mathrm{low}}}
\providecommand{\linkNeS}{\eta\Sigma}
\providecommand{\linkNloc}{\eta_{\mathrm{loc}}}
\providecommand{\linkNvcc}{v_{\mathrm{cc}}}
\providecommand{\linkNmult}{\mathrm{mult}}
\providecommand{\linkNgood}{Z^{\mathrm{good}}}
\providecommand{\linkNdZ}{d_{Z_0}}
\providecommand{\linkNlw}{\mathrm{lw}}
\providecommand{\linkNs}{s_\Pi}
\providecommand{\linkNfin}{\mathrm{Fin}_{\Sha}}

This section proves \cref{thm:linkN}, the Galois link. Part~(a) of that theorem says that a family of data all of
whose members satisfy the \emph{Galois preconditions} and the condition $\sfname{Tie4}$ satisfies the link
$\sfname{LinkN}$ of \cref{def:families:linkN}: the weighted divisibility defect $\linkNeS(L,j)$ at the twisted places is
bounded by $\linkNdZ(j)\bigl(C(fmb\Lambda^2+(\log x)^2N_c)+\sum_\ell\log\ell\sum_{0<e<t}\mathrm{sh}(\ell,e)\bigr)$.
Part~(b) transports this to the parameter sets of the construction of \cref{sec:families:nonempty}.

\paragraph{The argument.}
The local inequalities behind $\sfname{LinkN}$ have the shape \eqref{eq:families:bd}: the length of a congruence module at
a good eigensystem is bounded by the length of a Selmer group plus error terms. They are established in
\lref{GlobalTheory.NLinkClose.main_alt_alt} for every $\sfname{TwGood}$ family whose members carry the
\emph{norm-choice facts} and two further conditions with uniform constants. The present section supplies everything
around this input.
\begin{enumerate}[label=(\arabic*)]
\item \emph{Uniform constants} (\cref{sec:pf:linkN:extra}). With $C_T=40$ and $C=320$, every datum carrying the
norm-choice facts and $\sfname{Tie4}$ satisfies a commuting endomorphism condition and a paired character condition. The
endomorphism comes from central elements of a group ring whose images are $\#G$ times central idempotents
(\cref{pf:linkN:cenidem}); the paired character condition is reduced to local results of $p$-adic Hodge theory.
\item \emph{Cuts} (\cref{sec:pf:linkN:lattice,sec:pf:linkN:cuts,sec:pf:linkN:goodcuts}). The Galois preconditions do not
imply the norm-choice facts, but a \emph{good cut} of the datum does: one keeps the local indices having good
eigensystems, restricts the Hecke algebras to a subset of eigensystems, and replaces each Hecke module $M$ by its quotient
by the part lying in the eigenspaces of the discarded eigensystems. The vanishing of local Galois cohomology at
auxiliary places of Taylor--Wiles type (\cref{pf:linkN:tamevanish}) enters here.
\item \emph{Summation} (\cref{sec:pf:linkN:sum}). For a family satisfying the local bounds, the bounds are summed over the
twisted places with the weights $\log\Nm w/[\mathit{LL}:\Q]$, giving $\sfname{LinkN}$ with constant $2000(|C|+1)^3$.
\item \emph{Return from the cut} (\cref{sec:pf:linkN:cuts}). A lattice lemma (\cref{pf:linkN:lengthle}) shows that the
$\eta$-invariants of a datum exceed those of its good cut by at most $\linkNdZ(j)\,v_\ell(\#G)$ at each twisted place
above $\ell$, and $\log\#G$ is absorbed by the main term of $\sfname{LinkN}$. This proves part~(a).
\item \emph{Rows} (\cref{sec:pf:linkN:b}). Part~(b) follows formally from part~(a), a transport of properties of the
members of layer points to all linked bundles, and two theorems saying that the Galois preconditions and $\sfname{Tie4}$
hold at all members.
\end{enumerate}

\paragraph{What is proved here and what is cited.}
The local bounds are proved in \lref{GlobalTheory.NLinkClose.main_alt_alt}, and the bound at the primes $2$ and $3$ is
\cref{lem:families:at23}. Several local results used for the uniform constants (the lifting theorem for crystalline lattices of tame
type behind the crystalline period bound of \cref{sec:pf:linkN:thmT}, a local pushout bound, congruences at the type primes and a
sixteen-step statement derived from Artin reciprocity), the estimates for the sums over primes in \cref{pf:linkN:pointbound}, and the
large theorems asserting that the Galois preconditions and $\sfname{Tie4}$ hold at members, are cited with their statements. The
algebraic and cohomological statements used in the proof of the local bounds are proved in \cref{sec:pf:linkN:congmod,sec:pf:linkN:cocycles}.
Peculiarities of the formal statements are collected in \cref{sec:pf:linkN:remarks}.

% ---------------------------------------------------------------------------
\subsection{Data, invariants and hypotheses}\label{sec:pf:linkN:data}

We use the notation of \cref{sec:families}, in particular of \cref{sec:families:galois}, and recall the components of a
datum $D$ that occur below.

\paragraph{Ranks, host and rows.}
$D$ has rank parameters $k,m,t,f,n,r,n_V$ with $2\le m\le f$; a host with number fields $F\subset E$, $[E:\Q]=2f$,
distinguished primes $q,d$, the \emph{type primes} $g_0,g_1$ (we write $\mathrm{gset}=\{g_0,g_1\}$), a finite set
$\mathrm{small}$ of small primes containing $2$, a level $Q_S$, and the class number $h_E$ of $E$. A prime
$p\notin\mathrm{small}\cup\{q,d\}$ is \emph{split} if the Jacobi symbol $\bigl(\frac{-d}{p}\bigr)$ equals $1$; otherwise the host assigns
it one of two \emph{inert} classes. For a place $\lambda$ of $E$, $n_\lambda=e_\lambda f_\lambda$ denotes its local degree. Its rows datum
contains the coefficient field $\mathit{LL}$, Galois over $\Q$, its finite places $w$ with residue characteristic
$\ell(w)$ and ramification index $e(w)$, the characters $\mu_w\in\{\mu_0,\mu^{\mathrm{tw}}_w\}$ ($w$ is \emph{twisted} if
$\mu_w\ne\mu_0$), complex spaces $Z_j(\mu)$, rational numbers $\eta_j(w)\ge0$ and $\mathrm{sh}(\ell,e)$, and the positive
integer
\begin{equation}\label{eq:pf:linkN:G}
\#G=\prod_{k\in\{0,1\}}\bigl(\#\GL_r(\F_{g_k})\bigr)^{f},
\end{equation}
the order of $G=\prod_{k\in\{0,1\}}\prod_{v\in\mathrm{Pl}_k}\GL_r(\F_{g_k})$, where $\mathrm{Pl}_k$ is a fixed set of $f$
places attached to $g_k$. We put $VV(\ell)=\sum_{0<e<t}\mathrm{sh}(\ell,e)$.

\paragraph{Local indices.}
The Galois data of $D$ contain an algebraic closure $\Omega$ of $E$, for each prime $p$ the set $S_p$ of places of $E$
above $p$ with the cyclotomic character of $G_{E,S_p}$, and a finite set $\mathrm{Loc}(D)$ of \emph{local indices}. A
local index $\iota$ carries: a prime $\ell=\ell_\iota$, an integer $1\le j_\iota\le m$ and a twisted place $w_\iota$ of
$\mathit{LL}$ with $\ell(w_\iota)=\ell_\iota$; a complete discrete valuation ring $\OO_\iota$, finite over $\Zl$, with
fraction field $K_\iota$ and an embedding of $\mathit{LL}$ into $K_\iota$; a finite non-empty set $X_\iota$
(\emph{eigensystems}); a local $\OO_\iota$-subalgebra $T_\iota\subseteq\OO_\iota^{X_\iota}$ separating the points of
$X_\iota$; a free $\OO_\iota$-module $M_\iota$ of finite rank with an injective homomorphism of $\OO_\iota$-algebras
$T_\iota\to\End(M_\iota)$; finite subsets $\linkNgood_\iota\subseteq Z_\iota\subseteq X_\iota$; characters
$\chi_{\iota,c}$ of $G_E$ with values in $\OO_\iota^\times$; and disjoint finite sets $S_\iota,Q_\iota$ of split places of
$E$ not above $\ell$, with elements $\gamma_{\iota,v}$ of the inertia groups at $v$ generating tame inertia for
$v\in S_\iota\cup Q_\iota$. We write $\mathrm{Loc}(j,w)$ for the set of $\iota$ with $(j_\iota,w_\iota)=(j,w)$.

\begin{definition}[Hecke eigenspaces and local invariants]\label{pf:linkN:localinv}
Let $\iota$ be a local index, $\OO=\OO_\iota$, $K=K_\iota$, $M=M_\iota$, $M_K=K\otimes_\OO M$, and let $M^\circ\subseteq M_K$
be the lattice $1\otimes M$.
\begin{enumerate}[label=(\alph*)]
\item For $\pi\in X_\iota$ the \emph{Hecke eigenspace} is $M_K[\pi]=\{v\in M_K:\ tv=t(\pi)v\text{ for all }t\in T_\iota\}$,
and $\linkNmult(\iota,\pi)=\dim_K M_K[\pi]$. For $Y\subseteq X_\iota$ put $M_K[Y]=\sum_{\pi\in Y}M_K[\pi]$.
\item The \emph{normalised length} of an $\OO$-module $N$ is $\lambda_\iota(N)=\len_\OO(N)/\len_\OO(\OO/\ell\OO)$ if
$\len_\OO(N)<\infty$, and $\lambda_\iota(N)=0$ otherwise.
\item $\linkNup_\iota=M_K[\linkNgood_\iota]\cap\bigl(M^\circ+M_K[X_\iota\smallsetminus\linkNgood_\iota]\bigr)$,
$\linkNlow_\iota=M^\circ\cap M_K[\linkNgood_\iota]$, and $\linkNloc(\iota)=\lambda_\iota(\linkNup_\iota/\linkNlow_\iota)$.
\item Let $\mathfrak c_\iota\subseteq T_\iota$ be the kernel of $t\mapsto(t(y))_{y\notin Z_\iota}$, and for $\pi\in X_\iota$
let $\mathfrak c_\pi\subseteq\OO$ be the ideal generated by the $t(\pi)$, $t\in\mathfrak c_\iota$. Then
$\linkNvcc(\iota,\pi)=\lambda_\iota(\OO/\mathfrak c_\pi)$.
\end{enumerate}
\formal{Formalism.BasicNotions.NData.etaLoc}
\formal{Formalism.BasicNotions.NData.vcc}
\end{definition}

When $M_K$ is the direct sum of the $M_K[\pi]$ (\cref{pf:linkN:decomp}), $\linkNup_\iota$ is the projection of the lattice
to the $\linkNgood_\iota$-part and $\linkNlow_\iota$ its intersection with that part, so $\linkNup_\iota/\linkNlow_\iota$ is
the congruence module of the lattice between the good eigensystems and the others; $\mathfrak c_\pi$ measures the
congruences between $\pi$ and the eigensystems outside $Z_\iota$. Two compatibility conditions, which are part of the
definition of a datum, tie these invariants to the rows datum:
\begin{equation}\label{eq:pf:linkN:ties}
\eta_j(w)=\sum_{\iota\in\mathrm{Loc}(j,w)}\linkNloc(\iota)\quad(1\le j\le m),
\end{equation}
and $\mathrm{sh}(p,e)=N$ whenever the $\Z_p$-length of $\Sha^1(G_{E,S_p},\Q_p/\Z_p(-e))$, $0<e<t$, is a natural number
$N$. Every local index is twisted (its place $w_\iota$ is twisted). Given representations
$\rho_{\iota,\pi}:G_E\to\GL_r(\OO_\iota)$, $\pi\in X_\iota$, the Selmer length $\linkNs(\iota,\rho,\pi)$ is the normalised
length of the Selmer group with Bloch--Kato conditions at $\ell$, relaxed at $S_\iota\cup Q_\iota$, of
\cref{sec:families:galois}.

\paragraph{Size data and $\sfname{LinkN}$.}
Size data $L$ for $D$ carry a scale $(x,\eps,J)$ with $x>1$ and $0<\eps<1/52$, a finite support $\supp(L)$ of places of
$\mathit{LL}$, and derived quantities: $\Lambda=x^{\eps}$, $b=e^{\Lambda}$, a non-negative real $N_c$, the dimensions
$\linkNdZ(j)=\dim_\C Z_j(\mu_0)$, the \emph{zero-space ranks} $\dim_\C Z_j(\mu^{\mathrm{tw}}_w)$ at twisted places $w$, and
$\linkNlw(w)=\log\Nm w/[\mathit{LL}:\Q]$.

\begin{lemma}[Host and sizes]\label{pf:linkN:hostsize}
\begin{enumerate}[label=(\roman*)]
\item Every prime $p\notin\mathrm{small}\cup\{q,d\}$ is split or inert.
\item For every prime $\ell$, $\sum_{\lambda\mid\ell}n_\lambda\le2f$, the sum over the places $\lambda$ of $E$ above $\ell$.
\item For size data $L$, $\Lambda=x^\eps\ge1$.
\end{enumerate}
\formallabel{Formalism.MainNotions.TwistedPlaces.split_or_inert}
\formallabel{Formalism.MainNotions.TwistedPlaces.sum_nLam_le}
\formallabel{Formalism.MainNotions.TwistedPlaces.L_ge_one}
\end{lemma}

\begin{proof}
(i) holds by the definition of the classes. (ii) The places above $\ell$ correspond to the prime ideals of the ring of integers of $E$
above $\ell$, and $\sum_{\mathfrak p\mid\ell}e_{\mathfrak p}f_{\mathfrak p}=[E:\Q]=2f$ \cite[Ch.~I, Prop.~8.2]{Neukirch1999} (so equality holds).
(iii) $x>1$ and $\eps>0$.
\end{proof}

\begin{definition}[The link]\label{pf:linkN:LinkN}
For size data $L$ of $D$ and $j\in\N$ put
\begin{equation}\label{eq:pf:linkN:etaSum}
\linkNeS(L,j)=\sum_{\substack{w\in\supp(L)\\ w\ \text{twisted}}}\linkNlw(w)\,e(w)\,\eta_j(w),\qquad
\mathrm{vv}(D,\ell)=\begin{cases}VV(\ell)\log\ell,&\ell\text{ prime},\\0,&\text{otherwise}.\end{cases}
\end{equation}
For $C\in\R$ let $\mathrm{LN}(L,j,C)$ be the statement: if $\sum_{\ell\in\N}\mathrm{vv}(D,\ell)$ converges, then
\begin{equation}\label{eq:pf:linkN:linkNat}
\linkNeS(L,j)\le \linkNdZ(j)\Bigl(C\bigl(f\,m\,b\,\Lambda^2+(\log x)^2N_c\bigr)+\sum_{\ell}\mathrm{vv}(D,\ell)\Bigr).
\end{equation}
Let $\linkNfin(D)$ mean that $\Sha^1(G_{E,S_p},\Q_p/\Z_p(-e))$ has finite length for every prime $p$ and $0<e<t$. A
family $\mathcal F$ satisfies $\sfname{LinkN}$ \emph{with constant $C$} if there is $x_0$ such that $\mathrm{LN}(L_{x,i},j,C)$
holds for all $x\ge x_0$, all $i\in I_x$ with $\linkNfin(D_{x,i})$ and all $1\le j\le m$.
\formal{Formalism.MainNotions.LinkN_at}
\formal{Formalism.MainNotions.LinkN}
\end{definition}

\paragraph{Families and norm choices.}
Families $\mathcal F\in\sfname{FamP}(\eps)$, classes, threshold packages, largeness and $\sfname{Closed}$ are those of
\cref{def:families:closed}; $D3(3,\mathcal A)$ is the class of designed families of a class $\mathcal A$, and every member
of $D3(3,\mathcal A)$ belongs to $\mathcal A$ and is $\sfname{TwGood}$. The member $(x,i)$ of $\mathcal F$ has datum
$D_{x,i}$, size data $L_{x,i}$, and for $x>x_1$ an extra record with norm choice $\mathrm{nch}_{x,i}$ and a parameter
choice $P_{x,i}$ (constants of the member).

\paragraph{Hypotheses.}
The following propositions occur as hypotheses of the formal statements. \Cref{pf:linkN:H2,pf:linkN:Hdec,pf:linkN:Hcheb,pf:linkN:Hdz}
are classical theorems; all of them, and the packages of \cref{pf:linkN:Hlayer}, are proved elsewhere in the formal development
\cite{LeanSource2026} (\cref{sec:families:theta}).

\begin{hypothesis}[$\sfname{Hypotheses}_2$]\label{pf:linkN:H2}
The conjunction of: the properties of Fontaine's period ring $B_{\mathrm{cris}}$ \cite{Fontaine1982,Fontaine1994}; global
Artin reciprocity \cite{Artin1927,Neukirch1999}; the local Bloch--Kato theory (dimension formula and dualities for the
finite parts $H^1_f$) \cite[\S3]{BlochKato1990}; finiteness of local Galois cohomology, local Tate duality and the local
Euler characteristic formula for finite modules over finite extensions of $\Ql$ \cite{Tate1963},
\cite[Ch.~II, \S5]{Serre1997}; and an integral comparison theorem of Bhatt--Scholze type for crystalline lattices
\cite{BhattScholze2022,BhattScholze2023}.
\formal{GlobalTheory.NLink.Hypotheses₂}
\end{hypothesis}

\begin{hypothesis}[Local cohomology of decomposition groups]\label{pf:linkN:Hdec}
For every number field $F$, algebraic closure $\Omega$, prime $\ell$, $r\ge1$, finite discrete $\Z/\ell^r$-module $M$,
finite place $v$ of $F$ and valuation ring $A$ of $\Omega$ above $v$, with decomposition group $D_A$, and every continuous
action of $D_A$ on $M$: (finiteness) $H^0,H^1,H^2$ are finite and $H^n=0$ for $n\ge3$; (duality)
$H^2(D_A,\mu_{\ell^r})\cong\Z/\ell^r$, and the cup products induce bijections
$H^0(M)\to\Hom(H^2(M^\vee(1)),H^2(\mu_{\ell^r}))$, $H^1(M^\vee(1))\to\Hom(H^1(M),H^2(\mu_{\ell^r}))$ and its transpose,
and $H^2(M)\to\Hom(H^0(M^\vee(1)),H^2(\mu_{\ell^r}))$, where $M^\vee(1)$ is $\Hom(M,\Z/\ell^r)$ twisted by the cyclotomic
character; (Euler characteristic) $\#H^1=\#H^0\,\#H^2\,(\#M)^{[F_v:\Ql]}$ if $v\mid\ell$ and $\#H^1=\#H^0\,\#H^2$ if
$v\nmid\ell$. These are Tate's theorems \cite{Tate1963}, \cite[Ch.~II, \S5]{Serre1997},
\cite[Ch.~VII]{NeukirchSchmidtWingberg2008}.
\formal{ReciprocityTheory.Gal.LocalFiniteness}
\formal{ReciprocityTheory.Gal.LocalTateDuality}
\formal{ReciprocityTheory.Gal.LocalEulerChar}
\end{hypothesis}

\begin{hypothesis}[Chebotarev]\label{pf:linkN:Hcheb}
For every number field $F$, open normal subgroup $U\subseteq G_F$, $g\in G_F$ and finite set $A$ of finite places, there
are a place $v\notin A$ of absolute degree one, unramified over $\Q$, a place of $\Omega$ above $v$ whose inertia group
lies in $U$, and a Frobenius element $\sigma$ at it with $\sigma g^{-1}\in U$. This follows from the Chebotarev density
theorem \cite{Chebotarev1926}, \cite[Ch.~VII, \S13]{Neukirch1999}.
\formal{GlobalTheory.TwSys.Cheb}
\end{hypothesis}

\begin{hypothesis}[Automorphic packages]\label{pf:linkN:Hlayer}
The three bundles of results of \cref{sec:families:theta} used for the construction: the \emph{third-layer inputs}
$\mathcal P_3$ (which contain \cref{pf:linkN:Hcheb}, local theta inputs, Harish-Chandra finiteness and rapid decay of cusp
forms); the \emph{compatibility package} $\mathcal P_{\mathrm{Gal}}$ for Galois representations attached to polarised
cuspidal representations (which contains a height-counting statement); and the \emph{twist-series package}
$\mathcal P_{\mathrm{ser}}$ (rationality and integrality of twist series).
\formal{GlobalTheory.HeightCounting.Up.PublishedInputsThirdLayerAlt}
\formal{GlobalTheory.NRows.GeometricGaloisRepCompatibility}
\formal{GlobalTheory.HeightCounting.Up.SerN}
\end{hypothesis}

\begin{hypothesis}[Depth-zero support]\label{pf:linkN:Hdz}
If a smooth irreducible representation of $\GL_n(k)$, $k$ a non-archimedean local field, has non-zero vectors fixed by
$1+\varpi M_n(\OO_k)$, then so has every member of its cuspidal support. This is a consequence of the results of Moy and
Prasad \cite{MoyPrasad1994,MoyPrasad1996}.
\formal{GlobalTheory.Published.DepthZeroSupport}
\end{hypothesis}

Finally $h_L$ always denotes a proof of $\sfname{RecTExists}$ (\cref{def:rectexists}).

\begin{remark}\label{pf:linkN:interp}
The local inequality $\linkNvcc\le\linkNs+(\text{error})$ has the shape of the numerical criterion of Wiles and Lenstra
\cite{Wiles1995,Lenstra1995} (see also \cite{TaylorWiles1995,DarmonDiamondTaylor1997}); this is an interpretation, and
the arguments below do not use it.
\end{remark}

% ---------------------------------------------------------------------------
\subsection{Hecke modules and their cuts}\label{sec:pf:linkN:lattice}

Throughout this subsection $\OO$ is a discrete valuation ring with fraction field $K$, $X$ is a finite set,
$T\subseteq\OO^X$ is an $\OO$-subalgebra separating the points of $X$, $M$ is a finitely generated free $\OO$-module and
$t\mapsto t_M$ is a homomorphism of $\OO$-algebras $T\to\End_\OO(M)$. We use $M_K$, $M^\circ$, $M_K[x]$, $M_K[Y]$ and
$\linkNmult(x)=\dim_KM_K[x]$ as in \cref{pf:linkN:localinv}, and for $\phi\in\End_\OO(M)$ we write $\phi_K$ for its
extension to $M_K$. For $Y\subseteq X$ put
\[
U_Y=M_K[Y]\cap\bigl(M^\circ+M_K[X\smallsetminus Y]\bigr),\qquad \mathit{Lo}_Y=M^\circ\cap M_K[Y].
\]

\begin{lemma}[Idempotents up to a scalar]\label{pf:linkN:decomp}
There are $d\in\OO\smallsetminus\{0\}$ and elements $g_x\in T$ ($x\in X$) with $g_x(y)=0$ for $y\ne x$ and
$dv=\sum_{x\in X}(g_x)_Mv$ for all $v\in M_K$. If $t\in T$ vanishes at every $y\ne x$, then $t_Mv\in M_K[x]$ for every
$v\in M_K$. In particular $M_K=M_K[X]$.
\formallabel{UpperTheory.GalC.Alg.decomp}
\end{lemma}

\begin{proof}
Fix $x$. For $y\ne x$ choose $t_y\in T$ with $t_y(x)\ne t_y(y)$; then $f_x=\prod_{y\ne x}(t_y-t_y(y))\in T$ satisfies
$f_x(x)\ne0$, as $\OO$ is a domain, and $f_x(y)=0$ for $y\ne x$. Put $d=\prod_xf_x(x)\ne0$ and
$g_x=\bigl(\prod_{y\ne x}f_y(y)\bigr)f_x$. At a point $z$ only $g_z$ is non-zero, and $g_z(z)=d$; hence $\sum_xg_x=d\cdot1$
in $T$, and $dv=\sum_x(g_x)_Mv$. If $t$ vanishes at every $y\ne x$, then $st=s(x)t$ in $\OO^X$ for all $s\in T$, so
$s_M(t_Mv)=s(x)\,t_Mv$, i.e.\ $t_Mv\in M_K[x]$. The last assertion follows.
\end{proof}

\begin{definition}[The cut module]\label{pf:linkN:cutmodule}
For $S\subseteq X$ let $\mathrm{out}_S\subseteq M$ be the set of $m\in M$ with $t_Mm=0$ for all $t\in T$ vanishing on
$X\smallsetminus S$. Put $M_S=M/\mathrm{out}_S$, with quotient map $q$ and $q_K:M_K\to(M_S)_K$; let $T_S\subseteq\OO^S$ be the
image of $T$ under restriction to $S$, acting on $M_S$ by $(t|_S)_{M_S}(q(m))=q(t_Mm)$.
\formallabel{UpperTheory.GalC.Alg.cutM}
\end{definition}

The action is well defined: $\mathrm{out}_S$ is $T$-stable because $T$ is commutative, and if $u\in T$ vanishes on $S$ and
$f\in T$ vanishes outside $S$ then $fu=0$, so $u_M(M)\subseteq\mathrm{out}_S$. If $a\in\OO\smallsetminus\{0\}$ and
$am\in\mathrm{out}_S$, then $a\,t_Mm=0$, hence $t_Mm=0$, for every $t$ vanishing outside $S$; so $M_S$ is torsion free and
therefore free. $T_S$ separates the points of $S$, and it is local if $T$ is.

\begin{lemma}[Kernel of the cut]\label{pf:linkN:qK}
Let $S\subseteq X$.
\begin{enumerate}[label=(\roman*)]
\item For $x\in S$, $q_K(M_K[x])$ lies in the eigenspace $(M_S)_K[x]$ of $T_S$.
\item $\ker q_K=M_K[X\smallsetminus S]$.
\item If $Y\subseteq X$ and $Y'\subseteq S$ satisfy $Y\cap S\subseteq Y'$, then $q_K(M_K[Y])\subseteq(M_S)_K[Y']$.
\end{enumerate}
\formallabel{UpperTheory.GalC.Alg.qK_eq_zero_iff}
\end{lemma}

\begin{proof}
(i) For $t\in T$ and $v\in M_K[x]$, $(t|_S)_{M_S}q_K(v)=q_K(t_Mv)=t(x)q_K(v)$.
(ii) Let $x\notin S$ and $v\in M_K[x]$. The element $f_x$ of the proof of \cref{pf:linkN:decomp} vanishes on $S$, so
$q_K((f_x)_Mv)=(f_x|_S)_{M_S}q_K(v)=0$; as $(f_x)_Mv=f_x(x)v$ with $f_x(x)\ne0$, $q_K(v)=0$. Conversely let $q_K(v)=0$ and
write $dv=\sum_x(g_x)_Mv$. If $x\in S$, then $g_x$ vanishes outside $S$, so $(g_x)_M$ kills $\mathrm{out}_S$ and factors as
$\ell_x\circ q$ with $\ell_x:M_S\to M$ linear; hence $(g_x)_Mv=(\ell_x)_K(q_Kv)=0$. If $x\notin S$, then $(g_x)_Mv\in M_K[x]$.
Thus $dv\in M_K[X\smallsetminus S]$, and so $v\in M_K[X\smallsetminus S]$.
(iii) For $x\in Y\cap S$ use (i) and $x\in Y'$; for $x\in Y\smallsetminus S$, $q_K(M_K[x])=0$ by (ii).
\end{proof}

\begin{lemma}[Comparison of congruence modules]\label{pf:linkN:lengthle}
Let $Z,S\subseteq X$, let $Z'=Z\cap S$, viewed as a subset of $S$, and let $U'_{Z'}$, $\mathit{Lo}'_{Z'}$ be the modules
defined as above for $(M_S,T_S)$. Let $\alpha\in\End_\OO(M)$ and $A\in\OO\smallsetminus\{0\}$ be such that $\alpha_K$ is
multiplication by $A$ on $M_K[x]$ for $x\in Z$ and $0$ on $M_K[x]$ for $x\notin S$. Then $U_Z/\mathit{Lo}_Z$ and
$U'_{Z'}/\mathit{Lo}'_{Z'}$ have finite length and
\[
\len_\OO\bigl(U_Z/\mathit{Lo}_Z\bigr)\le\len_\OO\bigl(U'_{Z'}/\mathit{Lo}'_{Z'}\bigr)
+\Bigl(\sum_{x\in Z}\linkNmult(x)\Bigr)\len_\OO(\OO/A\OO).
\]
\formallabel{UpperTheory.GalC.Alg.length_le}
\end{lemma}

\begin{proof}
\emph{Finiteness.} Let $d,g_x$ be as in \cref{pf:linkN:decomp}; $(g_y)_M$ acts on $M_K[x]$ by $g_y(x)$, which is $0$ for
$x\ne y$. Let $u\in U_Z$, $u=m+w$ with $m\in M^\circ$ and $w\in M_K[X\smallsetminus Z]$. Then $(g_y)_Mu=0$ for $y\notin Z$ and
$(g_x)_Mw=0$ for $x\in Z$, so $du=\sum_{x\in Z}(g_x)_Mm\in M^\circ\cap M_K[Z]=\mathit{Lo}_Z$. Hence $U_Z\subseteq d^{-1}M^\circ$ is
finitely generated, and $U_Z/\mathit{Lo}_Z$ is finitely generated and killed by $d\ne0$, so it has finite length. The same
argument applies to $M_S$, since $T_S$ separates the points of $S$.

\emph{The map.} We have $q_K(M^\circ)=M_S^\circ$, and by \cref{pf:linkN:qK}(iii), applied to $(Z,Z')$ and to
$(X\smallsetminus Z,S\smallsetminus Z')$, $q_K(U_Z)\subseteq U'_{Z'}$ and $q_K(\mathit{Lo}_Z)\subseteq\mathit{Lo}'_{Z'}$. Let
$\Phi:U_Z/\mathit{Lo}_Z\to U'_{Z'}/\mathit{Lo}'_{Z'}$ be the induced map.

\emph{$A$ kills $\ker\Phi$.} Let $u\in U_Z$ with $q_K(u)\in\mathit{Lo}'_{Z'}$. Then $q_K(u)=q_K(1\otimes m_0)$ for some $m_0\in M$,
so $u-1\otimes m_0\in M_K[X\smallsetminus S]$ by \cref{pf:linkN:qK}(ii). By linearity $\alpha_K$ vanishes on
$M_K[X\smallsetminus S]$ and is multiplication by $A$ on $M_K[Z]$; hence $Au=\alpha_K(u)=\alpha_K(1\otimes m_0)=1\otimes\alpha(m_0)\in
M^\circ$. Since $Au\in M_K[Z]$ as well, $Au\in\mathit{Lo}_Z$.

\emph{Length of the kernel.} Let $P\subseteq U_Z$ be the preimage of $\ker\Phi$. It is finitely generated and torsion free,
hence free, and contained in $M_K[Z]$; a basis of $P$ is $K$-linearly independent, so
$\rank P\le\dim_KM_K[Z]\le\sum_{x\in Z}\linkNmult(x)$. The module $\ker\Phi$ is generated by the images of a basis of $P$,
each killed by $A$, so it is a quotient of $(\OO/A\OO)^{\rank P}$ and $\len(\ker\Phi)\le\rank P\cdot\len(\OO/A\OO)$.

\emph{Conclusion.} By additivity of length on $0\to\ker\Phi\to U_Z/\mathit{Lo}_Z\to\operatorname{im}\Phi\to0$
\cite[Prop.~6.9]{AtiyahMacdonald1969}, $\len(U_Z/\mathit{Lo}_Z)=\len(\ker\Phi)+\len(\operatorname{im}\Phi)$, and
$\len(\operatorname{im}\Phi)\le\len(U'_{Z'}/\mathit{Lo}'_{Z'})$.
\end{proof}

\begin{lemma}\label{pf:linkN:lenunit}
Let $\iota$ be a local index of a datum, $\ell=\ell_\iota$, and $n\ge1$ an integer. Then
$\lambda_\iota(\OO_\iota/n\OO_\iota)\le v_\ell(n)$.
\formallabel{GlobalTheory.CutAtGood.lenUnit_natCast_le}
\end{lemma}

\begin{proof}
Write $n=\ell^an'$ with $\ell\nmid n'$. Since $\OO=\OO_\iota$ is local and finite over $\Zl$, $n'$ is a unit of $\OO$, so
$n\OO=\ell^a\OO$. If $\varpi$ is a uniformiser and $\ell\OO=\varpi^e\OO$, then $\len(\OO/\ell^a\OO)=ea=a\len(\OO/\ell\OO)$,
so $\lambda_\iota(\OO/n\OO)=a=v_\ell(n)$.
\end{proof}

% ---------------------------------------------------------------------------
\subsection{Congruence modules}\label{sec:pf:linkN:congmod}

This subsection proves a general bound for the length of congruence modules and two comparison statements; they are used
in the proof of the local bounds. Throughout, $\OO$ is a discrete valuation ring with fraction field $K$ and normalised
valuation $v$ (with $v(0)=\infty$), $X$ is a set, $T\subseteq\OO^X$ is an $\OO$-subalgebra, $V$ is a finite-dimensional
$K$-vector space and $\mathrm{act}:T\to\End_K(V)$ is a ring homomorphism. Lengths are lengths of $\OO$-modules, with values in
$\N\cup\{\infty\}$, and $0\cdot\infty=0$. For $x\in X$ and $S\subseteq X$ put
\[
V_x=\{u\in V:\ \mathrm{act}(f)u=f(x)u\text{ for all }f\in T\},\qquad V_S=\sum_{x\in S}V_x,
\]
and for an $\OO$-submodule $L\subseteq V$ define the \emph{congruence module}
\[
\mathrm{CM}(L,S)=L\big/\bigl((V_S\cap L)+(V_{X\smallsetminus S}\cap L)\bigr).
\]
If $V=M_K$ for a Hecke module $M$ as in \cref{sec:pf:linkN:lattice}, then $V_x=M_K[x]$ and $V_S=M_K[S]$. The first statement
only rewrites \cref{pf:linkN:localinv}(c) in this notation.

\begin{lemma}\label{pf:linkN:etaLoceq}
For every local index $\iota$ of a datum,
\[
\linkNloc(\iota)=\lambda_\iota\Bigl(\bigl((M^\circ+M_K[X_\iota\smallsetminus\linkNgood_\iota])\cap M_K[\linkNgood_\iota]\bigr)\big/
\bigl(M^\circ\cap M_K[\linkNgood_\iota]\bigr)\Bigr).
\]
\formallabel{Formalism.MainNotions.TruncatedSeriesCountingIdentities.Br.etaLoc_eq}
\end{lemma}

\begin{proof}
The numerator is $\linkNup_\iota$ (intersection is commutative) and the denominator is $\linkNlow_\iota\subseteq\linkNup_\iota$.
\end{proof}

\begin{lemma}[Eigenspaces span]\label{pf:linkN:eigspan}
Let $X$ be finite and suppose that $T$ separates the points of $X$.
\begin{enumerate}[label=(\roman*)]
\item There is $C\in\OO\smallsetminus\{0\}$ such that for every $Y\subseteq X$ some $e_Y\in T$ equals $C$ on $Y$ and $0$ on
$X\smallsetminus Y$.
\item Let $M$ be a finitely generated $\OO$-module with a homomorphism of $\OO$-algebras $T\to\End_\OO(M)$, and let $V=K\otimes_\OO M$
with the induced action. Then $\mathrm{act}(e_{\{x\}})V\subseteq V_x$ for $x\in X$, and $V=V_X$.
\end{enumerate}
\formallabel{Formalism.MainNotions.TruncatedSeriesCountingIdentities.Br.exists_ind}
\formallabel{Formalism.MainNotions.TruncatedSeriesCountingIdentities.Br.eigSpan}
\end{lemma}

\begin{proof}
(i) As in the proof of \cref{pf:linkN:decomp}, choose $f_x\in T$ with $f_x(x)\ne0$ and $f_x(y)=0$ for $y\ne x$, put
$C=\prod_xf_x(x)\ne0$ and $e_Y=\sum_{x\in Y}\bigl(\prod_{x'\ne x}f_{x'}(x')\bigr)f_x$.
(ii) For $f\in T$ we have $fe_{\{x\}}=f(x)e_{\{x\}}$ in $\OO^X$, so $\mathrm{act}(f)\mathrm{act}(e_{\{x\}})u=f(x)\mathrm{act}(e_{\{x\}})u$,
i.e.\ $\mathrm{act}(e_{\{x\}})u\in V_x$. Since $\sum_xe_{\{x\}}=C\cdot1$, every $u\in V$ satisfies $Cu=\sum_x\mathrm{act}(e_{\{x\}})u\in V_X$,
and $C$ is invertible in $K$.
\end{proof}

\begin{lemma}[Projection to the good part]\label{pf:linkN:projcm}
In the situation of \cref{pf:linkN:eigspan}(ii), let $L\subseteq V$ be an $\OO$-submodule and $S\subseteq X$. Then
\[
\len\Bigl(\bigl((L+V_{X\smallsetminus S})\cap V_S\bigr)\big/(L\cap V_S)\Bigr)\le\len\mathrm{CM}(L,S).
\]
\formallabel{Formalism.MainNotions.TruncatedSeriesCountingIdentities.Br.len_le}
\formallabel{Formalism.MainNotions.TruncatedSeriesCountingIdentities.Br.quot_le}
\end{lemma}

\begin{proof}
Let $C$ and $e_Y$ be as in \cref{pf:linkN:eigspan} and $P=C^{-1}\mathrm{act}(e_S)$, a $K$-linear endomorphism of $V$. On $V_x$,
$\mathrm{act}(e_S)$ is multiplication by $e_S(x)$; hence $P$ is the identity on $V_S$ and $0$ on $V_{X\smallsetminus S}$. As
$e_S=\sum_{x\in S}e_{\{x\}}$, \cref{pf:linkN:eigspan}(ii) gives $P(V)\subseteq V_S$, and since $e_S+e_{X\smallsetminus S}=C\cdot1$,
$u-Pu=C^{-1}\mathrm{act}(e_{X\smallsetminus S})u\in V_{X\smallsetminus S}$. Write $A=V_S$ and $B=V_{X\smallsetminus S}$, and consider
$L\to((L+B)\cap A)/(L\cap A)$, $u\mapsto Pu$. It is well defined, because $Pu=u-(u-Pu)\in L+B$ and $Pu\in A$. It is surjective: if
$a=u+b\in A$ with $u\in L$ and $b\in B$, then $Pu=Pa-Pb=a$. It kills $(A\cap L)+(B\cap L)$, since $P(a+b)=a\in L\cap A$ for $a\in A\cap L$
and $b\in B\cap L$. Hence the quotient on the left is a quotient of $\mathrm{CM}(L,S)$.
\end{proof}

The bound for congruence modules rests on the following three statements.

\begin{lemma}\label{pf:linkN:split}
Let $A\subseteq B$ be submodules of an $\OO$-module $N$, let $i:P\to N$ be $\OO$-linear and $C\subseteq P$ a submodule with $i(P)\subseteq B$,
$i(C)\subseteq A$ and $B\subseteq A+i(P)$. Then $\len(N/A)\le\len(N/B)+\len(P/C)$.
\formallabel{Formalism.NG.QuotientPowerLengths.split}
\end{lemma}

\begin{proof}
The kernel of $N/A\to N/B$ is $B/A$, and $\len(N/A)=\len(N/B)+\len(B/A)$. Every element of $B$ is $a+i(p)$ with $a\in A$, so $B/A$ is
contained in the image of the map $P/C\to N/A$ induced by $i$; hence $\len(B/A)\le\len(P/C)$.
\end{proof}

\begin{lemma}\label{pf:linkN:freequot}
Let $Q$ be a finitely generated torsion-free $\OO$-module and $l\in\OO$. Then $\len(Q/lQ)=\rank(Q)\,v(l)$.
\formallabel{Formalism.NG.QuotientPowerLengths.length_quot_smulRange_free}
\formallabel{Formalism.NG.QuotientPowerLengths.length_quot_smulRange_pi}
\end{lemma}

\begin{proof}
$Q$ is free, $Q\cong\OO^n$ with $n=\rank Q$, so $Q/lQ\cong(\OO/l\OO)^n$, and $\len(\OO/l\OO)=v(l)$: for $l=u\varpi^k$ with $u$ a unit
the chain $\varpi^j\OO/\varpi^k\OO$ is a composition series, and $\OO/0=\OO$ has infinite length because $\OO$ is not a field.
\end{proof}

\begin{lemma}\label{pf:linkN:core}
Let $\mathcal E$ be a set of $K$-linear endomorphisms of $V$, $S$ a finite set, and for $x\in S$ let $W_x\subseteq V$ be a subspace and
$w(x)\in\N\cup\{\infty\}$. Assume that every $g\in\mathcal E$ preserves every $W_x$, and that for every $x\in S$ there are $g\in\mathcal E$
and $l\in\OO$ such that $g$ is multiplication by $l$ on $W_x$ and $v(l)\le w(x)$. Let $U\subseteq V$ be a subspace preserved by $\mathcal E$
and $N\subseteq V$ a finitely generated $\OO$-submodule preserved by $\mathcal E$ with $N\subseteq\sum_{x\in S}W_x+U$. Let $\mathrm{img}(N)$ be
the $\OO$-span of the $g(n)$, $g\in\mathcal E$, $n\in N$, and $D(N)=\mathrm{img}(N)+(U\cap N)\subseteq N$. Then
\[
\len\bigl(N/D(N)\bigr)\le\sum_{x\in S}\dim_KW_x\cdot w(x).
\]
\formallabel{Formalism.NG.QuotientPowerLengths.core}
\formallabel{Formalism.NG.QuotientPowerLengths.step1}
\end{lemma}

\begin{proof}
Induction on $\#S$. If $S=\emptyset$, then $N\subseteq U$, so $D(N)=N$. Let $S=S_0\cup\{x\}$ with $x\notin S_0$, and choose $g\in\mathcal E$ and
$l\in\OO$ with $g|_{W_x}=l$ and $v(l)\le w(x)$. Put $U_1=\sum_{y\in S_0}W_y+U$, a subspace preserved by $\mathcal E$, and $N'=N\cap U_1$. Then
$N'$ is finitely generated ($\OO$ is noetherian), preserved by $\mathcal E$ and contained in $\sum_{y\in S_0}W_y+U$, so by induction
$\len(N'/D(N'))\le\sum_{y\in S_0}\dim W_y\,w(y)$.

\emph{Step 1: $\len\bigl(N/(lN+N')\bigr)\le\dim W_x\cdot v(l)$.} Let $q:N\to V/U_1$ be the projection and $Q=q(N)$. Then $Q$ is finitely generated
and torsion free, as $V/U_1$ is a $K$-vector space, and $N/(lN+\ker q)\cong Q/lQ$. Since $N\subseteq W_x+U_1$, $Q$ lies in the image of $W_x$
in $V/U_1$, and an $\OO$-basis of $Q$ is $K$-linearly independent there; so $\rank Q\le\dim W_x$, and \cref{pf:linkN:freequot} gives the
claim, as $\ker q=N'$.

\emph{Step 2.} Apply \cref{pf:linkN:split} to $A=D(N)$, $B=D(N)+lN+N'$, the inclusion $i:N'\to N$ and $C=D(N')$. Clearly $i(N')\subseteq B$, and
$D(N')\subseteq D(N)$ since $\mathrm{img}(N')\subseteq\mathrm{img}(N)$. It remains to see $lN\subseteq D(N)+N'$. Let $n\in N$ and write
$n=a+b$ with $a\in W_x$, $b\in U_1$. Then $g(n)\in\mathrm{img}(N)$, and $g(n)-ln=g(b)-lb$ lies in $U_1$ and in $N$, hence in $N'$; so
$ln=g(n)-(g(n)-ln)\in D(N)+N'$. Therefore
\begin{align*}
\len(N/D(N))&\le\len\bigl(N/(D(N)+lN+N')\bigr)+\len(N'/D(N'))\\
&\le\dim W_x\,v(l)+\sum_{y\in S_0}\dim W_y\,w(y),
\end{align*}
using Step 1, and $v(l)\le w(x)$ concludes.
\end{proof}

\begin{theorem}[Length of congruence modules]\label{pf:linkN:congmodle}
Assume $V=V_X$. Let $L\subseteq V$ be a finitely generated $\OO$-submodule, $S\subseteq X$ finite, $\mathcal E$ a set of $K$-linear
endomorphisms of $V$ commuting with $\mathrm{act}(T)$ and mapping $L$ into $L\cap V_S$, and $w:X\to\N\cup\{\infty\}$ such that for every $x\in S$
some $g\in\mathcal E$ is multiplication on $V_x$ by an element $l\in\OO$ with $v(l)\le w(x)$. Then
\[
\len\mathrm{CM}(L,S)\le\sum_{x\in S}\dim_KV_x\cdot w(x).
\]
\formallabel{Formalism.NG.QuotientPowerLengths.congMod_le}
\end{theorem}

\begin{proof}
An endomorphism $g$ commuting with $\mathrm{act}(T)$ preserves every $V_x$: for $u\in V_x$, $\mathrm{act}(f)g(u)=g(\mathrm{act}(f)u)=f(x)g(u)$;
hence it preserves every $V_Y$. Apply \cref{pf:linkN:core} with $W_x=V_x$ ($x\in S$), $U=V_{X\smallsetminus S}$ and $N=L$; the inclusion
$L\subseteq V_S+V_{X\smallsetminus S}$ holds because $V=V_X$. Since $\mathcal E$ maps $L$ into $L\cap V_S$, $D(L)\subseteq(V_S\cap L)+(V_{X\smallsetminus S}\cap L)$,
so $\mathrm{CM}(L,S)$ is a quotient of $L/D(L)$.
\end{proof}

% ---------------------------------------------------------------------------
\subsection{Automorphic data and the Galois preconditions}\label{sec:pf:linkN:gp}

Besides the Galois data, a datum carries rows, automorphic and type data. We recall the parts used below.

\paragraph{Rows.}
For $1\le j\le m$ the rows datum has a finite-dimensional complex space $M_j$ with subspaces $M_j(\sigma)$ indexed by the
\emph{rows} $\sigma$ (some rows are called \emph{good}), and for each character $\mu$ a complex space $M^t_{\mu,j}$ with a map
$\theta^P_{\mu,j}:M^t_{\mu,j}\to M_j$; $\theta^T_{\mu,j}$ is the adjoint of $\theta^P_{\mu,j}$, and $Z_j(\mu)$ is its image. For a
local index $\iota$ put $\mu_\iota=\mu^{\mathrm{tw}}_{w_\iota}$ and $M^t_\iota=M^t_{\mu_\iota,j_\iota}$.

\paragraph{Automorphic points.}
A local index $\iota$ has a set of \emph{automorphic points} $P$. Each has an eigensystem $\mathrm{pt}(P)\in X_\iota$ or none, and
every $x\in X_\iota$ is the eigensystem of some point; each has a non-zero space $\mathrm{Typ}(P)\subseteq M^t_\iota$ of \emph{type
vectors} and local parameters at the places of $E$. The \emph{admissibility clauses} (Adm) of the datum are conditions at the
local indices; among other things they attach to each automorphic point base change parameters, compatible with the character
choice, which determine its local parameters at the split places. The \emph{permutation provision} (Perm) (the forty-first
provision of the datum) contains, among other things, the statement that two automorphic points with the same eigensystem
have base change parameters which are permutations of each other, and the Galois family conditions. A local
parameter at a place above a type prime is \emph{type-admissible} if it is tame, monomial of depth zero and squarefree, in the
sense of the formal definition.

A row $\sigma$ \emph{localises at} $\iota$ if one of its automorphic parameters has, at the relevant split places, Hecke
polynomials congruent modulo the maximal ideal of $\OO_\iota$ to those of some eigensystem in $X_\iota$. Let $\mathrm{Rows}(\iota)$ be
the set of good rows localising at $\iota$ and $d_M(\iota)=\dim_\C\sum_{\sigma\in\mathrm{Rows}(\iota)}M_{j_\iota}(\sigma)$. We consider
the following conditions:
\begin{itemize}[leftmargin=3.2em]
\item[(ID)] (\emph{isotypic decomposition}) $M_j=\bigoplus_\sigma M_j(\sigma)$ for $1\le j\le m$;
\item[(SC)] (\emph{scalar clause}) for every $\iota$ and every $\sigma$ with $M_{j_\iota}(\sigma)\ne0$, the operator
$\theta^P_{\mu_\iota,j_\iota}\theta^T_{\mu_\iota,j_\iota}$ acts on $M_{j_\iota}(\sigma)$ as a positive real scalar;
\item[(MS)] (\emph{multiplicity statement}) if $\linkNgood_\iota\ne\emptyset$, then
$\sum_{\pi\in\linkNgood_\iota}\linkNmult(\iota,\pi)=d_M(\iota)$;
\item[(LU)] (\emph{uniqueness of localisations}) if $\iota,\iota'$ have the same pair $(j,w)$ and some good row lies in
$\mathrm{Rows}(\iota)\cap\mathrm{Rows}(\iota')$, then $\iota=\iota'$.
\end{itemize}
In the formal development (SC) and (MS) are parts of longer statements.

\paragraph{Type data.}
The group $G$ of \eqref{eq:pf:linkN:G} acts on the rows modules: central elements $z$ of $\C[G]$ act on $M^t_\iota$ by operators
$F\mapsto z\cdot F$. For $k\in\{0,1\}$ and a monic polynomial $f\in\F_{g_k}[X]$ (a \emph{label}) the type data contain a central
idempotent $e_{\mathrm{Ser}}(k,f)\in\C[\GL_r(\F_{g_k})]$, the idempotent of the series labelled by $f$; for $v\in\mathrm{Pl}_k$ we
write $e_{\mathrm{Ser}}(k,v,f)$ for its image in the factor of $\C[G]$ indexed by $(k,v)$. The \emph{good idempotent} is
\[
e_0=\sum_{p\ \mathrm{good}}\ \prod_{(k,v)}e_{\mathrm{Ser}}(k,v,p(k,v)),
\]
the sum over the good families $p=(p(k,v))$ of labels; it is central in $\C[G]$. The type data also contain \emph{parameters}
$\sigma'$ with subspaces $M^t_\iota(\sigma')$ and labels $\mathrm{lab}(k,\sigma',v)$; a parameter is good if its labels form a good
family. \emph{Series localisation} (Ser) is the condition: if $e_{\mathrm{Ser}}(k,v,f)\cdot F\ne0$ for some
$F\in M^t_\iota(\sigma')$, then $f=\mathrm{lab}(k,\sigma',v)$.

\paragraph{Pending data.}
A \emph{pending datum} $Q$ for $D$ contains, among other data: for every $\iota$ an element $\zeta_\iota\in\OO_\iota$ whose image under
the fixed embedding $\iota_\OO:\OO_\iota\to\C$ is a primitive $\#G$-th root of unity; an $\OO_\iota$-algebra homomorphism $z\mapsto z_Q$
from the centre of $\OO_\iota[G]$ to $\End_{\OO_\iota}(M_\iota)$ whose image commutes with $T_\iota$; and an expression of each
$e_{\mathrm{Ser}}(k,f)$ as the sum of the central idempotents of the irreducible representations $\tau$ of $\GL_r(\F_{g_k})$ with
label $f$. A \emph{type} of an automorphic point $P$ at $(k,v)$ is such a $\tau$ whose idempotent, placed at $(k,v)$, acts non-trivially
on some type vector of $P$. Two conditions on $Q$ are used:
\begin{itemize}[leftmargin=3.2em]
\item[(Eig)] (\emph{eigen-spanning}) for $\iota$, central $z\in\OO_\iota[G]$, $x\in X_\iota$ and $c\in\OO_\iota$: if $\iota_\OO(z)\cdot F=\iota_\OO(c)F$
for every type vector $F$ of every automorphic point with eigensystem $x$, then $(z_Q)_K$ is multiplication by $c$ on $M_K[x]$
(here $\iota_\OO(z)$ is obtained by applying $\iota_\OO$ to the coefficients);
\item[(TA)] (\emph{type admissibility}) if $\tau$ is a type of $P$ at $(k,v)$ whose label $\mathrm{lab}(\tau)$ is squarefree and
vanishes neither at a prescribed unit $u_0\in\F_{g_k}^\times$ nor at $u_0^{-1}$, then for $u=u_0$ or $u=u_0^{-1}$ the local
parameter of $P$ at the place attached to $v$ is monomial of depth zero with polynomial $\mathrm{lab}(\tau)(X-u)^t$, and at
the conjugate places it is monomial of depth zero, of degree $n_V$, with the inverse roots.
\end{itemize}

\begin{definition}[Galois preconditions, $\sfname{Tie4}$, norm-choice facts]\label{pf:linkN:GP}
Let $D$ be a datum and $\mathrm{nch}$ a norm choice for $D$.
\begin{enumerate}[label=(\alph*)]
\item The \emph{Galois preconditions} $\mathrm{GP}(D,\mathrm{nch})$ are eighteen conditions: (GP1) the point conditions (rows
conditions, non-vanishing of types, and $\mathrm{RE}(L)$ for all size data $L$); (GP2) admissibility of the datum, consisting of
(Adm), a residual irreducibility clause (including the non-eigenvalue condition at the auxiliary places above the type primes
used in \cref{pf:linkN:g5}) and (ID); (GP3) the first provision, consisting of (Perm), an unramifiedness statement for cuspidal
Galois representations and (SC); (GP4) (MS); (GP5) irreducibility of the residual representations of rank $r$; (GP6)
admissibility of residual tame eigenvalues; (GP7) a weight parameter condition; (GP8) a block character parameter condition;
(GP9) unramifiedness of the characters $\mu$; (GP10) a depth-zero condition; (GP11) a spherical inert base change condition;
(GP12) a condition on the parameters $\lambda$ used; (GP13) (LU); (GP14) largeness of the coefficient data (three conditions
on the coefficient rings); (GP15) a tie between $\mathrm{nch}$ and $D$; (GP16) a condition on triples; (GP17) a same-twist
condition for $\mathrm{nch}$; (GP18) a goodness condition for all rows.
\item $\sfname{Tie4}(D)$ is the conjunction of the \emph{character tie} (Chi) (at local indices with $\ell_\iota=g_k$, the
characters attached to the places $v\in\mathrm{Pl}_k$ are de Rham, with Weil--Deligne polynomial $X-u$ at the place attached to
$v$ and $X-u^{-1}$ at the conjugate places, for a unit $u$ equal to a prescribed character value or its inverse); the existence
of a pending datum $Q$ satisfying (TA) and (Eig); the \emph{inertia-polynomial tie} (Ine) (an annihilation statement for the
inertia polynomials acting on crystalline periods of potentially crystalline representations); and (Ser).
\item The \emph{norm-choice facts} $\mathrm{NF}(D,\mathrm{nch})$ consist of (GP1)--(GP3), (GP5)--(GP18); the
\emph{auxiliary-place clauses} (AP) of \cref{pf:linkN:g5}; the \emph{eigenvector clause} (EV) (if $r_x$ has an eigenvector with
eigencharacter $\chi_{\iota,1}$, then $x\in Z_\iota$ and its parameters are cuspidal); (MS); a de Rham condition (DR) (the good
representation families are de Rham away from the bad places, with distinct jumps); admissibility of the $J$-parameters of all
admissible size data; and the \emph{duties} (a hypothesis on the coefficient rings, absolute irreducibility with distinct
eigenvalues of the good residual representations, all design conditions, and a type-prime wild idempotent choice).
\end{enumerate}
\formallabel{UpperTheory.GalC.GaloisPreconditions}
\formallabel{UpperTheory.GalC.Tie4}
\formal{GlobalTheory.NFacts.NormChoiceFacts}
\end{definition}
% ---------------------------------------------------------------------------
\subsection{Cuts of data and the cut family}\label{sec:pf:linkN:cuts}

\begin{definition}[Cuts]\label{pf:linkN:cut}
A \emph{cut} $c$ of a datum $D$ consists of a set of \emph{kept} local indices and subsets $X^c_\iota\subseteq X_\iota$,
non-empty for kept $\iota$. The datum $\mathrm{cut}(D,c)$ is obtained from $D$ as follows. Its local indices are the kept
local indices of $D$; at a kept $\iota$ the set of eigensystems is $X^c_\iota$, the Hecke algebra is $(T_\iota)_{X^c_\iota}$
acting on $(M_\iota)_{X^c_\iota}$ (\cref{pf:linkN:cutmodule}), and the distinguished subsets are $Z_\iota\cap X^c_\iota$ and
$\linkNgood_\iota\cap X^c_\iota$; all other components of the local indices ($\ell_\iota,j_\iota,w_\iota,\OO_\iota$,
characters, auxiliary places) and all data indexed by local indices are restricted to the kept indices; in particular the
automorphic points at a kept $\iota$ are those of $D$, and such a point has eigensystem $x\in X^c_\iota$ in $\mathrm{cut}(D,c)$
if and only if it has eigensystem $x$ in $D$ (it has no eigensystem in $\mathrm{cut}(D,c)$ if its eigensystem in $D$ does
not lie in $X^c_\iota$); the values $\eta_j(w)$ are replaced by $\eta^c_j(w)=\sum\linkNloc^c(\iota)$, the sum over the kept
$\iota\in\mathrm{Loc}(j,w)$ of the invariants of the cut; all other components of $D$ are unchanged.
\formallabel{UpperTheory.GalC.cut}
\formallabel{UpperTheory.GalC.ndCut}
\end{definition}

The construction checks that the cut Hecke data are again of the required kind: $(M_\iota)_{X^c_\iota}$ is free,
$(T_\iota)_{X^c_\iota}$ is local and separates points, and its action is injective (if $t|_{X^c_\iota}$ acts by $0$, then by
\cref{pf:linkN:qK}(i),(ii) $t(x)$ kills $q_K(M_K[x])\cong M_K[x]$ for each $x\in X^c_\iota$, and these eigenspaces are
non-zero by the definition of a datum). Since the cut leaves the scale, the support, the twisted places, the dimensions
$\linkNdZ(j)$, the quantities $N_c,b,\Lambda$, the numbers $\mathrm{sh}(\ell,e)$ and the Galois data of $E$ unchanged, the
quantities $\mathrm{vv}(D,\ell)$ and the condition $\linkNfin$ are the same for $D$ and $\mathrm{cut}(D,c)$, and
$\linkNeS$ changes only through the values $\eta_j(w)$.

\begin{definition}[Good cuts]\label{pf:linkN:goodcut}
A cut $c$ of $D$ is \emph{good} if
\begin{enumerate}[label=(\alph*)]
\item $\iota$ is kept if and only if $\linkNgood_\iota\ne\emptyset$;
\item $\linkNgood_\iota\subseteq X^c_\iota$ and $Z_\iota\cap X^c_\iota\subseteq\linkNgood_\iota$;
\item for every kept $\iota$ there is $\alpha\in\End_{\OO_\iota}(M_\iota)$ commuting with $T_\iota$ such that $\alpha_K$ is
multiplication by $\#G$ on $M_K[\pi]$ for $\pi\in\linkNgood_\iota$ and $0$ on $M_K[\pi]$ for $\pi\notin X^c_\iota$;
\item for every kept $\iota$, every automorphic point of $\iota$ with eigensystem in $X^c_\iota$, every $k\in\{0,1\}$ and every
place $w$ of $E$ above $g_k$, the local parameter of the point at $w$ is type-admissible.
\end{enumerate}
The \emph{type-place condition} $\mathrm{TP}(D)$ is condition (d) for all local indices and all automorphic points having an
eigensystem.
\formal{UpperTheory.GalC.CutD.Good}
\formallabel{UpperTheory.GalC.CutRec}
\end{definition}

\begin{lemma}\label{pf:linkN:cutTP}
If $c$ is a good cut of $D$, then $\mathrm{TP}(\mathrm{cut}(D,c))$ holds.
\formallabel{UpperTheory.GalC.Tr.cutRec_cut}
\end{lemma}

\begin{proof}
Let $\iota$ be a kept local index and $P$ an automorphic point of $\iota$ having an eigensystem $x$ in $\mathrm{cut}(D,c)$.
Then $x\in X^c_\iota$, and by \cref{pf:linkN:cut} $P$ has eigensystem $x$ in $D$. Since the places above $g_k$ and the local
parameters of $P$ are not changed by the cut, \cref{pf:linkN:goodcut}(d) gives the type-admissibility of the local parameter
of $P$ at every place above $g_k$, for $k\in\{0,1\}$.
\end{proof}

\begin{definition}\label{pf:linkN:CGA}
$\sfname{CutGivesAlt}$ is the statement: for every datum $D$ and norm choice $\mathrm{nch}$ with $\mathrm{GP}(D,\mathrm{nch})$
and $\sfname{Tie4}(D)$ there is a good cut $c$ of $D$ such that $\mathrm{NF}(\mathrm{cut}(D,c),\mathrm{nch})$ and
$\sfname{Tie4}(\mathrm{cut}(D,c))$ hold, the norm choice being transported unchanged.
\formallabel{UpperTheory.GalC.CutGivesAlt}
\end{definition}

\begin{lemma}[$\eta$-invariants under a good cut]\label{pf:linkN:etale}
Let $c$ be a good cut of $D$, and write $\linkNloc^c$ for the invariants of $\mathrm{cut}(D,c)$.
\begin{enumerate}[label=(\roman*)]
\item If $\linkNgood_\iota=\emptyset$ then $\linkNloc(\iota)=0$.
\item For a kept local index $\iota$, with $\ell=\ell_\iota$,
$\linkNloc(\iota)\le\linkNloc^c(\iota)+\bigl(\sum_{\pi\in\linkNgood_\iota}\linkNmult(\iota,\pi)\bigr)v_\ell(\#G)$.
\item For $1\le j\le m$ and every place $w$ of $\mathit{LL}$,
\[
\eta_j(w)\le\eta^c_j(w)+\Bigl(\sum_{\iota\in\mathrm{Loc}(j,w)}\sum_{\pi\in\linkNgood_\iota}\linkNmult(\iota,\pi)\Bigr)
v_{\ell(w)}(\#G).
\]
\end{enumerate}
\formallabel{UpperTheory.GalC.eta_zero}
\formallabel{UpperTheory.GalC.eta_le}
\formallabel{UpperTheory.GalC.etaJ_le}
\end{lemma}

\begin{proof}
(i) Then $M_K[\linkNgood_\iota]=0$, so $\linkNup_\iota=0$.
(ii) Let $\alpha$ be as in \cref{pf:linkN:goodcut}(c). The integer $\#G$ is non-zero in $\OO_\iota$, since $\mathit{LL}$ embeds
into $K_\iota$. Apply \cref{pf:linkN:lengthle} with $Z=\linkNgood_\iota$, $S=X^c_\iota$ and $A=\#G$. By (b),
$Z\cap S=\linkNgood_\iota$, which is the good set of the cut, so the two quotients are those defining $\linkNloc(\iota)$ and
$\linkNloc^c(\iota)$. All lengths are finite, so dividing by $\len(\OO/\ell\OO)$ gives
$\linkNloc(\iota)\le\linkNloc^c(\iota)+(\sum_\pi\linkNmult(\iota,\pi))\lambda_\iota(\OO/\#G\OO)$, and \cref{pf:linkN:lenunit}
concludes.
(iii) By \eqref{eq:pf:linkN:ties}, $\eta_j(w)=\sum_{\iota\in\mathrm{Loc}(j,w)}\linkNloc(\iota)$, and by definition $\eta^c_j(w)$ is the
sum of $\linkNloc^c(\iota)$ over the kept $\iota\in\mathrm{Loc}(j,w)$; for $\iota\in\mathrm{Loc}(j,w)$ we have $\ell_\iota=\ell(w)$.
For kept $\iota$ use (ii). For $\iota$ not kept, $\linkNgood_\iota=\emptyset$ by \cref{pf:linkN:goodcut}(a), so
$\linkNloc(\iota)=0$ by (i), and the added term is non-negative. Summing gives the claim.
\end{proof}

\begin{lemma}\label{pf:linkN:meanw}
Let $L$ be size data for $D$ and $\ell$ a prime. Then $\sum_w\linkNlw(w)e(w)\le\log\ell$, the sum over the twisted places
$w\in\supp(L)$ above $\ell$. Consequently, for every integer $N\ge1$,
\[
\sum_{\substack{w\in\supp(L)\\ w\text{ twisted}}}\linkNlw(w)\,e(w)\,v_{\ell(w)}(N)\le\log N.
\]
\formallabel{Formalism.MainNotions.TwistedPlaces.admissibleMeanCondition_holds}
\formallabel{UpperTheory.GalC.sum_val_le}
\end{lemma}

\begin{proof}
If $w\mid\ell$ has residue degree $f(w)$, then $\Nm w=\ell^{f(w)}$ and $\linkNlw(w)e(w)=e(w)f(w)\log\ell/[\mathit{LL}:\Q]$;
since $\sum_{w\mid\ell}e(w)f(w)=[\mathit{LL}:\Q]$, the first claim follows. For the second, group the places according to
$\ell(w)$: only primes $\ell\mid N$ contribute, and the sum is at most $\sum_{\ell\mid N}v_\ell(N)\log\ell=\log N$.
\end{proof}

\begin{lemma}\label{pf:linkN:etaSumAdd}
Let $L$ be size data for $D$, $j\in\N$, $\eta'$ a real function on the places of $\mathit{LL}$, $d\ge0$ and $N\ge1$, and
suppose that $\eta_j(w)\le\eta'(w)+d\,v_{\ell(w)}(N)$ for every twisted place $w$. Then
\[
\linkNeS(L,j)\le\sum_{\substack{w\in\supp(L)\\ w\text{ twisted}}}\linkNlw(w)e(w)\eta'(w)+d\log N.
\]
\formallabel{UpperTheory.GalC.etaSum_le_add}
\end{lemma}

\begin{proof}
Multiply the hypothesis by $\linkNlw(w)e(w)\ge0$, sum, and apply \cref{pf:linkN:meanw}.
\end{proof}

\paragraph{Multiplicities.}
For size data $L$ of $D$ consider the \emph{multiplicity bound} $\mathrm{MB}(L)$ and the \emph{rank equality} $\mathrm{RE}(L)$:
for $1\le j\le m$ and every twisted place $w$,
\[
\mathrm{MB}:\ \sum_{\iota\in\mathrm{Loc}(j,w)}\sum_{\pi\in\linkNgood_\iota}\linkNmult(\iota,\pi)\le\dim_\C Z_j(\mu^{\mathrm{tw}}_w),
\qquad \mathrm{RE}:\ \dim_\C Z_j(\mu^{\mathrm{tw}}_w)=\linkNdZ(j).
\]
The rank equality is one of the point conditions (GP1) of the Galois preconditions; the multiplicity bound follows from four
other preconditions.

\begin{lemma}\label{pf:linkN:multle}
If $D$ satisfies (ID), (SC), (MS) and (LU), then $\mathrm{MB}(L)$ holds for all size data $L$ of $D$.
\formallabel{GlobalTheory.T.goodPointMultiplicityLeRank_proof}
\formallabel{GlobalTheory.NLink.multLe_of_rank}
\end{lemma}

\begin{proof}
Fix $1\le j\le m$ and a twisted $w$; we may assume $\mathrm{Loc}(j,w)\ne\emptyset$. By (MS), and trivially when
$\linkNgood_\iota=\emptyset$, $\sum_{\pi\in\linkNgood_\iota}\linkNmult(\iota,\pi)\le d_M(\iota)$, and by (ID)
$d_M(\iota)=\sum_{\sigma\in\mathrm{Rows}(\iota)}\dim M_j(\sigma)$. By (LU) the sets $\mathrm{Rows}(\iota)$, $\iota\in\mathrm{Loc}(j,w)$,
are pairwise disjoint, so by (ID) again
\[
\sum_{\iota\in\mathrm{Loc}(j,w)}d_M(\iota)=\dim_\C\bigoplus_{\sigma\in\bigcup_\iota\mathrm{Rows}(\iota)}M_j(\sigma)\le\dim_\C M_j.
\]
Let $\mu=\mu_\iota=\mu^{\mathrm{tw}}_w$ for $\iota\in\mathrm{Loc}(j,w)$. By (SC), $\theta^P_{\mu,j}\theta^T_{\mu,j}$ preserves each
$M_j(\sigma)$ and acts on it by a non-zero scalar, so by (ID) it is injective on $M_j$, and hence so is $\theta^T_{\mu,j}$.
Therefore $\dim M_j=\dim Z_j(\mu^{\mathrm{tw}}_w)$.
\end{proof}

\begin{definition}[The cut family]\label{pf:linkN:gd}
Let $\mathcal F\in\sfname{FamP}(\eps)$, and let $c_{x,i}$ be a cut of $D_{x,i}$ for every $x$ and $i\in I_x$. The family
$\mathrm{gd}(\mathcal F,c)$ has the index sets, threshold and constants of $\mathcal F$, members
$\mathrm{cut}(D_{x,i},c_{x,i})$, and size data, parameter choices, ties, lattice instances and extra records (with the norm
choices) transported along the cuts.
\formallabel{UpperTheory.GalC.gd}
\end{definition}

\begin{lemma}\label{pf:linkN:twgd}
If $\mathcal F$ is $\sfname{TwGood}$, so is $\mathrm{gd}(\mathcal F,c)$.
\formallabel{UpperTheory.GalC.twGood_gd}
\end{lemma}

\begin{proof}
The norm choices and the twisting data at twisted places are unchanged by the cuts. The only ingredient of the goodness
conditions of \cref{sec:families:galois} that involves the local Galois data of a datum is a family of lengths of
Selmer-type groups attached to the twisting characters (in the Bloch--Kato reduction conditions); these do not involve
the local indices and coincide for $D$ and $\mathrm{cut}(D,c)$. Hence every condition of $\sfname{TwGood}$ at the members of
$\mathcal F$ transfers to $\mathrm{gd}(\mathcal F,c)$, with the same threshold.
\end{proof}

\begin{lemma}\label{pf:linkN:evpow}
For $C\in\R$ and $\eps>0$ there is $x_0$ such that $x^C\le\exp(x^\eps)$ and $1\le x^\eps$ for all $x\ge x_0$.
\formallabel{UpperTheory.GalC.ev_pow_le_exp}
\end{lemma}

\begin{proof}
Since $\log x=o(x^\eps)$, eventually $C\log x\le x^\eps$ and $x\ge1$.
\end{proof}

We use the polynomial bound for the parameters of a family proved in
\lref{GlobalTheory.T.admissible_parameters_family_const}: \emph{for every family $\mathcal F$ there are $C_1\ge1$ and
$x'_1$ such that for $x\ge x'_1$, $x>x_1$ and $i\in I_x$ the member satisfies $k,n,t\le x^{C_1}$, $\log h_E\le C_1x^{C_1}$ and}
\begin{equation}\label{eq:pf:linkN:logG}
\log\#G\le C_1x^{C_1}\Lambda .
\end{equation}

\begin{proposition}[Return from the cut family]\label{pf:linkN:return}
Let $\mathcal F\in\sfname{FamP}(\eps)$, and let $c_{x,i}$ be cuts such that for all $x>x_1$ and $i\in I_x$ the cut $c_{x,i}$ is
good and $\mathrm{MB}(L_{x,i})$ and $\mathrm{RE}(L_{x,i})$ hold. If $\mathrm{gd}(\mathcal F,c)$ satisfies $\sfname{LinkN}$ with
constant $C$, then $\mathcal F$ satisfies $\sfname{LinkN}$ with constant $C+C_1$, where $C_1$ is the constant of
\eqref{eq:pf:linkN:logG}.
\formallabel{UpperTheory.GalC.linkN_of_gd}
\end{proposition}

\begin{proof}
Let $x_0$ be the threshold of $\sfname{LinkN}$ for the cut family, $x'_1$ that of \eqref{eq:pf:linkN:logG}, and $x_2$ that of
\cref{pf:linkN:evpow} for $(C_1,\eps)$; here $\eps>0$ because the scales have positive exponent. Let
$x\ge\max(x_0,x'_1,x_2,x_1+1)$, $i\in I_x$ with $\linkNfin(D_{x,i})$, $1\le j\le m$, and assume that
$\sum_\ell\mathrm{vv}(D_{x,i},\ell)$ converges. Write $D=D_{x,i}$ and $L=L_{x,i}$; $L$ has scale $(x,\eps)$, so $\Lambda=x^\eps\ge1$
(\cref{pf:linkN:hostsize}(iii))
and $x^{C_1}\le b$.

By \cref{pf:linkN:etale}(iii), $\mathrm{MB}(L)$ and $\mathrm{RE}(L)$, for every twisted place $w$
\[
\eta_j(w)\le\eta^c_j(w)+\linkNdZ(j)\,v_{\ell(w)}(\#G).
\]
\Cref{pf:linkN:etaSumAdd} with $\eta'=\eta^c_j$, $d=\linkNdZ(j)$ and $N=\#G$ gives $\linkNeS(L,j)\le\linkNeS^c(L,j)+\linkNdZ(j)\log\#G$,
where $\linkNeS^c$ is the $\eta$-sum of the cut member. The cut member has the same $\linkNfin$, $\mathrm{vv}$, $\linkNdZ$, $f$, $m$,
$b$, $\Lambda$, $x$ and $N_c$ (\cref{pf:linkN:cut}), so the hypothesis gives
$\linkNeS^c(L,j)\le\linkNdZ(j)\bigl(C(fmb\Lambda^2+(\log x)^2N_c)+\sum_\ell\mathrm{vv}(D,\ell)\bigr)$. Finally, by
\eqref{eq:pf:linkN:logG}, $x^{C_1}\le b$, $\Lambda\le\Lambda^2$, $fm\ge1$ (as $1\le m\le f$) and $(\log x)^2N_c\ge0$,
\[
\log\#G\le C_1x^{C_1}\Lambda\le C_1\,b\Lambda^2\le C_1\bigl(fmb\Lambda^2+(\log x)^2N_c\bigr).
\]
Adding the two bounds gives $\mathrm{LN}(L,j,C+C_1)$.
\end{proof}

% ---------------------------------------------------------------------------
\subsection{Summation of the local bounds}\label{sec:pf:linkN:sum}

In this subsection $D$ is a datum, $L$ size data with scale $(x,\eps)$, $P$ a parameter choice for $(D,L)$, and $C\in\R$.
For a local index $\iota$ with $\ell=\ell_\iota$ and a function $f$ on the places of $E$ we write $\sum_{\lambda\mid\ell}f(\lambda)$
for the sum over the places $\lambda$ of $E$ above $\ell$. For a prime $\ell$ and $1\le j\le m$ put
\[
\eta_\ell(j)=\sum_{\substack{w\in\supp(L)\text{ twisted}\\ w\mid\ell}}\linkNlw(w)e(w)\eta_j(w),\qquad\text{so that}\qquad
\linkNeS(L,j)=\sum_{\ell\in S(L)}\eta_\ell(j),
\]
where $S(L)$ is the set of residue characteristics of the twisted places in $\supp(L)$. For a prime $\ell$ let
$\nu_L(\ell)=\min_{g\in G_E}v_\ell\bigl(\chi_\ell(g)^{tr}-1\bigr)$, $\chi_\ell$ the $\ell$-adic cyclotomic character of $G_E$,
and let $\nu_{\mathrm{budget}}(D)=\log Q_S+\log(tr)$. The size data provide further positive quantities: a threshold
$\ell^*=\ell^*(L)$, two numbers $S_c(L)$ and $\eta_{\mathrm{sz}}(L)$, and the host provides a finite set $TT$. The
\emph{polynomial parameter bound} $\mathrm{PB}(L,C)$ is the statement: the constants of $P$ are at most $C$ (one of them is
positive), $k,n,t\le x^C$, $x^{-\eps}$ is at most a size quantity of $L$, $\log h_E\le Cx^C$ and $\log\#G\le Cx^C\Lambda$.

\begin{definition}[The local bounds]\label{pf:linkN:Concl}
For a family $\rho$ of representations $\rho_{\iota,\pi}$, a local index $\iota$ ($\ell=\ell_\iota$), $\pi\in X_\iota$ and $\beta\in\R$,
the \emph{local bound with error $\beta$} is \eqref{eq:families:bd}:
\[
\mathrm{Bd}(\rho,\iota,\pi,\beta):\qquad\linkNvcc(\iota,\pi)\le\linkNs(\iota,\rho,\pi)+\beta+VV(\ell)+v_\ell(h_E)+\ell^{-2}.
\]
\emph{Hub data} $X$ for $D$ consist of representation families $r,\rho$, local lengths $\mathrm{len}_0(\iota,\pi,\lambda)$, error terms
$\mathrm{err}(\iota,\pi)$, a predicate $\mathrm{band}(\iota,\lambda)$, a finite set $\mathrm{exc}$ of primes and further data. A prime
$\ell$ is in the \emph{middle range} if $\ell\notin\mathrm{small}\cup\{3,q,d\}$ and not ($\ell>\ell^*$ and $\ell\notin\mathrm{exc}$). Put
$\mathrm{mid}(\ell)=Cf\bigl(r(b/\ell+S_c+t+\eta_{\mathrm{sz}}+\Lambda)+mb\bigr)+C(frt)^C\Lambda$. The \emph{local bounds}
$\mathcal T(L,P,X,C)$ are the following statements, where $\iota$ ranges over the local indices, $\ell=\ell_\iota$, and $\pi$ over
$\linkNgood_\iota$:
\begin{description}[style=nextline,leftmargin=1.2em,font=\normalfont\itshape]
\item[positivity] $C\ge0$;
\item[passage] if $\ell\notin\mathrm{gset}$: $\linkNloc(\iota)\le\sum_{\pi\in\linkNgood_\iota}\linkNmult(\iota,\pi)
\bigl(\linkNvcc(\iota,\pi)+v_\ell(\#G)\bigr)$;
\item[multiplicity and rank] $\mathrm{MB}(L)$ and $\mathrm{RE}(L)$;
\item[killing] the Selmer groups of the $\rho_{\iota,\pi}$ are finite, $\linkNs(\iota,\rho,\pi)\le\nu_L(\ell)+\ell^{-2}$, and
$t\sum_{v\in Q_\iota}\lambda_\iota(\OO_\iota/(\zeta_v-1))\le\ell^{-2}$, where $\zeta_v=\chi_{\iota,1}(\gamma_{\iota,v})$;
\item[numerical bounds] four explicit upper bounds for the lengths $\max_\pi\mathrm{len}_0(\iota,\pi,\lambda)$, at split primes
$\ell\ge5$ and at inert primes in three ranges, and an approximation property of the band places;
\item[middle range] if $\ell\notin\mathrm{gset}$ is in the middle range:
$\linkNvcc\le\linkNs+\sum_{\lambda\mid\ell}\mathrm{len}_0(\iota,\pi,\lambda)+VV(\ell)+v_\ell(h_E)+\mathrm{err}(\iota,\pi)+\ell^{-2}$, and
$\mathrm{err}(\iota,\pi)\le Cfnt^2n_V\Lambda^2$;
\item[large primes] $\mathrm{Bd}(\rho,\iota,\pi,0)$ if $\ell>\ell^*$ and $\ell\notin\mathrm{gset}\cup\mathrm{exc}$;
\item[exceptional primes] $\#\mathrm{exc}\le(\#TT+2)(S_c+3\eta_{\mathrm{sz}}+t+8b/\ell^*+9)$, and every element of $\mathrm{exc}$ is
at most $3b$;
\item[small primes] there is a table $\mathrm{tf}$ (\cref{pf:linkN:tab}) with
$\sum_{\lambda\mid\ell}\mathrm{len}^{\mathrm{tf}}(\iota,\lambda)+BB^{\mathrm{tf}}(\iota)\le Cfmb\log x+Cx^C\Lambda^2$ and
$\mathrm{Bd}\bigl(\rho,\iota,\pi,\sum_{\lambda\mid\ell}\mathrm{len}^{\mathrm{tf}}(\iota,\lambda)+BB^{\mathrm{tf}}(\iota)\bigr)$ for
$\ell\in\mathrm{small}\cup\{3\}$;
\item[ramified primes] $\mathrm{Bd}(\rho,\iota,\pi,\mathrm{mid}(\ell))$ for $\ell\in\{q,d\}$;
\item[type primes] for $1\le j\le m$ and $\ell\in\mathrm{gset}$:
$\eta_\ell(j)\le\log\ell\cdot\linkNdZ(j)\bigl(Cx^C\eta_{\mathrm{sz}}+VV(\ell)+v_\ell(h_E)+2\ell^{-2}\bigr)$;
\item[parameters] $\mathrm{PB}(L,C)$;
\item[sizes] $3,q,d\le\ell^*$, every small prime is $\le\ell^*$, $3$ is not inert in the imaginary quadratic field of the host,
and $\mathrm{mid}(\ell)\le2Cfmb$ for $\ell\in\{q,d\}$.
\end{description}
A family $\mathcal F$ \emph{satisfies the local bounds with constant $C$} if there is $x_0$ such that for all $x\ge x_0$ with
$x>x_1$ and all $i\in I_x$ with $\linkNfin(D_{x,i})$ there are hub data $X$ with $\mathcal T(L_{x,i},P_{x,i},X,C)$.
\formallabel{GlobalTheory.NLink.Concl}
\formallabel{GlobalTheory.NLink.ConclF}
\end{definition}

\begin{definition}[Tables and admitted statements]\label{pf:linkN:tab}
A \emph{table} $\mathrm{tb}$ for $L$ consists of reals $s(\iota,\pi)$, $\mathrm{len}(\iota,\lambda)$, $BB(\iota)$,
$\mathrm{gcost}(j,\ell)$, $\mathrm{gs}(j,\ell)$, a predicate $\mathrm{band}(\iota,\lambda)$ and a finite set $\mathrm{exc}$ of
primes. The \emph{admitted statements} $\mathcal A(L,\mathrm{tb},C,P)$ are seventeen statements; with $\iota$, $\ell=\ell_\iota$ and
$\pi\in\linkNgood_\iota$ as before, the ones used below are:
\begin{enumerate}[label=(A\arabic*)]
\item\label{pf:linkN:A1} the passage inequality of \cref{pf:linkN:Concl};
\item\label{pf:linkN:A2} the inequality $\mathrm{MB}$ at twisted places $w$ with $\ell(w)\notin\mathrm{gset}$;
\item\label{pf:linkN:A3} $\mathrm{RE}(L)$;
\item\label{pf:linkN:A4} (\emph{signature matching}) for $\ell\notin\mathrm{gset}$:
$\linkNvcc(\iota,\pi)\le s(\iota,\pi)+\sum_{\lambda\mid\ell}\mathrm{len}(\iota,\lambda)+VV(\ell)+v_\ell(h_E)+BB(\iota)+\ell^{-2}$;
\item\label{pf:linkN:A5} $s(\iota,\pi)+\ell^{-2}\le\nu_L(\ell)+2\ell^{-2}$, and $\mathrm{gs}(j,\ell)+\ell^{-2}\le2\ell^{-2}$ for
$\ell\in\mathrm{gset}$;
\item\label{pf:linkN:A6} (\emph{type primes}) for $1\le j\le m$ and $\ell\in\mathrm{gset}$:
$\eta_\ell(j)\le\log\ell\cdot\linkNdZ(j)\bigl(\mathrm{gcost}(j,\ell)+\mathrm{gs}(j,\ell)+VV(\ell)+v_\ell(h_E)+\ell^{-2}\bigr)$ and
$\mathrm{gcost}(j,\ell)\le Cx^C\eta_{\mathrm{sz}}$;
\item\label{pf:linkN:A7} $\mathrm{PB}(L,C)$;
\end{enumerate}
the others are: $BB(\iota)\le Cfnt^2n_V\Lambda^2$ for $\ell\notin\mathrm{small}\cup\{q,d\}\cup\mathrm{gset}$; the five numerical
statements of \cref{pf:linkN:Concl} for the lengths $\mathrm{len}$; vanishing of $\mathrm{len}(\iota,\cdot)$ and $BB(\iota)$ for $\ell>\ell^*$,
$\ell\notin\mathrm{gset}\cup\mathrm{exc}$, together with the bound for $\#\mathrm{exc}$; $\mathrm{exc}\subseteq[0,3b]$;
$\sum_{\lambda\mid\ell}\mathrm{len}+BB\le Cfmb\log x+Cx^C\Lambda^2$ for $\ell\in\mathrm{small}\cup\{3\}$; and
$\sum_{\lambda\mid\ell}\mathrm{len}+BB\le\mathrm{mid}(\ell)\le2Cfmb$ for $\ell\in\{q,d\}$.
\formal{Formalism.MainNotions.TwistedPlaces.AdmT}
\end{definition}

\begin{lemma}[A table from the local bounds]\label{pf:linkN:admitted}
If $\mathcal T(L,P,X,C)$ holds, then there is a table $\mathrm{tb}$ with $\mathcal A(L,\mathrm{tb},C,P)$.
\formallabel{GlobalTheory.NLink.admittedStatements_of_localBoundsConclusion}
\formallabel{GlobalTheory.T.admittedStatements_of_localBoundsConclusion_alt}
\formallabel{GlobalTheory.NLink.admissible_table_of_conclusion}
\end{lemma}

\begin{proof}
Let $\mathrm{tf}$ be the table of the small-prime clause, write $\mathrm{sup}_\iota(g)=\max_{\pi\in\linkNgood_\iota}g(\pi)$ ($0$ if
$\linkNgood_\iota=\emptyset$), and define $\mathrm{tb}$ by $s(\iota,\pi)=\linkNs(\iota,\rho,\pi)$, $\mathrm{band}$ and $\mathrm{exc}$
those of $X$, $\mathrm{gcost}(j,\ell)=Cx^C\eta_{\mathrm{sz}}$, $\mathrm{gs}(j,\ell)=\ell^{-2}$, and, by cases on $\ell=\ell_\iota$,
\[
\bigl(\mathrm{len}(\iota,\lambda),BB(\iota)\bigr)=\begin{cases}
\bigl(0,\ \sum_{\lambda'\mid\ell}\mathrm{len}^{\mathrm{tf}}(\iota,\lambda')+BB^{\mathrm{tf}}(\iota)\bigr),&\ell\in\mathrm{small},\\
\bigl((\sum_{\lambda'\mid3}\mathrm{len}^{\mathrm{tf}}(\iota,\lambda')+BB^{\mathrm{tf}}(\iota))/\#\{\lambda'\mid3\},\ 0\bigr),&\ell=3,\\
\bigl(0,\ \mathrm{mid}(\ell)\bigr),&\ell\in\{q,d\},\\
\bigl(\mathrm{sup}_\iota(\mathrm{len}_0(\iota,\cdot,\lambda)),\ \mathrm{sup}_\iota(\mathrm{err}(\iota,\cdot))\bigr),&\ell\text{ in the middle range},\\
(0,0),&\text{otherwise}.
\end{cases}
\]
(Here $q,d\ne3$ by the host axioms, and $q,d$ are not small.) The statements \ref{pf:linkN:A1}--\ref{pf:linkN:A3} and
\ref{pf:linkN:A7} are clauses of $\mathcal T$. For \ref{pf:linkN:A4} let $\ell\notin\mathrm{gset}$ and $\pi\in\linkNgood_\iota$; in each
case a local bound $\mathrm{Bd}(\rho,\iota,\pi,\beta)$ with $\beta\le\sum_{\lambda\mid\ell}\mathrm{len}(\iota,\lambda)+BB(\iota)$ is
available, which is \ref{pf:linkN:A4} since $s(\iota,\pi)=\linkNs(\iota,\rho,\pi)$: for $\ell$ small or $\ell=3$ by the small-prime
clause (for $\ell=3$ the sum over $\lambda\mid3$ of the constant length is the numerator); for $\ell\in\{q,d\}$ by the
ramified-prime clause; in the middle range by the middle-range clause, as
$\mathrm{len}_0(\iota,\pi,\lambda)\le\mathrm{sup}_\iota(\mathrm{len}_0(\iota,\cdot,\lambda))$ and $\mathrm{err}(\iota,\pi)\le\mathrm{sup}_\iota(\mathrm{err})$;
otherwise $\ell>\ell^*$ and $\ell\notin\mathrm{exc}$, and the large-prime clause gives $\beta=0$. \ref{pf:linkN:A5} follows from
the killing clause and the definition of $\mathrm{gs}$, and \ref{pf:linkN:A6} is the type-prime clause. Of the remaining
statements, the bound for $BB$ follows from the middle-range clause; the numerical bounds hold in the middle range by the
numerical clause (there $\mathrm{tb}$ agrees with the table of the hub data) and trivially elsewhere (there $\mathrm{len}=0$,
since $3$ lies in none of the ranges concerned: the split range requires $\ell\ge5$, and $3$ is not inert by the size clause;
and the right-hand sides are non-negative); the vanishing for
$\ell>\ell^*$ outside $\mathrm{gset}\cup\mathrm{exc}$ holds by construction, since such $\ell$ is not small, not $3$, $q$ or
$d$ (these are $\le\ell^*$) and not in the middle range; the statements on $\mathrm{exc}$ are the exceptional-prime clause;
the bound at the small primes and $3$ is that of the small-prime clause; and at $q,d$ the first inequality holds by
construction and the second is the size clause.
\end{proof}

\begin{lemma}[Bound at one prime]\label{pf:linkN:etaAt}
Let $1\le j\le m$ and let $\ell$ be a prime.
\begin{enumerate}[label=(\roman*)]
\item If $\ell\notin\mathrm{gset}$, $U\ge0$, \ref{pf:linkN:A1}--\ref{pf:linkN:A3} hold, and
$\linkNvcc(\iota,\pi)+v_\ell(\#G)\le U$ for all $\iota$ with $\ell_\iota=\ell$ and all $\pi\in\linkNgood_\iota$, then
$\eta_\ell(j)\le\log\ell\cdot\linkNdZ(j)\,U$.
\item If $\ell\in\mathrm{gset}$ and \ref{pf:linkN:A5}, \ref{pf:linkN:A6} hold, then
$\eta_\ell(j)\le\log\ell\cdot\linkNdZ(j)\bigl(Cx^C\eta_{\mathrm{sz}}+2\ell^{-2}+VV(\ell)+v_\ell(h_E)\bigr)$.
\end{enumerate}
\formallabel{Formalism.MainNotions.TwistedPlaces.etaAt_le}
\formallabel{Formalism.MainNotions.TwistedPlaces.etaAt_le_g}
\end{lemma}

\begin{proof}
(i) Let $w\mid\ell$ be twisted. By \eqref{eq:pf:linkN:ties}, \ref{pf:linkN:A1}, \ref{pf:linkN:A2} and \ref{pf:linkN:A3},
\[
\eta_j(w)=\sum_{\iota\in\mathrm{Loc}(j,w)}\linkNloc(\iota)\le U\sum_{\iota\in\mathrm{Loc}(j,w)}\sum_{\pi\in\linkNgood_\iota}
\linkNmult(\iota,\pi)\le U\dim Z_j(\mu^{\mathrm{tw}}_w)=\linkNdZ(j)\,U .
\]
Multiply by $\linkNlw(w)e(w)\ge0$, sum over $w$ and use \cref{pf:linkN:meanw}. (ii) is \ref{pf:linkN:A6} combined with
\ref{pf:linkN:A5}.
\end{proof}

For a table $\mathrm{tb}$ define $\lambda_G(\ell)=Cx^C\eta_{\mathrm{sz}}$ if $\ell\in\mathrm{gset}$ and $\lambda_G(\ell)=\lambda_B(\ell)$
otherwise, where
\begin{equation}\label{eq:pf:linkN:lamB}
\lambda_B(\ell)=\begin{cases}
Cfmb\log x+Cx^C\Lambda^2,&\ell\in\mathrm{small}\cup\{3\},\\
2Cfmb,&\ell\in\{q,d\},\\
0,&\ell>\ell^*\text{ and }\ell\notin\mathrm{exc},\\
Cfnt^2n_V\Lambda^2+\mathrm{len}_B(\ell),&\text{otherwise},
\end{cases}
\end{equation}
with $\mathrm{len}_B(\ell)\ge0$ an explicit expression in $C$, the ranks and the size quantities of $L$, which distinguishes
split primes and two classes of inert primes. Two estimates proved in the development are used. First, by
\lref{Formalism.MainNotions.TwistedPlaces.loc_le_lamB}, \emph{if $C>0$ and the admitted statements about $BB$, the
numerical bounds, large primes, small primes and $q,d$ hold, then
$\sum_{\lambda\mid\ell}\mathrm{len}(\iota,\lambda)+BB(\iota)\le\lambda_B(\ell_\iota)$ for every $\iota$ with $\ell_\iota\notin\mathrm{gset}$}; in
the last case of \eqref{eq:pf:linkN:lamB} this combines the bound for $BB$ with the numerical bounds and the identity
$\sum_{\lambda\mid\ell}[E_\lambda:\Q_\ell]=2f$. Second, by \lref{UpperTheory.KStar.sum_nuL_le}, \emph{for every finite set
$S$ of primes, $\sum_{\ell\in S}\nu_L(\ell)\log\ell\le\nu_{\mathrm{budget}}(D)$}; this follows from $\nu_L(\ell)\le v_\ell(Q_S)+v_\ell(tr)$
for $\ell\ne q,d$ and $\nu_L(q)=\nu_L(d)=0$ (\lref{UpperTheory.Mu.nuLe} and \lref{UpperTheory.Mu.nuQD}).

\begin{lemma}\label{pf:linkN:smallsums}
For every finite set $S$ of integers $\ge2$, $\sum_{\ell\in S}\log\ell/\ell^2\le2$. For every finite set $S$ of primes and every
integer $n\ge1$, $\sum_{\ell\in S}v_\ell(n)\log\ell\le\log n$.
\formallabel{Formalism.MainNotions.TwistedPlaces.Sums.sum_log_div_sq_le}
\formallabel{Formalism.MainNotions.TwistedPlaces.Sums.sum_padicVal_log_le}
\end{lemma}

\begin{proof}
The first sum is at most $\sum_{n\ge2}\log n/n^2<1$. The second is at most $\sum_{\ell\mid n}v_\ell(n)\log\ell=\log n$.
\end{proof}

\begin{proposition}[Summation over primes]\label{pf:linkN:etaSumle}
Assume $\mathcal A(L,\mathrm{tb},C,P)$ and that $\sum_\ell\mathrm{vv}(D,\ell)$ converges. Then for $1\le j\le m$
\[
\linkNeS(L,j)\le\linkNdZ(j)\Bigl(6+\nu_{\mathrm{budget}}(D)+\log h_E+\log\#G+\sum_{\ell\in S(L)}\log\ell\cdot\lambda_G(\ell)\Bigr)
+\linkNdZ(j)\sum_\ell\mathrm{vv}(D,\ell).
\]
\formallabel{Formalism.MainNotions.TwistedPlaces.etaSum_le}
\end{proposition}

\begin{proof}
By \ref{pf:linkN:A7}, $C$ is at least a positive constant of $P$, so $C>0$. Put
$U(\ell)=\nu_L(\ell)+2\ell^{-2}+\lambda_G(\ell)+VV(\ell)+v_\ell(h_E)+v_\ell(\#G)\ge0$. We claim that
$\eta_\ell(j)\le\log\ell\cdot\linkNdZ(j)\,U(\ell)$ for $\ell\in S(L)$. If $\ell\in\mathrm{gset}$ this follows from
\cref{pf:linkN:etaAt}(ii), since $\nu_L(\ell),v_\ell(\#G)\ge0$. If $\ell\notin\mathrm{gset}$, then for $\iota$ with $\ell_\iota=\ell$ and
good $\pi$, \ref{pf:linkN:A4}, \ref{pf:linkN:A5} and the first estimate above give
\[
\linkNvcc(\iota,\pi)+v_\ell(\#G)\le\nu_L(\ell)+2\ell^{-2}+\lambda_B(\ell)+VV(\ell)+v_\ell(h_E)+v_\ell(\#G)=U(\ell),
\]
and \cref{pf:linkN:etaAt}(i) applies. Summing over $S(L)$: $\sum\log\ell\,(\nu_L(\ell)+2\ell^{-2})\le\nu_{\mathrm{budget}}(D)+4$ by the
second estimate and \cref{pf:linkN:smallsums}; $\sum\log\ell\,VV(\ell)=\sum_{\ell\in S(L)}\mathrm{vv}(D,\ell)\le\sum_\ell\mathrm{vv}(D,\ell)$,
the terms being non-negative and the series convergent; and $\sum\log\ell\,v_\ell(h_E)\le\log h_E$,
$\sum\log\ell\,v_\ell(\#G)\le\log\#G$ by \cref{pf:linkN:smallsums}.
\end{proof}

\begin{definition}\label{pf:linkN:Large}
$x$ is \emph{large for $(\eps,C)$}, written $\mathrm{Lg}(x,\eps,C)$, if $x\ge100$, $8x\le\exp(x^{\eps/3})$,
$4\le\exp(x^\eps)$, $x^{C+2}\exp(x^{\eps/3})\le\exp(x^\eps)$, $48C\le x^{1/2}$, $(C+3)^2(1+\log x)^2\le x^\eps$,
$x^{20C+40}\exp(20x^{\eps/4})\le\exp(x^{\eps/3})$ and $10^8(C+1)^3\le x$.
\formallabel{Formalism.MainNotions.TwistedPlaces.Large}
\end{definition}

\begin{lemma}\label{pf:linkN:largeev}
For $\eps>0$ and $C\in\R$ there is $x_0$ with $\mathrm{Lg}(x,\eps,C)$ for all $x\ge x_0$.
\formallabel{Formalism.MainNotions.TwistedPlaces.large_eventually}
\end{lemma}

\begin{proof}
Each condition holds for large $x$ because $\log x=o(x^a)$ for every $a>0$. For instance,
$x^{C+2}\exp(x^{\eps/3})\le\exp(x^\eps)$ follows from $(C+2)\log x\le x^{\eps/3}\le x^{\eps/3}(x^{2\eps/3}-1)$ once
$x^{2\eps/3}\ge2$, and $x^{20C+40}\exp(20x^{\eps/4})\le\exp(x^{\eps/3})$ from $(20C+40)\log x\le x^{\eps/4}\le
x^{\eps/4}(x^{\eps/12}-20)$ once $x^{\eps/12}\ge40$; the other conditions are similar or immediate.
\end{proof}

Put $U_b=x^{C+2}\exp(x^{\eps/4})$.

\begin{proposition}[Point bound]\label{pf:linkN:pointbound}
Assume $C>0$, $\mathrm{PB}(L,C)$, $\mathrm{Lg}(x,\eps,C)$ and $\mathcal A(L,\mathrm{tb},C,P)$, and let $S$ be a finite set of primes. Then
\[
6+\nu_{\mathrm{budget}}(D)+\log h_E+\log\#G+\sum_{\ell\in S}\log\ell\cdot\lambda_G(\ell)\le2000(C+1)^3\,fmb\Lambda^2.
\]
\formallabel{Formalism.MainNotions.TwistedPlaces.point_bound}
\end{proposition}

\begin{proof}
Under these hypotheses the following estimates hold for every finite set $T$ of primes with the stated property:
\begin{enumerate}[label=(\roman*)]
\item if $T\subseteq\mathrm{gset}$: $\sum_{\ell\in T}\log\ell\cdot Cx^C\eta_{\mathrm{sz}}\le2(C+1)^2U_b^{18}$
(\lref{Formalism.MainNotions.TwistedPlaces.sum_g});
\item if $T\subseteq\mathrm{small}\cup\{3\}$: $\sum_{\ell\in T}\log\ell\,(Cfmb\log x+Cx^C\Lambda^2)\le8Cfmb\Lambda^2$
(\lref{Formalism.MainNotions.TwistedPlaces.sum_23});
\item if $T\subseteq\{q,d\}$: $\sum_{\ell\in T}\log\ell\cdot2Cfmb\le4Cfmb\Lambda^2$
(\lref{Formalism.MainNotions.TwistedPlaces.sum_qd});
\item if $T$ consists of primes $\ell\ne2,3$ with $\ell\le\ell^*$:
$\sum_{\ell\in T}\log\ell\,(Cfnt^2n_V\Lambda^2+\mathrm{len}_B(\ell))\le1000(C+1)^3fmb\Lambda^2$
(\lref{Formalism.MainNotions.TwistedPlaces.sum_small});
\item if $T\subseteq\mathrm{exc}\cap(\ell^*,\infty)$, using the admitted statements on large primes and on $\mathrm{exc}$:
$\sum_{\ell\in T}\log\ell\,(Cfnt^2n_V\Lambda^2+\mathrm{len}_B(\ell))\le1164000(C+1)^2U_b^{18}$
(\lref{Formalism.MainNotions.TwistedPlaces.sum_exc});
\item $6+\nu_{\mathrm{budget}}(D)+\log h_E+\log\#G\le20(C+1)^2U_b^{18}$ (\lref{Formalism.MainNotions.TwistedPlaces.glob_le});
\item $4\cdot10^6(C+1)^2U_b^{18}\le b$ (\lref{Formalism.MainNotions.TwistedPlaces.big_le_b}).
\end{enumerate}
By \eqref{eq:pf:linkN:lamB}, for every prime $\ell$ the number $\log\ell\cdot\lambda_G(\ell)$ is at most the sum of the following
non-negative terms: $\log\ell\cdot Cx^C\eta_{\mathrm{sz}}$ if $\ell\in\mathrm{gset}$; $\log\ell\,(Cfmb\log x+Cx^C\Lambda^2)$ if
$\ell\in\mathrm{small}\cup\{3\}$; $\log\ell\cdot2Cfmb$ if $\ell\in\{q,d\}$; and $\log\ell\,(Cfnt^2n_V\Lambda^2+\mathrm{len}_B(\ell))$ if
$\ell\ne2,3$ and $\ell\le\ell^*$, or if $\ell\in\mathrm{exc}$ and $\ell>\ell^*$. Indeed, if none of the first three cases applies
($2$ is small), then either $\ell>\ell^*$ and $\ell\notin\mathrm{exc}$, where $\lambda_G(\ell)=0$, or one of the last two cases
applies. Summing over $S$ and applying (i)--(vi) to the five corresponding subsets of $S$, the left-hand side is at most
\[
(20+2+1164000)(C+1)^2U_b^{18}+8CX+4CX+1000(C+1)^3X,\qquad X=fmb\Lambda^2 .
\]
Since $1\le m\le f$ and $\Lambda\ge1$, $b\le X$, so by (vii) the first term is at most $X$; and $1+12C+1000(C+1)^3\le2000(C+1)^3$
for $C\ge0$.
\end{proof}

\begin{proposition}[Export bound]\label{pf:linkN:export}
Assume $\mathrm{Lg}(x,\eps,C)$ and that $\linkNfin(D)$ implies the existence of a table $\mathrm{tb}$ with
$\mathcal A(L,\mathrm{tb},C,P)$. If $\linkNfin(D)$ holds, then $\mathrm{LN}(L,j,2000(|C|+1)^3)$ holds for $1\le j\le m$.
\formallabel{Formalism.MainNotions.TwistedPlaces.export_bound_of_admissible}
\end{proposition}

\begin{proof}
Let $\mathrm{tb}$ be the table. By \ref{pf:linkN:A7}, $C>0$, so $|C|=C$. Combine \cref{pf:linkN:etaSumle} with
\cref{pf:linkN:pointbound} for $S=S(L)$, and use $\linkNdZ(j)\ge0$ and $(\log x)^2N_c\ge0$.
\end{proof}

\begin{theorem}[Summation]\label{pf:linkN:sumthm}
Let $\mathcal F\in\sfname{FamP}(\eps)$ satisfy the local bounds with constant $C$. Then $\mathcal F$ satisfies $\sfname{LinkN}$ with
constant $2000(|C|+1)^3$.
\formallabel{GlobalTheory.T.linkN_along_family_of_localBounds}
\formallabel{GlobalTheory.LinkFam.main}
\formallabel{Formalism.MainNotions.TwistedPlaces.stmtU_holds}
\formallabel{Formalism.MainNotions.TwistedPlaces.stmtF_holds}
\end{theorem}

\begin{proof}
By \cref{pf:linkN:admitted}, for $x\ge x_0$ with $x>x_1$ and every $i\in I_x$, $\linkNfin(D_{x,i})$ implies the existence of a
table with $\mathcal A(L_{x,i},\cdot,C,P_{x,i})$. Since the scales of the members have exponent $\eps>0$,
\cref{pf:linkN:largeev} gives $x_L$ with $\mathrm{Lg}(x,\eps,C)$ for $x\ge x_L$. For $x\ge\max(x_0,x_1+1,x_L)$ the size data
$L_{x,i}$ have scale $(x,\eps)$, and \cref{pf:linkN:export} gives $\mathrm{LN}(L_{x,i},j,2000(|C|+1)^3)$ whenever
$\linkNfin(D_{x,i})$, for $1\le j\le m$. This is $\sfname{LinkN}$ with threshold $\max(x_0,x_1+1,x_L)$.
\end{proof}

% ---------------------------------------------------------------------------
\subsection{Uniform constants: the commuting endomorphism and paired character conditions}\label{sec:pf:linkN:extra}

\begin{lemma}[Integrality of central idempotents]\label{pf:linkN:cenidem}
Let $\Gamma$ be a finite group of order $n$ and $\zeta\in\C$ a primitive $n$-th root of unity. If an endomorphism $\phi$ of a
finite-dimensional complex vector space satisfies $\phi^{n+1}=\phi$, then $\tr\phi\in\Z[\zeta]$. If $e\in\C[\Gamma]$ is a central
idempotent, then $n\cdot e(h)\in\Z[\zeta]$ for every $h\in\Gamma$, where $e(h)$ is the coefficient of $h$.
\formallabel{GlobalTheory.CentralIdem.trace_mem_adjoin}
\formallabel{GlobalTheory.CentralIdem.centralIdemCoeff}
\end{lemma}

\begin{proof}
Every eigenvalue $\mu$ of $\phi$ satisfies $\mu^{n+1}=\mu$, so $\mu=0$ or $\mu$ is an $n$-th root of unity, a power of $\zeta$; the
trace is the sum of the eigenvalues with multiplicity. For the second statement let $y=h^{-1}e$. Since $e$ is central and
idempotent and $h^{-n}=1$, $y^{n+1}=h^{-(n+1)}e=y$, so the trace of left multiplication by $y$ on $\C[\Gamma]$ lies in $\Z[\zeta]$. In
the basis $\Gamma$ all diagonal entries of this operator equal the coefficient $y(1)=e(h)$, so the trace is $n\,e(h)$.
\end{proof}

\begin{lemma}[Central lifts]\label{pf:linkN:centralelt}
Let $D$ be a datum with a pending datum $Q$, $\iota$ a local index, and $s=(s(k,v))$ a family of labels. There is a central element
$z$ of $\OO_\iota[G]$ with $\iota_\OO(z)=\#G\cdot\prod_{(k,v)}e_{\mathrm{Ser}}(k,v,s(k,v))$. More generally, if $y\in\C[G]$ is
central and all its coefficients lie in $\iota_\OO(\OO_\iota)$, then $y=\iota_\OO(z)$ for a central $z\in\OO_\iota[G]$.
\formallabel{GlobalTheory.Pass.CentralStep.main}
\formallabel{GlobalTheory.CentralIdem.card_mul_idem_from_center}
\formallabel{GlobalTheory.CentralIdem.exists_central_lift}
\end{lemma}

\begin{proof}
For the second statement let $z$ be the coefficientwise preimage of $y$, which exists and is unique since $\iota_\OO$ is injective.
The coefficientwise map $\OO_\iota[G]\to\C[G]$ is an injective ring homomorphism, and $\iota_\OO(zb)=y\,\iota_\OO(b)=\iota_\OO(b)\,y=\iota_\OO(bz)$
for $b\in\OO_\iota[G]$, so $z$ is central. For the first statement, the factors $e_{\mathrm{Ser}}(k,v,s(k,v))$ are central
idempotents of $\C[G]$ lying in different factors, so their product $e$ is a central idempotent. By \eqref{eq:pf:linkN:G}, $\#G$ is
the order of $G$, and $\iota_\OO(\zeta_\iota)$ is a primitive $\#G$-th root of unity. By \cref{pf:linkN:cenidem} the coefficients of
$\#G\cdot e$ lie in $\Z[\iota_\OO(\zeta_\iota)]\subseteq\iota_\OO(\OO_\iota)$, and the second statement applies to $y=\#G\cdot e$.
\end{proof}

\begin{corollary}\label{pf:linkN:goodelt}
In the situation of \cref{pf:linkN:centralelt} there is a central element $z$ of $\OO_\iota[G]$ with $\iota_\OO(z)=\#G\cdot e_0$.
\formallabel{UpperTheory.GalC.exists_unit_one_add_generator}
\end{corollary}

\begin{proof}
Take the sum of the elements of \cref{pf:linkN:centralelt} over the good families $s$.
\end{proof}

The following lemma is proved in \lref{GlobalTheory.PassOp.onMt_e0_eq_self_of_mem_zgood_and_eq_zero_of_notMem}.

\begin{lemma}[Type vectors]\label{pf:linkN:typevec}
Assume (Ser), (Adm) and (Perm), and let $Q$ be a pending datum. Let $x\in Z_\iota$, let $P$ be an automorphic point with
eigensystem $x$ and $F\in\mathrm{Typ}(P)$. Then $e_0\cdot F=F$ if $x\in\linkNgood_\iota$, and $e_0\cdot F=0$ if $x\notin\linkNgood_\iota$.
\end{lemma}

The mechanism is as follows: each type vector of $P$ lies in a subspace $M^t_\iota(\sigma')$ for a parameter $\sigma'$ matching
the local parameters of $x$; by (Ser) the factors of $e_0$ act on $M^t_\iota(\sigma')$ through the labels of $\sigma'$, so that $e_0$
acts as the identity if $\sigma'$ is good and as $0$ otherwise; and by the definition of the data, $x\in\linkNgood_\iota$ if and
only if $x\in Z_\iota$ has a good parameter, which by (Perm) does not depend on the automorphic point over $x$.

\begin{proposition}[The operator]\label{pf:linkN:operator}
Assume the hypotheses of \cref{pf:linkN:typevec} and (Eig) for $Q$. For every local index $\iota$ there is
$\alpha\in\End_{\OO_\iota}(M_\iota)$ commuting with $T_\iota$ such that $\alpha_K$ is multiplication by $\#G$ on $M_K[x]$ for
$x\in\linkNgood_\iota$, and $\alpha_K=0$ on $M_K[x]$ for every $x$ such that $e_0$ kills all type vectors of all automorphic points
with eigensystem $x$; in particular $\alpha_K=0$ on $M_K[x]$ for $x\in Z_\iota\smallsetminus\linkNgood_\iota$. Consequently $D$ satisfies
the \emph{commuting endomorphism condition} $\mathrm{CE}(D)$: at every local index $\iota$ with $\ell_\iota\notin\mathrm{gset}$
there is an endomorphism of $M_\iota$ commuting with $T_\iota$ which acts by $\#G$ on the eigenspaces of $\linkNgood_\iota$ and by
$0$ on those of $Z_\iota\smallsetminus\linkNgood_\iota$.
\formallabel{GlobalTheory.PassOp.main_alt_alt}
\formallabel{GlobalTheory.T.commutingEndomorphismCondition_of_serLoc}
\formallabel{Formalism.MainNotions.N.KillingConditions.CommutingEndomorphismCondition}
\end{proposition}

\begin{proof}
Let $z$ be as in \cref{pf:linkN:goodelt} and $\alpha=z_Q$, which commutes with $T_\iota$. Since $e_0$ is central,
$\iota_\OO(z)\cdot F=\#G\,(e_0\cdot F)$ for every $F$. If $x\in\linkNgood_\iota\subseteq Z_\iota$, then by \cref{pf:linkN:typevec}
$\iota_\OO(z)\cdot F=\iota_\OO(\#G)F$ for all type vectors over $x$, and (Eig) with $c=\#G$ gives $\alpha_K=\#G$ on $M_K[x]$. If $e_0$
kills all type vectors over $x$, then $\iota_\OO(z)\cdot F=0$, and (Eig) with $c=0$ gives $\alpha_K=0$ on $M_K[x]$; by
\cref{pf:linkN:typevec} this applies to $x\in Z_\iota\smallsetminus\linkNgood_\iota$.
\end{proof}

\paragraph{The paired character condition.}
At the type primes the bound for $\linkNloc$ is not obtained from congruence modules but from the following condition.

\begin{definition}[Paired character condition]\label{pf:linkN:PCC}
For $C\in\R$ and $C_T\in\N$, the \emph{paired character condition} $\mathrm{PC}(D,C,C_T)$ is the statement: for every local index
$\iota$ with $\ell_\iota=g_k$ ($k\in\{0,1\}$, and $k'$ the other index), all families $r_\pi:G_E\to\GL_{n_V}(\OO_\iota)$ and
$\rho_\pi:G_E\to\GL_r(\OO_\iota)$ ($\pi\in X_\iota$) such that $r$ is good and $r_\pi$ has a semisimple decomposition in terms of $\rho_\pi$
for $\pi\in Z_\iota$, and all integers $\eta,s,S\ge0$, signs $\mathrm{sg}$ and finite sets $U_v$ of integers indexed by the places of
$E$ satisfying
\[
\eta+t+1<s,\qquad\log g_{k'}\le2\log g_k,\qquad C_T\bigl((s+S)(n_V+1)^2+(n_V+1)^3\bigr)<g_k,
\]
a sign condition and a host condition on $(r,\eta,s,S,\mathrm{sg},U)$: if the Selmer groups of the $\rho_\pi$, $\pi\in\linkNgood_\iota$, are
finite and $\Sha^1(G_{E,S_{g_k}},\Q_{g_k}/\Z_{g_k}(-e))$ is finite for $0<e<t$, then
\[
\begin{aligned}
\linkNloc(\iota)\le\sum_{\pi\in\linkNgood_\iota}\linkNmult(\iota,\pi)\Bigl(&4fn_V\,Cn_V^2(n_V+\eta+t)+2fr^2+\linkNs(\iota,\rho,\pi)+VV(g_k)\\
&+v_{g_k}(h_E)+4n_Vv_\delta(\iota)+t\sum_{v\in Q_\iota}\lambda_\iota\bigl(\OO_\iota/(\zeta_v-1)\bigr)\Bigr),
\end{aligned}
\]
where $v_\delta(\iota)=\sum_{0<e<t}(s_0(\iota)+v_{g_k}(e))$ with the local level exponent $s_0(\iota)$, and $\zeta_v=\chi_{\iota,1}(\gamma_{\iota,v})$.
\formallabel{Formalism.MainNotions.PairedCharacterCondition.PairedCharacterConditionAlt}
\end{definition}

\begin{definition}\label{pf:linkN:extraDef}
For $C\in\R$ and $C_T\in\N$ let $\mathrm{E}(D,C,C_T)$ be the conjunction of $\mathrm{CE}(D)$ and $\mathrm{PC}(D,C,C_T)$. A family
$\mathcal F$ \emph{has the facts with constants $(C,C_T)$} if every member with $x>x_1$ satisfies $\mathrm{NF}(D_{x,i},\mathrm{nch}_{x,i})$
and $\mathrm{E}(D_{x,i},C,C_T)$.
\formallabel{GlobalTheory.NLink.ExtraAlt}
\formallabel{GlobalTheory.NLink.FactsFAlt}
\end{definition}

\begin{proposition}[Paired character condition from local inputs]\label{pf:linkN:closinglocal}
Assume global Artin reciprocity, the Bhatt--Scholze comparison, Fontaine's $B_{\mathrm{cris}}$, the local Bloch--Kato theory,
local Tate duality and the local Euler characteristic formula (all part of \cref{pf:linkN:H2}). Let $D$ be a datum with local
Hodge--Witt, Fontaine and comparison data, let $C_T\ge40$ and $C\ge8C_T$, and assume that $D$ has a pending datum $Q$ with (TA)
and (Eig), and satisfies (Adm), (Perm), (AP), (EV), irreducibility of the residual representations of rank $r$, (Ser), the type-prime
wild idempotent choice, (Ine) and (Chi). Then $\mathrm{PC}(D,C,C_T)$ holds.
\formallabel{GlobalTheory.ClosingLocalInputs.main_alt_alt}
\formallabel{GlobalTheory.T.conclusion41_of_pAdicHodge_alt}
\end{proposition}

\begin{proof}
The following statements about $D$ are established first:
\begin{enumerate}[label=(\roman*)]
\item by \cref{pf:linkN:thmT}, the local theory of $D$ satisfies the crystalline period bound $\mathcal T(K',L,C_T)$ at every
prime (the development calls it Theorem~T), as $C_T\ge40$;
\item by \lref{GlobalTheory.T.local_absorption_length_bound_of_blochKato}, a local pushout proposition (two length bounds of
Bloch--Kato duality type) holds, by the local Bloch--Kato theory, local duality and the Euler characteristic formula;
\item by \lref{GlobalTheory.Pass.residualRep_spans_matrixAlgebra_of_goodPoint}, the residual representations at good points
are absolutely irreducible (their images span the matrix algebra), by residual irreducibility;
\item by \lref{GlobalTheory.typePrimeCongruences_forall}, the congruence lemmas at the type primes (for good families: a
characteristic-polynomial lemma, a wild idempotent choice, a vanishing lemma, a non-vanishing bound for quotient homomorphisms,
and a coboundary lemma) hold, by (Adm), (Perm), residual irreducibility and the wild idempotent choice;
\item by \lref{GlobalTheory.T.theorem_tb_of_local_blochKato_duality_euler} and
\lref{GlobalTheory.T.blochKatoLocal_all_of_tateDuality_eulerChar}, two further local statements about the local theory of $D$
and its Fontaine data hold, by the local Bloch--Kato theory, local duality and the Euler characteristic formula;
\item the Galois family conditions, which are part of (Perm).
\end{enumerate}
The conclusion is then \lref{GlobalTheory.T.admT_of_artin_reciprocity_and_local_conditions}: from global Artin reciprocity
one obtains a sixteen-step statement (\lref{GlobalTheory.T.no_cut_passes_of_artin_reciprocity}), which comprises the two
statements of \cref{pf:linkN:cenidem,pf:linkN:centralelt}, local counts of direct and dual liftings, potentially crystalline type
families with a common tame extension, a lifting step, residues and a global statement; together with $C\ge8C_T$, (i)--(vi) and the
hypotheses, it yields a congruence statement at the type primes, an $\eta$-bound for the cut at the type locus, and finally the
inequality of \cref{pf:linkN:PCC}.
\end{proof}

\begin{lemma}[Comparison data]\label{pf:linkN:wach}
Assume the integral comparison theorem of Bhatt--Scholze type for $\Zl$-lattices in crystalline representations contained in
$\sfname{Hypotheses}_2$. Then every datum admits local Hodge--Witt data, Fontaine data and comparison data for its local theory.
\formallabel{GlobalTheory.T.exists_wachFunctor_comparison_of_bhattScholze}
\formallabel{GlobalTheory.HWL.exists_hw_fon_cmp_of_integral}
\end{lemma}

\begin{proof}
By \lref{GlobalTheory.T.bhattScholze_coefficients_of_padicInt_lattices}, the comparison statement for $\Zl$-lattices implies the
corresponding statement for lattices over the rings of integers of finite extensions of $\Ql$. By
\lref{GlobalTheory.HWL.exists_hw_fon_cmp_ties}, the latter provides, for every local theory with $A_{\inf}$-data satisfying the tie
conditions, local Hodge--Witt data, Fontaine data and comparison data (together with further compatibilities, which are not used here). The
local theory of a datum, with its $A_{\inf}$-data and its ties, is such a local theory.
\end{proof}

\begin{theorem}[Uniform constants]\label{pf:linkN:extra}
Assume $\sfname{Hypotheses}_2$. Put $C_T=40$ and $C=8C_T=320$. Every datum $D$ which satisfies $\mathrm{NF}(D,\mathrm{nch})$ for some
norm choice and $\sfname{Tie4}(D)$ satisfies $\mathrm{E}(D,C,C_T)$.
\formallabel{GlobalTheory.ExtraFromFacts.main_alt_alt}
\formallabel{GlobalTheory.T.exists_const_tieChi_bound_of_facts}
\formallabel{GlobalTheory.T.exists_bound_tieChi_of_prints}
\formallabel{GlobalTheory.NormFromFacts.main_alt_alt}
\end{theorem}

\begin{proof}
$\sfname{Tie4}(D)$ provides (Chi), a pending datum $Q$ with (TA) and (Eig), (Ine) and (Ser). The norm-choice facts contain (Adm),
(Perm), (AP), (EV), residual irreducibility and the duties, which contain the type-prime wild idempotent choice.

\emph{Paired character condition.} By \cref{pf:linkN:wach}, the
Bhatt--Scholze comparison provides every datum with local Hodge--Witt, Fontaine and comparison data. \Cref{pf:linkN:closinglocal}
with $C_T=40$ and $C=320$ gives $\mathrm{PC}(D,320,40)$.

\emph{Commuting endomorphism condition.} \Cref{pf:linkN:operator} applies with the pending datum $Q$, since (Ser), (Adm), (Perm) and
(Eig) are available. (In the formal proof the central elements of \cref{pf:linkN:centralelt} are taken from the first two
components of the sixteen-step statement.)
\end{proof}

% ---------------------------------------------------------------------------
\subsection{The crystalline period bound}\label{sec:pf:linkN:thmT}

This subsection gives the top level of the proof of the crystalline period bound used in \cref{pf:linkN:closinglocal}(i) (the
development calls it Theorem~T). Its main ingredient, a lifting theorem for crystalline lattices of tame type proved with
Breuil--Kisin modules, and three further steps are cited.

\begin{definition}[Setting]\label{pf:linkN:Tsetting}
Fix a prime $\ell$ and let $G_\ell=\Gal(\Qlbar/\Ql)$. The local theory of a datum provides period data at $\ell$ (Fontaine's rings
$B_{\mathrm{cris}}\subseteq B_{\mathrm{dR}}$ with their actions of $G_\ell$, Frobenius and filtration); relative to them a continuous
representation $V$ of $G_\ell$ over a finite extension $L$ of $\Ql$ is \emph{crystalline}, its \emph{Hodge--Tate jumps} at an embedding $\tau$ are
defined, and so are the \emph{finite} (Bloch--Kato) cocycles $G_\ell\to V$. Let $K'/\Ql$ be a finite Galois extension whose maximal unramified
subextension has $G_\ell$-stable image in $B_{\mathrm{cris}}$, $G_{K'}=\Gal(\Qlbar/K')$, and $I\subseteq G_\ell$ the inertia group. $V$ is
\emph{potentially crystalline} if $V|_{G_{K'}}$ is crystalline; then $I$ acts on $D_{\mathrm{pcris}}(V)$. For a finite set $\Sigma$ of characters
$I\to L^\times$, \emph{inertia acts through $\Sigma$} on $V$ if every $\chi\in\Sigma$ is trivial on $I\cap G_{K'}$ and of order prime to $\ell$, and
$D_{\mathrm{pcris}}(V)$ is the sum of the $\chi$-eigenspaces of $I$, $\chi\in\Sigma$. For a Frobenius element $\mathrm{Fr}\in G_\ell$ and a character
$\chi$ of $I$, $\mathrm{orb}(\chi)$ is the least $k\ge1$ with $\chi(\mathrm{Fr}^{-k}g\,\mathrm{Fr}^k)=\chi(g)$ for all $g\in I$, and
$r'(\Sigma)=\max_{\chi,\psi\in\Sigma}\mathrm{orb}(\chi\psi^{-1})$.

For $G_\ell$-stable $\OO_L$-lattices $T_i\subseteq V_i$, the \emph{lattice finite cocycles} of $T_i$ are the finite cocycles $G_\ell\to V_i$ with
values in $T_i$. For a surjective homomorphism of lattices $F:T_1\to T_2$ and $e\in\N$, $\mathrm{Lift}_F(e)$ is the statement: for every lattice
finite cocycle $c$ of $T_2$, $\ell^ec$ is the sum of the image under $F$ of a lattice finite cocycle of $T_1$ and a coboundary with values in
$T_2$.
\formal{Formalism.WeaklyAdmissibleFiltrations.PotCrysData}
\formal{Formalism.WeaklyAdmissibleFiltrations.LatHom.LiftF}
\formal{Formalism.MainNotions.ElementDenominator.rPrime}
\end{definition}

\begin{definition}\label{pf:linkN:Tstmt}
For $K'$ as above, a finite extension $L$ of $\Ql$ and $C\in\N$, $\mathcal T(K',L,C)$ is the statement: let $T_1,T_2$ be lattices in
potentially crystalline $L$-representations, $F:T_1\to T_2$ a surjective homomorphism of lattices, $N,\eta,H\in\N$, $U\subseteq\Z$ finite,
$\Sigma$ a finite set of characters of $I$ and $\mathrm{Fr}$ a Frobenius element; if $\Sigma$ is stable under conjugation by $\mathrm{Fr}$,
$1\in\Sigma$, $\#\Sigma\le N$, $N\ge2$, inertia acts through $\Sigma$ on $V_1$ and $V_2$, $U\subseteq[-\eta,H]$, $0\in U$, $\#U\le N+1$,
$1\le\eta<H$, the Hodge--Tate jumps of $V_1$ and $V_2$ at every embedding lie in $U$, and $C(HN^2+N^3)<\ell$, then
\[
4r'(\Sigma)\le N^2\quad\text{and}\quad\mathrm{Lift}_F\bigl(C\,(N\,r'(\Sigma)+\eta\,(N+r'(\Sigma)))\bigr).
\]
\formal{Formalism.MainNotions.ElementDenominator.thmT}
\formal{Formalism.MainNotions.ElementDenominator.CorollaryALocal}
\end{definition}

\begin{lemma}\label{pf:linkN:liftmono}
If $\mathrm{Lift}_F(e)$ holds and $e\le e'$, then $\mathrm{Lift}_F(e')$ holds.
\formallabel{GlobalTheory.TheoremTClosing.liftF_of_le}
\end{lemma}

\begin{proof}
$\ell^{e'}c=\ell^{e'-e}(\ell^ec)$, and the set of sums in the definition is an $\OO_L$-module.
\end{proof}

\begin{lemma}[Orbit lengths]\label{pf:linkN:orbits}
Let $a\ge1$ and $e=\ell^a-1$, and for $c\in\Z/e$ let $\mathrm{per}(c)$ be the least period of $c$ under $x\mapsto\ell x$. Let $\Sigma'\subseteq\Z/e$
be stable under multiplication by $\ell$, with $0\in\Sigma'$, $\#\Sigma'\le N$ and $N\ge2$, and let $m=\max_{c,c'\in\Sigma'}\mathrm{per}(c-c')$.
Then $1\le m$, $4m\le N^2$, and $\mathrm{per}(c)\le m$ for $c\in\Sigma'$.
\formallabel{GlobalTheory.T.maxFrobeniusOrbitLengthDiff_bounds_zmod}
\formallabel{GlobalTheory.BK.maxPeriod_bounds}
\end{lemma}

\begin{proof}
As $\ell^a=1$ in $\Z/e$, multiplication by $\ell$ is a permutation of $\Z/e$ and every $c$ is periodic, with $\mathrm{per}(c)\ge1$ equal to the
length of its orbit; this gives $1\le\mathrm{per}(0)\le m$, and $\mathrm{per}(c)=\mathrm{per}(c-0)\le m$ for $c\in\Sigma'$. Let $c,c'\in\Sigma'$, with
periods $p,q$. Since multiplication by $\ell$ is additive, $\ell^{\mathrm{lcm}(p,q)}(c-c')=c-c'$, so $\mathrm{per}(c-c')\le\mathrm{lcm}(p,q)$. The orbits
of $c$ and $c'$ lie in $\Sigma'$, and the orbit of $0$ is $\{0\}$. If $c=0$ or $c'=0$, or if $c$ and $c'$ have the same orbit, then
$\mathrm{per}(c-c')\le\max(p,q)\le\max(1,N-1)$; otherwise the two orbits are disjoint subsets of $\Sigma'\smallsetminus\{0\}$, so $p+q\le N-1$ and
$\mathrm{lcm}(p,q)\le pq\le(N-1)^2/4$. In all cases $4\,\mathrm{per}(c-c')\le N^2$, since $4(N-1)\le N^2$.
\end{proof}

\begin{lemma}[Coefficients]\label{pf:linkN:tamecoeff}
Let $a\ge1$ with $\ell\nmid a$, put $e=\ell^a-1$, and let $L$ be a finite extension of $\Ql$.
\begin{enumerate}[label=(\roman*)]
\item The compositum $L'$ of $L$ with the unramified extension of $\Ql$ of degree $a$ is a finite extension of $\Ql$ containing both.
\item If $L$ contains the unramified extension of degree $a$, then $L$ contains a primitive $e$-th root of unity, so that $L$ carries
\emph{tame descent data} of exponent $(e,a)$ (namely $e+1=\ell^a$, $\ell\nmid a$ and such a root of unity).
\end{enumerate}
\formallabel{GlobalTheory.T.exists_coeffField_containsUnramifiedWitt_tameDescent}
\formallabel{GlobalTheory.TameCoeff.main}
\end{lemma}

\begin{proof}
(i) is clear. (ii) The unramified extension of degree $a$ is generated by a primitive $(\ell^a-1)$-th root of unity
\cite[Ch.~II, \S7]{Neukirch1999}.
\end{proof}

\begin{proposition}[Lifting step]\label{pf:linkN:liftstep}
Assume the properties of $B_{\mathrm{cris}}$ and the comparison theorem of Bhatt--Scholze type for lattices over the rings of integers of finite
extensions of $\Ql$. Let $D$ be a datum with local Hodge--Witt data and Fontaine data, and $C\ge40$ an integer. If the admitted local statements
of $D$ hold (see the proof of \cref{pf:linkN:thmT}), then $\mathcal T(K',L,C)$ holds for the period data of $D$ at every prime $\ell$, for every
$K'$ as above and every finite extension $L$ of $\Ql$.
\formallabel{GlobalTheory.T.liftingStep_theoremT_of_bhattScholze_fontaine}
\formallabel{GlobalTheory.TheoremTClosing.main}
\end{proposition}

\begin{proof}
Fix $\ell$, $K'$ and $L$; put $a=(\ell-1)!$ and $e=\ell^a-1$, so $\ell\nmid a$.

\emph{Coefficients.} Let $L'$ be the compositum of $L$ with the unramified extension of $\Ql$ of degree $a$; by \cref{pf:linkN:tamecoeff}, $L'$
carries tame descent data of exponent $(e,a)$. By \lref{GlobalTheory.T.theoremT_descend_smaller_coefficients}, $\mathcal T(K',L',C)$ implies
$\mathcal T(K',L,C)$ for $L\subseteq L'$; so we may replace $L$ by $L'$. Fix a compatible system $\varpi=(\varpi_n)$ in $\Qlbar$ with $\varpi_0^e=-\ell$ and
$\varpi_{n+1}^\ell=\varpi_n$, and let $K_\varpi$ be the field generated by $\varpi_0$ over the unramified extension of degree $a$.

\emph{Data.} Let $T_1,T_2,F,N,\eta,H,U,\Sigma,\mathrm{Fr}$ satisfy the hypotheses of $\mathcal T(K',L',C)$. Since $C\ge40$,
$40(HN^2+N^3)<\ell$; in particular $N<\ell$. The characters of $\Sigma$ have order prime to $\ell$ and are trivial on $I\cap G_{K'}$.

\emph{Tame field.} By \lref{GlobalTheory.T.exists_tameField_goodLattice_potentiallyCrystalline} (which uses the properties of
$B_{\mathrm{cris}}$), there is $\Sigma'\subseteq\Z/e$, stable under multiplication by $\ell$, with $0\in\Sigma'$, $\#\Sigma'\le N$ and
$\max_{c,c'\in\Sigma'}\mathrm{per}(c-c')=r'(\Sigma)$, such that every lattice $T$ in a potentially crystalline representation on which inertia acts
through $\Sigma$ becomes crystalline over $K_\varpi$ and is \emph{of tame type $\Sigma'$}: for $g\in I$ and $w\in L'$ with $g(\varpi_0)=w\varpi_0$,
$\prod_{c\in\Sigma'}(g-w^c)$ kills $D_{\mathrm{cris}}(V|_{\Gal(\Qlbar/K_\varpi)})$.

\emph{Lifting.} The numerical conditions ($\Sigma'$ as above, $N\ge2$, $U\subseteq[-\eta,H]$, $0\in U$, $\#U\le N+1$, $1\le\eta<H$ and
$40(HN^2+N^3)<\ell$) hold, and $T_1,T_2$ are of tame type $\Sigma'$ with Hodge--Tate jumps in $U$. The lifting theorem for good crystalline
lattices, \lref{GlobalTheory.T.lift_surjection_good_crystalline_lattices} (proved with Breuil--Kisin modules with tame descent, from the
comparison theorem and the properties of $B_{\mathrm{cris}}$), gives $\mathrm{Lift}_F\bigl(7(Nr'+\eta(N+r'))\bigr)$ with $r'=r'(\Sigma)$.

\emph{Conclusion.} By \cref{pf:linkN:orbits}, $4r'(\Sigma)\le N^2$, and by \cref{pf:linkN:liftmono} with $7\le C$,
$\mathrm{Lift}_F\bigl(C(Nr'+\eta(N+r'))\bigr)$ holds. The admitted local statements are not used.
\end{proof}

\begin{theorem}[Crystalline period bound]\label{pf:linkN:thmT}
Assume the hypotheses of $\sfname{Hypotheses}_2$ on $B_{\mathrm{cris}}$, on the comparison theorem of Bhatt--Scholze type for $\Zl$-lattices, on the
local Bloch--Kato theory, on local Tate duality and on the local Euler characteristic formula. Let $D$ be a datum with local Hodge--Witt data,
Fontaine data and comparison data, and $C\ge40$ an integer. Then $\mathcal T(K',L,C)$ holds for the period data of $D$ at every prime $\ell$, every
$K'$ as above and every finite extension $L$ of $\Ql$.
\formallabel{GlobalTheory.T.crystallinePeriodBound_all_constant_ge_forty}
\formallabel{GlobalTheory.InertiaEigenvalues.theoremT_all_records}
\end{theorem}

\begin{proof}
By \lref{GlobalTheory.T.bhattScholze_coefficients_of_padicInt_lattices}, the comparison theorem for $\Zl$-lattices implies the version for
lattices over the rings of integers of finite extensions of $\Ql$. The \emph{admitted local statements} of $D$ (five statements about the
local theory of $D$, one of which involves its Hodge--Witt and Fontaine data, and the surjectivity of an Artin--Schreier--Witt
map) follow from the local Bloch--Kato
theory, local duality and the Euler characteristic formula by \lref{GlobalTheory.T.localLifting_admittedHypotheses_all}. Now
\cref{pf:linkN:liftstep} applies.
\end{proof}

% ---------------------------------------------------------------------------
\subsection{Good cuts}\label{sec:pf:linkN:goodcuts}

\paragraph{Vanishing at auxiliary places.}
The following lemma is the classical vanishing of local cohomology at Taylor--Wiles type places
\cite[\S2]{DarmonDiamondTaylor1997}.

\begin{lemma}\label{pf:linkN:tamevanish}
Assume the duality and Euler characteristic statements of \cref{pf:linkN:Hdec}. Let $F$ be a number field, $\ell$ a prime, $v$ a
finite place of $F$ with $v\nmid\ell$, $A$ a valuation ring of $\Omega$ above $v$ with decomposition group $D_A$ and inertia group
$I_A$, and $M$ a finite discrete $\F_\ell$-module with a continuous action $\rho$ of $D_A$. If there is $\gamma\in I_A$ such that
$\rho(\gamma)-1$ is bijective on $M$, then $H^0(D_A,M)=H^1(D_A,M)=H^2(D_A,M)=0$.
\formallabel{GlobalTheory.StringRepresentationDecomposition.tame_vanish}
\formallabel{GlobalTheory.T.localCohomology_vanishes_tame_inertia_no_fixed}
\end{lemma}

\begin{proof}
If $m$ is invariant, then $(\rho(\gamma)-1)m=0$, so $m=0$. Let $\chi$ be the cyclotomic character modulo $\ell$ restricted to $D_A$;
since $v\nmid\ell$, the $\ell$-th roots of unity generate an extension unramified at $v$, so $\chi(\gamma)=1$. If $\varphi\in M^\vee(1)$
is invariant, then $\varphi(m)=(\gamma\varphi)(m)=\chi(\gamma)\varphi(\rho(\gamma)^{-1}m)=\varphi(\rho(\gamma)^{-1}m)$ for all $m$;
replacing $m$ by $\rho(\gamma)m$ gives $\varphi((\rho(\gamma)-1)m)=0$, and as $\rho(\gamma)-1$ is surjective, $\varphi=0$. Thus
$H^0(D_A,M^\vee(1))=0$, and the injectivity of $H^2(M)\to\Hom(H^0(M^\vee(1)),H^2(\mu_\ell))$ gives $H^2(D_A,M)=0$. Finally the Euler
characteristic formula for $v\nmid\ell$ gives $\#H^1=\#H^0\cdot\#H^2=1$.
\end{proof}

For a local index $\iota$ let $k_\iota$ be the residue field of $\OO_\iota$, $\bar\chi_1$ the reduction of $\chi_{\iota,1}$, and $\bar\eps$
the reduction of the cyclotomic character with values in $\OO_\iota^\times$. A \emph{residual representation} of $\iota$ is a
continuous representation $\bar\rho:G_E\to\GL_r(k_\iota)$ of the kind specified by the datum.

\begin{lemma}[Vanishing for residual representations]\label{pf:linkN:h12}
Assume \cref{pf:linkN:Hdec}. Let $\iota$ be a local index, $\bar\rho$ a residual representation of $\iota$ and $w\in S_\iota$, and put
$\gamma=\gamma_{\iota,w}$. If $\bar\chi_1(\gamma)$ is not an eigenvalue of $\bar\rho(\gamma)$, then for every $e\in\Z$ the restriction to
the decomposition group at $w$ of $\bar\rho\otimes\bar\chi_1^{-1}\bar\eps^{\,e}$ has vanishing $H^1$ and $H^2$.
\formallabel{GlobalTheory.HOneTwoLevel.hOneTwo_holds}
\formallabel{GlobalTheory.T.hOneTwo_of_localFiniteness_tateDuality_eulerChar}
\formallabel{Formalism.MainNotions.N.CohomologyScalarsVanish.cohScalarsVanish}
\end{lemma}

\begin{proof}
Let $M=k_\iota^r$ with the twisted action $\rho'$, viewed as an $\F_\ell$-module. The action is continuous, since the kernels of
$\bar\rho$ and of $\bar\chi_1^{-1}\bar\eps^{\,e}$ are open. The place $w$ is not above $\ell$ (the places of $S_\iota$ are prime to
$\ell_\iota$), and $\gamma$ lies in the inertia group at $w$; hence $\bar\eps(\gamma)=1$ and $\rho'(\gamma)=\bar\chi_1(\gamma)^{-1}\bar\rho(\gamma)$.
As $\bar\chi_1(\gamma)$ is not a root of the characteristic polynomial of $\bar\rho(\gamma)$, the map
$\rho'(\gamma)-1=\bar\chi_1(\gamma)^{-1}(\bar\rho(\gamma)-\bar\chi_1(\gamma))$ is bijective, and \cref{pf:linkN:tamevanish} gives the
vanishing for the $\F_\ell$-structure. The continuous cochains, cocycles and coboundaries of $M$ are the same for the
$k_\iota$-structure and the $\F_\ell$-structure, so the vanishing holds for the $k_\iota$-representation.
\end{proof}

\begin{lemma}[Auxiliary-place clauses]\label{pf:linkN:g5}
Assume \cref{pf:linkN:Hdec}. Suppose that $D$ satisfies (GP2), that residual representations exist at the local indices with good
eigensystems, and that $\mathrm{TP}(D)$ holds. Then $D$ satisfies the \emph{auxiliary-place clauses} (AP): at every local index
$\iota$ with $\linkNgood_\iota\ne\emptyset$, for every good family $r$ and every place $w\in S_\iota\cup Q_\iota$ not above $\ell_\iota$,
\begin{enumerate}[label=(\roman*)]
\item if $w\in Q_\iota$ or $w$ lies above $p_{\mathrm{de}}$ (or above $p'_{\mathrm{de}}$ when $\ell_\iota=p_{\mathrm{de}}$): for every $x$ the
inertia group at $w$ acts through $r_x$ by matrices simultaneously conjugate to $\mathrm{diag}(1,\dots,1,\chi_{\iota,1},\dots,\chi_{\iota,1})$,
with $r$ entries equal to $1$;
\item if $w\in S_\iota$ lies above $g_0$ or $g_1$: there is a polynomial over $T_\iota$, with unit value at $\chi_{\iota,1}(\gamma_{\iota,w})$,
which evaluated at $r_x(\gamma_{\iota,w})$ annihilates $r_x(\gamma)-\chi_{\iota,1}(\gamma)$ for all $x$ and all $\gamma$ in the inertia group;
\item if $w\in S_\iota$: an idempotent-choice vanishing statement holds, and the restrictions of the twisted residual representations
$\bar\rho\otimes\bar\chi_1^{-1}\bar\eps^{\,e}$ ($e\in\Z$) to the decomposition group at $w$ have vanishing $H^0$, $H^1$ and $H^2$.
\end{enumerate}
\formallabel{UpperTheory.GalS.prvG5_of_admissible_cutRec}
\end{lemma}

\begin{proof}
By \lref{UpperTheory.GalS.lgc_of_cut'_alt}, (GP2) and $\mathrm{TP}(D)$ imply local--global compatibility at the auxiliary
places. By \lref{UpperTheory.GalS.eigPd_holds_alt}, (GP2), local--global compatibility and the existence of residual
representations imply that $\bar\chi_1(\gamma_{\iota,w})$ is not an eigenvalue of $\bar\rho(\gamma_{\iota,w})$ for $w\in S_\iota$ above
$p_{\mathrm{de}}$ (resp.\ $p'_{\mathrm{de}}$ when $\ell_\iota=p_{\mathrm{de}}$). By
\lref{Formalism.MainNotions.SumRepresentations.galois_clause_of_sufficient_alt}, (AP) follows from (GP2) and a
\emph{sufficient clause}: local--global compatibility at the auxiliary places, and, for $w\in S_\iota$, the existence of residual
representations and, for each of them, the non-eigenvalue condition at $\gamma_{\iota,w}$ and the vanishing of \cref{pf:linkN:h12}.
By definition, $w\in S_\iota$ means that $w$ lies above $g_0$, $g_1$ or $p_{\mathrm{de}}$ (resp.\ $p'_{\mathrm{de}}$) and not above
$\ell_\iota$. Above $g_k$ the non-eigenvalue condition is part of the residual irreducibility clause of (GP2), as $\ell_\iota\ne g_k$;
above $p_{\mathrm{de}}$ it was obtained above. \Cref{pf:linkN:h12} then gives the vanishing.
\end{proof}

\begin{theorem}[Facts from the preconditions]\label{pf:linkN:facts}
Assume \cref{pf:linkN:Hcheb,pf:linkN:Hdec}. If $\mathrm{GP}(D,\mathrm{nch})$ and $\mathrm{TP}(D)$ hold, then $\mathrm{NF}(D,\mathrm{nch})$
holds.
\formallabel{UpperTheory.GalC.facts_of_galois_preconditions_of_cutRec}
\end{theorem}

\begin{proof}
The conditions (GP1)--(GP3), (MS) and (GP5)--(GP18) are among the preconditions. Residual representations exist by a point
condition (\lref{GlobalTheory.T.residualRankPartExists_of_jointCondition}). The de Rham condition (DR) holds by
\lref{UpperTheory.GalS.de_rham_points_weight_jumps_alt}: (Adm) and $\mathrm{TP}(D)$ make the local parameters above the type
primes type-admissible, which determines the de Rham property and the jumps. The auxiliary-place clauses hold by
\cref{pf:linkN:g5}. The eigenvector clause (EV) follows from (GP2) and the existence of residual representations by
\lref{GlobalTheory.T.prvG6_of_admissible'}. Admissibility of the $J$-parameters follows from (Adm), $\mathrm{TP}(D)$ and the
weight and block character parameter conditions by \lref{UpperTheory.GalS.admJ_of_admPar}. The duties follow from (GP14),
(Perm), the distinct-jumps part of (DR) and residual irreducibility, using \cref{pf:linkN:Hcheb} for the design conditions and the
type-prime condition, by \lref{UpperTheory.GalS.duties_of_large_aux41}.
\end{proof}

\begin{proposition}[Existence of a good cut]\label{pf:linkN:existsgood}
If $\mathrm{GP}(D,\mathrm{nch})$ and $\sfname{Tie4}(D)$ hold, then $D$ has a good cut.
\formallabel{UpperTheory.GalC.exists_good_alt}
\end{proposition}

\begin{proof}
Let $Q$ be the pending datum of $\sfname{Tie4}(D)$, with (TA) and (Eig); the preconditions give (Adm) and (Perm), and $\sfname{Tie4}$
gives (Ser). Keep $\iota$ if and only if $\linkNgood_\iota\ne\emptyset$, and let $X^c_\iota$ be the set of $x\in X_\iota$ for which some
automorphic point with eigensystem $x$ has a type vector $F$ with $e_0\cdot F\ne0$.

(b) Let $x\in\linkNgood_\iota$. Choose an automorphic point $P$ with eigensystem $x$ and a non-zero type vector $F$ of $P$; by
\cref{pf:linkN:typevec}, $e_0\cdot F=F\ne0$, so $x\in X^c_\iota$. If $x\in Z_\iota\cap X^c_\iota$ and $x\notin\linkNgood_\iota$, then
\cref{pf:linkN:typevec} gives $e_0\cdot F=0$ for every type vector over $x$, contradicting $x\in X^c_\iota$. In particular $X^c_\iota\ne\emptyset$
for kept $\iota$, so $c$ is a cut, and (a) holds by definition.

(c) By \cref{pf:linkN:operator} there is $\alpha$ commuting with $T_\iota$, acting by $\#G$ on $M_K[x]$ for $x\in\linkNgood_\iota$ and by
$0$ on $M_K[x]$ whenever $e_0$ kills every type vector over $x$, which is the case for every $x\notin X^c_\iota$.

(d) Let $\iota$ be kept, $P$ an automorphic point with eigensystem $x\in X^c_\iota$, $k\in\{0,1\}$ and $w\mid g_k$. Choose an automorphic
point $P_1$ with eigensystem $x$ and a type vector $F$ of $P_1$ with $e_0\cdot F\ne0$. Since $e_0$ is a sum of products over good
families, there is a good family $s$ with $e_{\mathrm{Ser}}(k',v,s(k',v))\cdot F\ne0$ for all $(k',v)$
(\lref{UpperTheory.GalC.exists_tuple}). Writing $e_{\mathrm{Ser}}(k',s(k',v))$ as a sum of central idempotents of irreducible
representations with label $s(k',v)$, one of them acts non-trivially on $F$; it is a type of $P_1$ at $(k',v)$ with label $s(k',v)$
(\lref{UpperTheory.GalC.exists_type}). By (TA) and (Adm) the local parameter of $P_1$ at every place above $g_k$ is
type-admissible (\lref{UpperTheory.GalC.nc_of_type}). By (Adm), $P_1$ and $P$ carry base change parameters, which are
permutations of each other by (Perm), since both points have eigensystem $x$. The places above $g_k$ are split, and the local
parameters at split places are determined by the base change parameters up to such a permutation
(\lref{GlobalTheory.RowThetaParam.isBCV_of_perm}). Hence the local parameter of $P$ at $w$ equals that of $P_1$ and is
type-admissible.
\end{proof}

\begin{theorem}[Good cuts with facts]\label{pf:linkN:goodcutfacts}
Assume \cref{pf:linkN:Hcheb,pf:linkN:Hdec}. Then $\sfname{CutGivesAlt}$ holds.
\formallabel{UpperTheory.GalC.exists_good_cut_facts_tie4}
\end{theorem}

\begin{proof}
Let $\mathrm{GP}(D,\mathrm{nch})$ and $\sfname{Tie4}(D)$ hold, and let $c$ be the good cut of \cref{pf:linkN:existsgood}. The Galois
preconditions pass to $\mathrm{cut}(D,c)$ (\lref{UpperTheory.GalC.pre_cut_alt}): each of the eighteen conditions restricts to
the kept local indices and to the eigensystems in $X^c_\iota$ (for (MS) and (LU) this uses (b) of \cref{pf:linkN:goodcut}), and the
first provision (GP3) of the cut is re-derived from the restricted conditions by \lref{UpperTheory.GalS.Parts.main_alt_alt}.
The cut satisfies $\mathrm{TP}$ by \cref{pf:linkN:cutTP}, so \cref{pf:linkN:facts} gives the norm-choice facts of the cut. Finally $\sfname{Tie4}$ passes to the cut by restriction of each component to the kept indices, the pending datum being
restricted accordingly; (Eig) for the restricted datum follows from (Eig) for $D$ because the eigenspaces of the cut are images of
eigenspaces of $D$ (\cref{pf:linkN:qK}(i)), by \lref{UpperTheory.GalC.compatibilityTies_eigenRestriction}.
\end{proof}

% ---------------------------------------------------------------------------
\subsection{Proof of the Galois link}\label{sec:pf:linkN:b}

We now restate and prove \cref{thm:linkN}, with the notation of \cref{sec:families:rows} and of this section. Recall that
$\sfname{NL}(\mathcal C)=\sfname{Closed}(\mathcal C,\sfname{LinkN})$, and write
\[
p_{\mathrm{dep}}(h_L,c)=\sfname{parametersOfDependence}(h_L,c)
\]
for the parameter of \cref{sec:families:nonempty}.

\begin{theorem}[\cref{thm:linkN}]\label{pf:linkN:main}
\begin{enumerate}[label=(\alph*)]
\item Assume \cref{pf:linkN:H2,pf:linkN:Hdec,pf:linkN:Hcheb}. Let $\mathcal F\in\sfname{FamP}(\eps)$ be $\sfname{TwGood}$, and suppose
that $\mathrm{GP}(D_{x,i},\mathrm{nch}_{x,i})$ and $\sfname{Tie4}(D_{x,i})$ hold for all $x>x_1$ and $i\in I_x$. Then $\mathcal F$
satisfies $\sfname{LinkN}$.
\item Assume \cref{pf:linkN:H2,pf:linkN:Hdec,pf:linkN:Hlayer,pf:linkN:Hdz}. For every $h_L$ and $c\in\R$, the parameter
$p=p_{\mathrm{dep}}(h_L,c)$ satisfies
\[
\begin{aligned}
&\sfname{LinkedAt}\bigl(p,\ \sfname{Row}(\cdot,\cdot,\ \mathcal A\mapsto\sfname{Closed}(D3(3,\mathcal A),\sfname{BoundAtTwoThree}))\bigr)\\
&\qquad\Longrightarrow\ \sfname{LinkedAt}\bigl(p,\ \sfname{Row}(\cdot,\cdot,\ \mathcal A\mapsto\sfname{NL}(D3(3,\mathcal A)))\bigr).
\end{aligned}
\]
\end{enumerate}
\formallabel{UpperTheory.GalC.final_alt_alt}
\formallabel{GlobalTheory.NLink.Final.mainTheorem}
\end{theorem}

\paragraph{Part (a) with a cut hypothesis.}

\begin{theorem}\label{pf:linkN:famthm}
Assume $\sfname{Hypotheses}_2$ and $\sfname{CutGivesAlt}$. Let $\mathcal F\in\sfname{FamP}(\eps)$ be $\sfname{TwGood}$, and suppose that
$\mathrm{GP}(D_{x,i},\mathrm{nch}_{x,i})$ and $\sfname{Tie4}(D_{x,i})$ hold for all $x>x_1$ and $i\in I_x$. Then $\mathcal F$
satisfies $\sfname{LinkN}$, with constant $2000(|C_s|+1)^3+C_1$, where $C_s$ is the constant of the local bounds of the cut family
below and $C_1$ that of \eqref{eq:pf:linkN:logG} for $\mathcal F$.
\formallabel{UpperTheory.GalC.linkN_of_twGood_galoisPreconditions}
\end{theorem}

\begin{proof}
\emph{Step 1: constants.} Let $C=320$ and $C_T=40$ be the constants of \cref{pf:linkN:extra}.

\emph{Step 2: cuts.} For $x>x_1$ and $i\in I_x$ let $c_{x,i}$ be a good cut of $D_{x,i}$ such that
$\mathrm{NF}(\mathrm{cut}(D_{x,i},c_{x,i}),\mathrm{nch}_{x,i})$ and $\sfname{Tie4}(\mathrm{cut}(D_{x,i},c_{x,i}))$ hold; it exists by
$\sfname{CutGivesAlt}$. For $x\le x_1$ let $c_{x,i}$ be the trivial cut (all local indices kept, $X^c_\iota=X_\iota$). Let
$\mathcal F^c=\mathrm{gd}(\mathcal F,c)$ (\cref{pf:linkN:gd}).

\emph{Step 3: the cut family.} By Step 2 and \cref{pf:linkN:extra}, $\mathcal F^c$ has the facts with constants $(C,C_T)$. By
\cref{pf:linkN:twgd} it is $\sfname{TwGood}$, and by \cref{lem:families:at23} (which applies to every family under
$\sfname{Hypotheses}_2$) it satisfies the bound at $2$ and $3$.

\emph{Step 4: local bounds and summation.} By \lref{GlobalTheory.NLinkClose.main_alt_alt} (under $\sfname{Hypotheses}_2$, every
$\sfname{TwGood}$ family which has the facts with constants $(C,C_T)$, $C\ge0$, and satisfies the bound at $2$ and $3$ satisfies the
local bounds of \cref{pf:linkN:Concl} with some constant $C_s$), $\mathcal F^c$ satisfies the local bounds, and by
\cref{pf:linkN:sumthm} it satisfies $\sfname{LinkN}$ with constant $2000(|C_s|+1)^3$.

\emph{Step 5: return.} For $x>x_1$ the cut $c_{x,i}$ is good, $\mathrm{RE}(L_{x,i})$ is one of the point conditions (GP1), and
$\mathrm{MB}(L_{x,i})$ holds by \cref{pf:linkN:multle}, since (ID), (SC), (MS) and (LU) are among the preconditions. By
\cref{pf:linkN:return}, $\mathcal F$ satisfies $\sfname{LinkN}$ with constant $2000(|C_s|+1)^3+C_1$.
\end{proof}

\paragraph{Rows.}
We make the notions of \cref{def:families:rows} precise. A \emph{parameter tuple} is
$\kappa=(\eps,\mathbf C,c_D,c_h,C_5,c_6,C_b)$, with $\eps$ real, $\mathbf C$ a constant package and the other entries real; it is
\emph{positive} if $C_5,c_6,C_b>0$. A context $K$ is \emph{of} $\kappa$ if its constants are those of $\kappa$ and it is good for
$(C_5,c_6,C_b)$. For a property $p(\kappa,K)$ let $\mathrm{Ev}_{\mathrm{par}}(p)$ be the statement: there is $\eps_0>0$ such that for all
$\eps\in(0,\eps_0]$ with $\eps<1/52$, all $C_8\ge1$ and all $C_5,c_6,C_b$ there is a threshold package $K_0$ such that for all standard
constant packages $\mathbf C$ with parameter $C_8$ which are large for $K_0$ there is $c_1$ such that for all $c_D\ge c_1$ there is $c_2$
such that for all $c_h\ge c_2$, if $\kappa=(\eps,\mathbf C,c_D,c_h,C_5,c_6,C_b)$ is positive, then there is $x_0$ with $p(\kappa,K)$
for all contexts $K$ of $\kappa$ with scale $x\ge x_0$.

\begin{definition}[Rows and linkedness]\label{pf:linkN:rows}
Let $Q$ be a slot and $\mathrm{Lk}$ a linkage for the bundle type of \cref{sec:families:rows}.
\begin{enumerate}[label=(\alph*)]
\item For a property $E(\kappa,D,O,c,b)$ of bundles, $\mathrm{Ev}_{Q,\mathrm{Lk}}(E)$ is $\mathrm{Ev}_{\mathrm{par}}$ of the statement: for every
point $\beta$ over $K$ lying in $Q$, every $c\in\R$ and every bundle $b$ over $(\beta.D,\beta.O,c)$ with $\mathrm{Lk}(K,\beta,c,b)$,
$E(\kappa,\beta.D,\beta.O,c,b)$ holds.
\item A \emph{raw selection} $\sigma$ for a class $\mathcal A$ assigns to every $\mathcal F\in\mathcal A$, $x>x_1$ and $i\in I_x$ a bundle
$\sigma(\mathcal F,x,i)$ over $(D_{x,i},O_{x,i},C_{\mathrm{sz}}(\mathcal F))$, where the object $O_{x,i}$ of the member consists of its
size data, lattice data and extra record (so its norm choice is $\mathrm{nch}_{x,i}$). It is a \emph{selection for $E$} if for every
$\mathcal F\in\mathcal A$ there is $\kappa$ whose constants are those of $\mathcal F$ such that $\mathcal F$ satisfies the seven duties
and every $\sigma(\mathcal F,x,i)$ is built from $Q$ (it is a bundle linked to a point in $Q$ over a context of $\kappa$) and satisfies
$E(\kappa,\cdot)$.
\item $\sfname{Row}(Q,\mathrm{Lk},H)$ means: there is $E$ with $\mathrm{Ev}_{Q,\mathrm{Lk}}(E)$ such that $H(\mathcal A,\sigma)$ holds for every
class $\mathcal A$ and every selection $\sigma$ for $E$. For a property $R$ of bundles, $\mathrm{All}(R)(\mathcal A,\sigma)$ means that
$R$ holds at every $\sigma(\mathcal F,x,i)$.
\item For slots $Q,R$, $Q\preceq R$ means: eventually in the parameters (with the extra constants) and for large twisting data
(\cref{def:families:twist}), for every base $B$ and frame $f$ admitting an exact field datum and a pin, every core over a stage tuple
with base $B$ which lies in $Q$ also lies in $R$.
\item For a parameter $p$ and a property $r$ of slots and linkages, $\sfname{LinkedAt}(p,r)$ means: there is a slot $R$ with
$\sfname{IsBuiltE}(p)\preceq R$ and $r\bigl(\sfname{IsGoodE}(p)\wedge R,\ \sfname{BundleLinkage}(p)\bigr)$.
\end{enumerate}
\formal{GlobalTheory.Fc.EvP}
\formal{GlobalTheory.Fc.RowE}
\formal{GlobalTheory.Fc.Row}
\formal{GlobalTheory.Fc.LargeBaseCondition}
\formal{GlobalTheory.Fc.RL.LinkedAt}
\end{definition}

\begin{lemma}[Row calculus]\label{pf:linkN:rowcalc}
Let $Q,Q'$ be slots and $\mathrm{Lk}$ a linkage.
\begin{enumerate}[label=(\roman*)]
\item If $p(\kappa,K)$ implies $q(\kappa,K)$ for all contexts $K$ of $\kappa$, then $\mathrm{Ev}_{\mathrm{par}}(p)$ implies $\mathrm{Ev}_{\mathrm{par}}(q)$.
\item If $\mathrm{Ev}_{Q,\mathrm{Lk}}(R)$ holds for a property $R$ not depending on $\kappa$, then $\sfname{Row}(Q,\mathrm{Lk},\mathrm{All}(R))$.
\item $\sfname{Row}(Q,\mathrm{Lk},H)$ and $\sfname{Row}(Q,\mathrm{Lk},H')$ imply $\sfname{Row}(Q,\mathrm{Lk},H\wedge H')$.
\item If $\sfname{Row}(Q,\mathrm{Lk},H)$ and $H$ implies $H'$ for every class and every selection for the trivial property, then
$\sfname{Row}(Q,\mathrm{Lk},H')$.
\item If every core in $Q$ lies in $Q'$, then $\sfname{Row}(Q',\mathrm{Lk},H)$ implies $\sfname{Row}(Q,\mathrm{Lk},H)$.
\item If $Q\preceq R$ and $Q\preceq R'$, then $Q\preceq R\wedge R'$.
\end{enumerate}
\formal{GlobalTheory.Fc.EvP.mono}
\formal{GlobalTheory.Fc.RowE.row}
\formal{GlobalTheory.Fc.Row.imp}
\formallabel{GlobalTheory.Fc.Row.strengthen}
\formal{GlobalTheory.Fc.LargeBaseCondition.and}
\end{lemma}

\begin{proof}
(i) The same thresholds work. (ii) Take $E=R$; every bundle of a selection for $R$ satisfies $R$. (iii) is
\lref{GlobalTheory.Fc.Row.and}: one takes the conjunction of the two properties $E,E'$, and a selection for $E\wedge E'$ is a
selection for each of them. (iv) Keep $E$; a selection for $E$ is a selection for the trivial property. (v) Keep $E$: points in $Q$ are
points in $Q'$, and bundles built from $Q$ are built from $Q'$. (vi) Both eventual statements hold simultaneously, with the smaller
$\eps_0$, the later thresholds and the unions of the exceptional sets.
\end{proof}

\begin{proposition}[The slot statement]\label{pf:linkN:slot}
Assume $\sfname{Hypotheses}_2$ and $\sfname{CutGivesAlt}$. Let $Q$ be a slot and $\mathrm{Lk}$ a linkage such that
$\mathrm{Ev}_{Q,\mathrm{Lk}}\bigl(\mathrm{GP}(D,O.\mathrm{nch})\wedge\sfname{Tie4}(D)\bigr)$ and
$\sfname{Row}\bigl(Q,\mathrm{Lk},\mathcal A\mapsto\sfname{Closed}(D3(3,\mathcal A),\sfname{BoundAtTwoThree})\bigr)$. Then
$\sfname{Row}\bigl(Q,\mathrm{Lk},\mathcal A\mapsto\sfname{NL}(D3(3,\mathcal A))\bigr)$.
\formallabel{UpperTheory.GalC.linkN_of_boundAtTwoThree_slot}
\end{proposition}

\begin{proof}
By \cref{pf:linkN:rowcalc}(ii),(iii) we have $\sfname{Row}(Q,\mathrm{Lk},H)$ with
$H(\mathcal A,\sigma)=\mathrm{All}(\mathrm{GP}\wedge\sfname{Tie4})(\mathcal A,\sigma)\wedge\sfname{Closed}(D3(3,\mathcal A),\sfname{BoundAtTwoThree})$.
By \cref{pf:linkN:rowcalc}(iv) it suffices to show $\sfname{NL}(D3(3,\mathcal A))$ for every class $\mathcal A$ and raw selection
$\sigma$ with $H(\mathcal A,\sigma)$. Let $\eps_1>0$ be given by the closedness in the second component of $H$, and for
$\eps\in(0,\eps_1]$ let $K$ be the threshold package given there. Let $\mathcal F\in\sfname{FamP}(\eps)$ belong to $D3(3,\mathcal A)$
and be large for $(K,C_H)$. Then $\mathcal F\in\mathcal A$ and $\mathcal F$ is $\sfname{TwGood}$, and by the first component of $H$,
for every $x>x_1$ and $i\in I_x$ the selected bundle satisfies $\mathrm{GP}(D_{x,i},\mathrm{nch}_{x,i})\wedge\sfname{Tie4}(D_{x,i})$. By
\cref{pf:linkN:famthm}, $\mathcal F$ satisfies $\sfname{LinkN}$.
\end{proof}

\begin{proposition}[Linkedness]\label{pf:linkN:linkedAt}
Assume $\sfname{Hypotheses}_2$ and $\sfname{CutGivesAlt}$, and let $p$ be a parameter. If
$\sfname{LinkedAt}\bigl(p,(Q,\mathrm{Lk})\mapsto\mathrm{Ev}_{Q,\mathrm{Lk}}(\mathrm{GP}(D,O.\mathrm{nch})\wedge\sfname{Tie4}(D))\bigr)$ and
$p$ is linked for $\sfname{Row}(\cdot,\cdot,\mathcal A\mapsto\sfname{Closed}(D3(3,\mathcal A),\sfname{BoundAtTwoThree}))$, then $p$ is
linked for $\sfname{Row}(\cdot,\cdot,\mathcal A\mapsto\sfname{NL}(D3(3,\mathcal A)))$.
\formallabel{UpperTheory.GalC.linkedAt_linkN_of_boundAtTwoThree}
\end{proposition}

\begin{proof}
Let $R_1$ and $R_3$ be the slots given by the two hypotheses, and $R=R_1\wedge R_3$. By \cref{pf:linkN:rowcalc}(vi),
$\sfname{IsBuiltE}(p)\preceq R$. The slot $\sfname{IsGoodE}(p)\wedge R$ is contained in $\sfname{IsGoodE}(p)\wedge R_1$ and in
$\sfname{IsGoodE}(p)\wedge R_3$, so the first hypothesis restricts to it by \cref{pf:linkN:rowcalc}(i), and the second by
\cref{pf:linkN:rowcalc}(v). \Cref{pf:linkN:slot} for this slot and $\sfname{BundleLinkage}(p)$ concludes.
\end{proof}

\paragraph{Properties at members of layer points.}
The construction produces data at \emph{layer points}: a context $K$, a base $B$, a frame $f$, twisting data $T$, a column choice
$w$, a layer $Z$, a set of twisted places, rows, a norm choice $\mathrm{nch}$, norm data and a core datum; the datum at the point
is the datum $w.D(Z)$ of the layer. For a property $P$ of such data, $\mathrm{Mem}(h_L,P)$ means: eventually in the parameters (with
the extra constants) and for large twisting data, for every base and frame admitting an exact field datum and a pin, $P$ holds at
all layer-point data whose default extension (built from $h_L$) and reading position are non-empty and whose column pin holds.
If $\mathrm{Mem}(h_L,P)$ and $\mathrm{Mem}(h_L,P')$, then $\mathrm{Mem}(h_L,P\wedge P')$, by
\lref{GlobalTheory.HeightCounting.Up.Model.AtMem.and_z} (one intersects the two eventual statements).

\begin{proposition}[Transport to linked bundles]\label{pf:linkN:transport}
Let $p$ be the parameter attached to a norm-choice extension $N$, an $S$-extension and a geometric property, and assume that the
witnesses of $p$ move along set division and that every $N$-datum yields an extension-with-law datum. If $P(D,\mathrm{nch})$ is a
property of data and norm choices with $\mathrm{Mem}(h_L,P(w.D(Z),\mathrm{nch}))$, then
$\sfname{LinkedAt}\bigl(p,(Q,\mathrm{Lk})\mapsto\mathrm{Ev}_{Q,\mathrm{Lk}}(P(D,O.\mathrm{nch}))\bigr)$.
\formallabel{GlobalTheory.Mem.AtPtTransport.atMem_at_wit_alt}
\formallabel{GlobalTheory.NRows.AtPt.lg1_entry_of_atWit}
\formallabel{GlobalTheory.NRows.AtPt.memSlot_entry_of_atWit}
\formallabel{GlobalTheory.NRows.AtPt.exists_wPt_entC_atWit}
\formallabel{GlobalTheory.NRows.AtPt.exists_wit_pt_alt}
\end{proposition}

\begin{proof}
Let $R_P$ be the \emph{membership location} of $P$: the slot of cores over a stage tuple with base $B$ such that $P$ holds at all
layer-point data over $B$ as in the definition of $\mathrm{Mem}$. By hypothesis, eventually in the parameters (with the extra
constants) and for large twisting data, every core lies in $R_P$; so the slot $R$ of cores which lie in $R_P$ whenever their
context is one at which this inclusion holds satisfies $\sfname{IsBuiltE}(p)\preceq R$, and eventually every context is of this kind
(\lref{GlobalTheory.Fc.LgModel.largeBaseProperty_conditional}: such conditional slots always satisfy the large base condition).
It therefore suffices to prove $P(D,O.\mathrm{nch})$ for every point $\beta$ in $\sfname{IsGoodE}(p)\wedge R_P$ and every bundle linked
to it. A good point of $p$ carries a witnessed point of the slot data of $p$; moving its witness to the set-division lattices of its
layer (possible by hypothesis) produces a witnessed point $q$ with the same datum $D$ and object $O$ as $\beta$ whose layer-point data
have non-empty default extension (an extension-with-law datum gives a default-extension datum), non-empty reading position and
column pin. Since $q$ lies over a base in $R_P$, $P$ holds at its layer-point data, i.e.\ $P(D,O.\mathrm{nch})$.
\end{proof}

\begin{lemma}[Moving positions]\label{pf:linkN:moving}
For every $h_L$, $c\in\R$ and geometric property $\Gamma$, the slot data of the parameter attached to the norm-choice extension with
list pin of level $c$, the trivial $S$-extension and the geometric closure of $\Gamma$ have the property that every position can be
moved to the set-division lattices of its layer. Consequently the witnesses of $p_{\mathrm{dep}}(h_L,c)$ move along
set division.
\formallabel{GlobalTheory.NLink.Final.slot_mv_extension_with_list_geo_cl}
\end{lemma}

\begin{proof}
Given a position, the rows pass to the moved division lattices and heights
(\lref{GlobalTheory.T.rows_pass_to_moved_divisionLatticesHeights}); there is a kept norm choice with upper lattice data on the
moved layer (\lref{GlobalTheory.T.exists_kept_choice_upper_lattice_alt}); the list pin is preserved, since the rows are the same
(\lref{GlobalTheory.Mem.GridMove.listPin_rowsSame}); the spectral stage data transfer to the moved lattices
(\lref{GlobalTheory.T.spectralStageData_transfer_lattices_heights}); and the division geometry scalars transfer along the set
division (\lref{GlobalTheory.T.divisionGeometryScalars_transfer_setDivision}). This is the moving property. By
\lref{GlobalTheory.T.witnesses_pass_along_set_division}, moving positions implies moving witnesses; and
$p_{\mathrm{dep}}(h_L,c)$ is by definition the parameter of the lemma for the trivial geometric property.
\end{proof}

\paragraph{The preconditions at members.}
The following two theorems are proved by combining member properties established elsewhere; we describe their structure and cite
the ingredients.

\begin{theorem}[Galois preconditions at members]\label{pf:linkN:GPmem}
Assume \cref{pf:linkN:Hdec,pf:linkN:Hlayer,pf:linkN:Hdz}. For every $h_L$, $\mathrm{Mem}(h_L,\mathrm{GP}(w.D(Z),\mathrm{nch}))$.
\formallabel{UpperTheory.GalS.galois_preconditions_at_mem}
\formallabel{UpperTheory.GalS.core0_of_local_duality_inputs}
\end{theorem}

\begin{proof}
Three member properties are combined by the conjunction rule:
(i) the existence of \emph{linked inputs} for the datum, pinned at the character choice of the point, together with a count of
ramification jumps (\lref{GlobalTheory.NRows.AtLayer.exists_inputs_pinAt_jumpsCard_at_mem}, from $\mathcal P_3$ and
$\mathcal P_{\mathrm{Gal}}$); (ii) the condition on triples (GP16) (\lref{GlobalTheory.HeightCounting.Up.TripFin.finiteTriples_satisfy_condition_alt},
from $\mathcal P_3$ and $\mathcal P_{\mathrm{ser}}$); (iii) the goodness condition (GP18)
(\lref{GlobalTheory.T.goodFTAll_at_mem_of_inputs}). At a member carrying these, the conditions (GP1)--(GP12) are obtained as
follows: the point conditions from the theta inputs, Harish-Chandra finiteness and rapid decay
(\lref{GlobalTheory.T.theta_conditions_of_inputs_of_jumpsCard}); (Adm) and (ID) are components of the linked inputs, and the
first provision follows from them by \lref{UpperTheory.GalS.Parts.main_alt_alt}; admissibility of the residual tame eigenvalues
(\lref{GlobalTheory.AdmTame.main_alt_alt_alt}); residual irreducibility in rank $r$ from a three-prime theorem and the
Brauer--Nesbitt theorem \cite{CurtisReiner1962}
(\lref{Formalism.MainNotions.AdmissibleThreePrime.residual_irreducible_of_theorem_three_prime_alt}), which also gives the
residual irreducibility clause of (GP2); (MS) from a reduced multiplicity statement
(\lref{GlobalTheory.T.reducedMultiplicity_of_inputs}); the weight and block character parameter conditions from the character
choice (\lref{GlobalTheory.T.weight_param_of_character_choice} and \lref{GlobalTheory.BlockCharInf.main}); unramifiedness of $\mu$
from the finiteness of automorphic forms in $\mathcal P_3$ (\lref{GlobalTheory.MuUnrTie.main_alt_alt}); and the depth-zero
condition from \cref{pf:linkN:Hdz} through canonical row modules
(\lref{GlobalTheory.T.depthZeroG_of_rowMod} and \lref{GlobalTheory.RowModCanon.main_alt_alt}); (GP11) and (GP12) are components of the
linked inputs. Finally (GP13) is a localisation-coincidence theorem (\lref{GlobalTheory.LocCoincide.main_alt_alt}), (GP14) is
part of the default extension, (GP15) is part of the norm data, (GP16) and (GP18) are (ii) and (iii), and (GP17) holds by
\lref{GlobalTheory.HeightCounting.SameTw.sameTwist_at_alt}.
\end{proof}

\begin{theorem}[$\sfname{Tie4}$ at members]\label{pf:linkN:Tie4mem}
Assume $\mathcal P_3$, the height-counting statement contained in $\mathcal P_{\mathrm{Gal}}$, and \cref{pf:linkN:Hdz}. For every
$h_L$, $\mathrm{Mem}(h_L,\sfname{Tie4}(w.D(Z)))$.
\formallabel{GlobalTheory.T.tieChi_at_members_of_depthZero}
\formallabel{GlobalTheory.TypePrimesFinished.locComp_of_isRc_classMapD_cenLoc}
\end{theorem}

\begin{proof}
Three member properties are combined by the conjunction rule and \lref{GlobalTheory.HeightCounting.Up.Model.AtMem.at_member_imp}
(member properties are closed under implication):
(i) (Chi) together with a pending datum satisfying (TA) and (Eig), by \lref{GlobalTheory.T.tieChi_at_mem_of_depth_zero_inputs},
which uses canonical row modules, Green's labelling of the characters of finite general linear groups \cite{Green1955}, depth-zero
support, rationality of roots of unity and of series idempotents, and central-locality, $\Z_\ell$-stability and central integrality
properties at all members (\lref{GlobalTheory.T.cenLoc_zlStable_cenInt_at_members});
(ii) the existence of a pinned character choice (\lref{GlobalTheory.T.exists_character_choice_pin_at_of_third_layer}), which
gives (Ine) by \lref{GlobalTheory.T.tieInePoly_of_characterChoice};
(iii) (Ser), from $\mathcal P_3$ (\lref{GlobalTheory.T.serLoc_at_member_of_published_inputs}).
\end{proof}

\begin{proposition}\label{pf:linkN:finalOfAlt}
Assume $\sfname{Hypotheses}_2$, $\sfname{CutGivesAlt}$, $\mathcal P_3$, $\mathcal P_{\mathrm{Gal}}$ and \cref{pf:linkN:Hdz}. Let $h_L$ be
given and assume $\mathrm{Mem}(h_L,\mathrm{GP}(w.D(Z),\mathrm{nch}))$. Then for every $c\in\R$ the implication of
\cref{pf:linkN:main}(b) holds for the parameter $p=p_{\mathrm{dep}}(h_L,c)$.
\formallabel{UpperTheory.GalC.final_of_alt}
\end{proposition}

\begin{proof}
By \cref{pf:linkN:moving} the witnesses of $p$ move along set division. By \cref{pf:linkN:Tie4mem}, the hypothesis and the
conjunction rule,
\[
\mathrm{Mem}\bigl(h_L,\ \mathrm{GP}(w.D(Z),\mathrm{nch})\wedge\sfname{Tie4}(w.D(Z))\bigr).
\]
The norm-choice extension of $p$ consists of an extension-with-law datum and a list pin, so every datum of it yields an
extension-with-law datum. \Cref{pf:linkN:transport} gives
\[
\sfname{LinkedAt}\bigl(p,\ (Q,\mathrm{Lk})\mapsto\mathrm{Ev}_{Q,\mathrm{Lk}}(\mathrm{GP}(D,O.\mathrm{nch})\wedge\sfname{Tie4}(D))\bigr),
\]
and \cref{pf:linkN:linkedAt} concludes.
\end{proof}

\begin{proof}[Proof of \cref{pf:linkN:main}]
(a) Under \cref{pf:linkN:Hcheb,pf:linkN:Hdec}, $\sfname{CutGivesAlt}$ holds by \cref{pf:linkN:goodcutfacts}; apply
\cref{pf:linkN:famthm}.
(b) $\mathcal P_3$ contains \cref{pf:linkN:Hcheb}, so $\sfname{CutGivesAlt}$ holds by \cref{pf:linkN:goodcutfacts}, and
\cref{pf:linkN:GPmem} gives $\mathrm{Mem}(h_L,\mathrm{GP}(w.D(Z),\mathrm{nch}))$. Apply \cref{pf:linkN:finalOfAlt}.
\end{proof}

% ---------------------------------------------------------------------------
\subsection{Cocycle bounds}\label{sec:pf:linkN:cocycles}

This subsection proves the statements about global Galois cocycles used in the global congruence count of the proof of the
local bounds (\lref{GlobalTheory.NLinkClose.main_alt_alt}).

\begin{definition}[Notation]\label{pf:linkN:cocnot}
Let $D$ be a datum and $\iota$ a local index; write $\ell=\ell_\iota$, $\OO=\OO_\iota$, $K$ for its fraction field and $\varpi$ for a
uniformiser. For $k\ge1$ let $W_k=K^k/\OO^k$, a discrete $\OO$-module. A homomorphism of monoids $\mathsf M:G_E\to M_k(\OO)$ defines an
action $g\cdot x=\mathsf M(g)x$ of $G_E$ on $W_k$; $Z^1(\mathsf M)$ denotes the $\OO$-module of continuous maps $f:G_E\to W_k$ with
$f(gh)=f(g)+g\cdot f(h)$, and $B^1(\mathsf M)\subseteq Z^1(\mathsf M)$ the coboundaries $g\mapsto g\cdot x-x$. For an $\OO$-module $C$ of maps
$G_E\to W_k$ put $\len_{\mathrm{cob}}(C)=\len_\OO(C\cap B^1(\mathsf M))$. For $1\le c\le t$ let $\chi_c=\chi_{\iota,c}$; by the definition of a datum
$\chi_c=\chi_1\cdot\chi_{\mathrm{cyc}}^{-(c-1)}$, where $\chi_{\mathrm{cyc}}:G_E\to\Zl^\times\to\OO^\times$ is the $\ell$-adic cyclotomic character, and
we write $\mathsf M_c$ for the $1\times1$ matrix-valued character $\chi_c\chi_1^{-1}$ (so $\mathsf M_1$ is trivial).

The Galois data of $D$ fix, for every finite place $v$ of $E$, a place of $\Omega$ above $v$, with absolute value $|\cdot|_v$; $D_v$ and $I_v$ are
its decomposition and inertia groups in $G_E=\Gal(\Omega/E)$. For a module $C$ of maps $G_E\to W_k$ let $\mathrm{res}_v(C)=\{f|_{I_v}:f\in C\}$.
An element $\sigma\in G_E$ is \emph{wild at $v$} if $|\sigma(y)/y-1|_v<1$ for all $y\in\Omega^\times$. An element $\gamma$ is a \emph{tame generator}
at $v$ if $\gamma\in I_v$ and $I_v$ lies in the closure of the subgroup generated by $\gamma$ and the elements wild at $v$; the local index
provides such elements $\gamma_{\iota,v}$ for $v\in S_\iota\cup Q_\iota$.
\formal{Formalism.MainNotions.N.KOG}
\formal{GlobalTheory.Cut.resIne}
\formal{Formalism.MainNotions.N.wildSet}
\formal{Formalism.BasicNotions.TameGenAt}
\end{definition}

\begin{lemma}\label{pf:linkN:cycunr}
Let $v$ be a finite place of $E$ not above $\ell$ and $\gamma\in I_v$. Then $\chi_{\mathrm{cyc}}(\gamma)=1$ and $\chi_c(\gamma)=\chi_1(\gamma)$ for
$1\le c\le t$.
\formallabel{GlobalTheory.T.cyclotomicChar_unramified_away_from_ell}
\formallabel{GlobalTheory.CycUnram.main}
\end{lemma}

\begin{proof}
Let $n\ge1$ and $\zeta\in\Omega$ a primitive $\ell^n$-th root of unity, so $\gamma(\zeta)=\zeta^a$ with $a\equiv\chi_{\mathrm{cyc}}(\gamma)\bmod\ell^n$.
As $\gamma$ is in the inertia group, $|\gamma(\zeta)-\zeta|_v<1$; since $\ell$ is a unit at $v$, the $\ell^n$-th roots of unity have distinct
images in the residue field, so $\zeta^a=\zeta$ and $\chi_{\mathrm{cyc}}(\gamma)\equiv1\bmod\ell^n$. (Formally this is
\lref{ReciprocityTheory.Gal.cycZ_toZModPow_eq}, which states that $\chi_{\mathrm{cyc}}(\sigma)\equiv q\bmod\ell^n$ whenever
$|\sigma(x)-x^q|_v<1$ for all $x$ with $|x|_v\le1$, the place being prime to $\ell$.) Hence $\chi_{\mathrm{cyc}}(\gamma)=1$, and
$\chi_c(\gamma)=\chi_1(\gamma)\chi_{\mathrm{cyc}}(\gamma)^{-(c-1)}=\chi_1(\gamma)$.
\end{proof}

\begin{lemma}[Wild elements]\label{pf:linkN:wildel}
Let $\sigma$ be wild at $v$.
\begin{enumerate}[label=(\roman*)]
\item $\sigma^i$ is wild at $v$ for every $i\ge0$.
\item $\sigma\in I_v$.
\item If $k\in\N$ satisfies $|k|_v=1$ and $\sigma^k$ fixes $x\in\Omega$, then $\sigma$ fixes $x$.
\end{enumerate}
\formal{GlobalTheory.WildInertiaChar.wild_pow}
\formal{GlobalTheory.WildInertiaChar.fixed_of_pow_fixed}
\formal{GlobalTheory.TameCharUnique.mem_Ine_of_wild}
\end{lemma}

\begin{proof}
(i) Induction on $i$: with $a=\sigma(\sigma^iy)/\sigma^iy-1$ and $b=\sigma^iy/y-1$ we have $\sigma^{i+1}y/y-1=ab+a+b$, of absolute value $<1$.
(ii) For $y\ne0$, $|\sigma(y)|_v=|y|_v\,|1+(\sigma(y)/y-1)|_v=|y|_v$, so $\sigma$ preserves the valuation ring of the place, i.e.\
$\sigma\in D_v$; and for $0\ne x$ with $|x|_v\le1$, $|\sigma(x)-x|_v=|x|_v|\sigma(x)/x-1|_v<1$, so $\sigma$ acts trivially on the residue
field. (iii) Suppose $d=\sigma(x)-x\ne0$. Then $\sum_{i<k}\sigma^i(d)=\sigma^k(x)-x=0$, so $k=-\sum_{i<k}\bigl(\sigma^i(d)/d-1\bigr)$, and by (i)
every term has absolute value $<1$; hence $|k|_v<1$, a contradiction.
\end{proof}

\begin{lemma}[Wild inertia]\label{pf:linkN:wildinertia}
Let $v$ be a finite place of $E$ not above $\ell$, and let $f:G_E\to K/\OO$ be locally constant with $f(ab)=f(a)+f(b)$ for $a,b\in I_v$.
Then $f(\sigma)=0$ for every $\sigma$ wild at $v$.
\formallabel{GlobalTheory.T.wildInertia_killed_by_additive_locallyConstant}
\formallabel{GlobalTheory.WildAdditive.main}
\formallabel{GlobalTheory.WildAdditive.eq_zero_of_wild}
\formallabel{GlobalTheory.WildAdditive.exists_pow_nsmul_eq_zero}
\end{lemma}

\begin{proof}
First, every element of $K/\OO$ is killed by a power of $\ell$. Indeed $\ell$ is not a unit of $\OO$ (otherwise $\OO=\ell\OO$, and
$\OO=0$ by Nakayama's lemma over $\Zl$, as $\OO$ is finite over $\Zl$), so $\ell\in\varpi\OO$; an element $a/b$ ($b\ne0$) of $K$ satisfies
$b\OO=\varpi^j\OO$ for some $j$, and then $\ell^ja/b\in\OO$.

By \cref{pf:linkN:wildel}(ii), $\sigma\in I_v$, and additivity gives $f(1)=0$ and $f(\sigma^n)=nf(\sigma)$. Since $f$ is locally constant,
$f^{-1}(0)$ is an open neighbourhood of $1$ for the Krull topology, so it contains an open normal subgroup $H$ of finite index $n\ge1$; then
$\sigma^n\in H$. Write $n=\ell^ad$ with $\ell\nmid d$. By \cref{pf:linkN:wildel}(i), $\tau=\sigma^d$ is wild at $v$, and $\tau^{\ell^a}=\sigma^n$
fixes the fixed field of $H$. As $v\nmid\ell$, $|\ell^a|_v=1$, so \cref{pf:linkN:wildel}(iii) shows that $\tau$ fixes this field, i.e.\
$\tau\in H$. Hence $d\,f(\sigma)=f(\sigma^d)=0$. Choosing $m$ with $\ell^mf(\sigma)=0$ and using $\gcd(\ell^m,d)=1$ we get $f(\sigma)=0$.
\end{proof}

\begin{lemma}[Tame generators]\label{pf:linkN:tamegen}
Let $\gamma$ be a tame generator at $v$, $A$ an abelian group and $f:G_E\to A$ locally constant, with $f(ab)=f(a)+f(b)$ for $a,b\in I_v$,
$f(\sigma)=0$ for every $\sigma$ wild at $v$, and $f(\gamma)=0$. Then $f$ vanishes on $I_v$.
\formallabel{GlobalTheory.Cut.Quad.eq_zero_of_tameGen}
\end{lemma}

\begin{proof}
By additivity, $\{g\in I_v:f(g)=0\}$ is a subgroup; it contains $\gamma$ and, by \cref{pf:linkN:wildel}(ii), the elements wild at $v$. Hence the
subgroup generated by $\gamma$ and the wild elements lies in $f^{-1}(0)$, which is closed because $f$ is locally constant; so its closure,
which contains $I_v$, lies in $f^{-1}(0)$.
\end{proof}

\begin{lemma}\label{pf:linkN:torsKO}
Let $\OO$ be an integral domain with fraction field $K$ and $a\in\OO$, and assume that $a\ne0$ or that $\OO$ is a discrete valuation ring. Then
the $a$-torsion submodule $(K/\OO)[a]$ satisfies $\len(K/\OO)[a]\le\len(\OO/a\OO)$.
\formallabel{GlobalTheory.Cut.Quad.length_torsionBy_le}
\formallabel{GlobalTheory.UnramHom.length_torsionBy_le}
\end{lemma}

\begin{proof}
If $a=0$ and $\OO$ is a discrete valuation ring, $\OO/a\OO=\OO$ has infinite length. If $a\ne0$, the map $\OO\to(K/\OO)[a]$,
$c\mapsto ca^{-1}\bmod\OO$, kills $a\OO$ and is surjective: if $ay\in\OO$ then $y=(ay)a^{-1}$.
\end{proof}

\begin{definition}[Cocycle modules of an idempotent choice datum]\label{pf:linkN:cocW}
Let $r=(r_x)_{x\in X_\iota}$ be homomorphisms $G_E\to\GL_{n_V}(\OO)$, and $\tilde r_x=r_x\chi_1^{-1}$. An idempotent choice datum $\mathcal S$ at $\iota$
provides, among other data, a finite set $X'$ with a map $p:X'\to X_\iota$, an $\OO$-subalgebra $T'\subseteq\OO^{X'}$ and matrices
$\mathsf E(x')\in M_{n_V}(\OO)$ ($x'\in X'$). Let $R'\subseteq M_{n_V}(\OO)^{X'}$ be the set of finite sums of functions
$x'\mapsto f(x')\,r_{p(x')}(g)$ with $f\in T'$ and $g\in G_E$. For $\gamma\in G_E$ let $y_\gamma\in M_{n_V}(\OO)^{X'}$,
$y_\gamma(x')=\tilde r_{p(x')}(\gamma)\mathsf E(x')-\mathsf E(x')$, and let $Y$ be the set of finite sums $\sum_ja_jy_{\gamma_j}$ with $a_j\in R'$
(products taken pointwise). For $X_s\subseteq X_\iota$, $k\ge1$ and a map $\psi:M_{n_V}(\OO)^{X'}\to M_k(\OO)$, let $\mathrm{Coc}_{X_s}(\psi)$ be the
$\OO$-module of maps $G_E\to W_k$ spanned by the functions $\gamma\mapsto\phi(y_\gamma)$, where $\phi:M_{n_V}(\OO)^{X'}\to W_k$ runs over the
$\OO$-linear maps such that
\begin{enumerate}[label=(\alph*)]
\item $\phi(ay)=\psi(a)\phi(y)$ for $a\in R'$ and $y\in Y$, and
\item $\phi(y)=0$ for every $y\in Y$ which vanishes at all $x'$ with $p(x')\in X_s$.
\end{enumerate}
\formal{GlobalTheory.Cut.cocW}
\end{definition}

\begin{lemma}[Quadratic relation]\label{pf:linkN:cocquad}
Let $v$ be a finite place such that $(r_x(\gamma)-1)(r_x(\gamma)-\chi_1(\gamma))=0$ for all $x\in X_s$ and $\gamma\in I_v$. Then for every
$f\in\mathrm{Coc}_{X_s}(\psi)$ and $\gamma\in I_v$, $\psi(r(\gamma))f(\gamma)=f(\gamma)$, where $r(\gamma)=(r_{p(x')}(\gamma))_{x'}\in R'$.
\formallabel{GlobalTheory.Cut.Check.cocW_quad}
\end{lemma}

\begin{proof}
Both sides are $\OO$-linear in $f$, so it suffices to treat $f=\phi(y_\cdot)$ with $\phi$ as in the definition. The functions $r(\gamma)$ and
$-1=(-1)\cdot r_{p(\cdot)}(1)$ lie in $R'$, so $z=r(\gamma)y_\gamma-y_\gamma\in Y$. For $x'$ with $x=p(x')\in X_s$,
\[
z(x')=(r_x(\gamma)-1)(\chi_1(\gamma)^{-1}r_x(\gamma)-1)\mathsf E(x')=\chi_1(\gamma)^{-1}(r_x(\gamma)-1)(r_x(\gamma)-\chi_1(\gamma))\mathsf E(x')=0 .
\]
By (b), $\phi(z)=0$, and by (a), $\phi(r(\gamma)y_\gamma)=\psi(r(\gamma))\phi(y_\gamma)$; hence $\psi(r(\gamma))f(\gamma)=f(\gamma)$.
\end{proof}

\begin{lemma}[Inertia away from $\ell$, $S_\iota$, $Q_\iota$]\label{pf:linkN:lgsel}
Assume that $1\in R'$ and that $\tilde r_x(\gamma)=1$ for all $x\in X_s$, $\gamma\in I_v$ and all finite places $v$ not above $\ell$ and not in
$S_\iota\cup Q_\iota$. Then $\mathrm{res}_v\mathrm{Coc}_{X_s}(\psi)=0$ for all such $v$.
\formallabel{GlobalTheory.Proof.LocalGlobalSelmer.main}
\end{lemma}

\begin{proof}
(The assumption $1\in R'$ always holds: $1=1\cdot r_{p(\cdot)}(1)$.) Let $\gamma\in I_v$. Then $y_\gamma=1\cdot y_\gamma\in Y$, and
$y_\gamma(x')=(\tilde r_{p(x')}(\gamma)-1)\mathsf E(x')=0$ whenever $p(x')\in X_s$; so $\phi(y_\gamma)=0$ by (b). Thus every generator of
$\mathrm{Coc}_{X_s}(\psi)$ vanishes on $I_v$.
\end{proof}

\begin{lemma}[Peeling]\label{pf:linkN:peel}
Let $C$ be an $\OO$-module of maps $G_E\to W_k$ and $Q$ a finite set of finite places. There is $C'\subseteq C$ with $\mathrm{res}_v(C')=0$ for
$v\in Q$ and $\len C\le\len C'+\sum_{v\in Q}\len\mathrm{res}_v(C)$.
\formallabel{GlobalTheory.Cut.Peel.main}
\formallabel{GlobalTheory.Cut.Peel.exists_le_map_eq_bot}
\end{lemma}

\begin{proof}
More generally, for linear maps $f_q$ ($q\in Q$) on a module containing $C$, $C'=C\cap\bigcap_q\ker f_q$ satisfies
$\len C\le\len C'+\sum_q\len f_q(C)$. By induction on $Q$: if $C_0$ works for $Q_0$ and $q\notin Q_0$, then
$\len C_0=\len(C_0\cap\ker f_q)+\len f_q(C_0)$ and $\len f_q(C_0)\le\len f_q(C)$. Apply this to the restriction maps to $I_v$.
\end{proof}

\begin{definition}[Places above $\ell$]\label{pf:linkN:bridgedef}
For every place $w$ of $E$ above $\ell$ the local theory of $D$ provides a \emph{local bridge}: a profinite group $G_w$ with an isomorphism
$\iota_w$ onto $D_w$, under which we let $I^{\mathrm{loc}}_w=\iota_w^{-1}(I_w)$; a complete discrete valuation ring $\OO_w$ with fraction field $L_w$,
algebraic over $\Ql$; a ring isomorphism $e_w:\OO\to\OO_w$ and a field homomorphism $K\to L_w$ extending it. It induces an $e_w$-semilinear map
$\phi_w:W_k\to L_w^k/\OO_w^k$. For $\mathsf M:G_E\to M_k(\OO)$, $G_w$ acts on $L_w^k/\OO_w^k$ through $e_w\circ\mathsf M\circ\iota_w$, and the
\emph{localisation} of $f:G_E\to W_k$ is $\mathrm{loc}_w(f)=\phi_w\circ f\circ\iota_w$. The local data provide $\OO_w$-modules
$B_w(\mathsf M)\subseteq F_w(\mathsf M)$ of maps $G_w\to L_w^k/\OO_w^k$: the coboundaries, and the \emph{finite part}, the sum of the coboundaries and
of the images of the Bloch--Kato finite cocycles of the period data of $D$ at $w$; and, for $r$ and $X_s$ as above, a set
$\mathrm{Cut}_w(\mathsf M)$ of \emph{cut cocycles}. For $e\in\N$, $\mathrm{LB}_w(\mathsf M,e)$ is the statement
$\len_{\OO_w}\bigl(\langle\mathrm{Cut}_w(\mathsf M)\rangle/(\langle\mathrm{Cut}_w(\mathsf M)\rangle\cap F_w(\mathsf M))\bigr)\le e\,\len(\OO_w/\ell\OO_w)$,
where $\langle\cdot\rangle$ is the $\OO_w$-span. A map $c:G_w\to L_w^k/\OO_w^k$ is an \emph{unramified cut cocycle} if there is $x$ with
$c(g)=g\cdot x-x$ for $g\in I^{\mathrm{loc}}_w$.
\formal{Formalism.BasicNotions.LocBridge}
\formal{GlobalTheory.Cut.locFun}
\formal{Formalism.WeaklyAdmissibleFiltrations.HfD}
\formal{GlobalTheory.Pass.LocBound}
\end{definition}

\begin{lemma}[Local bridges]\label{pf:linkN:bridge}
Let $w\mid\ell$.
\begin{enumerate}[label=(\roman*)]
\item The homomorphism $K\to L_w$ is bijective, so $\phi_w$ is bijective and $\phi_w(g\cdot x)=\iota_w^{-1}(g)\cdot\phi_w(x)$ for
$g\in D_w$.
\item If $c:D_w\to W_k$ satisfies $\phi_w\circ c\circ\iota_w\in B_w(\mathsf M)$, then there is $m\in W_k$ with $c(g)=g\cdot m-m$ for all
$g\in D_w$.
\end{enumerate}
\formallabel{GlobalTheory.CutTrivIne.eq_zero_of_phi}
\formallabel{GlobalTheory.Cut.ShaBound.toLLK_surj}
\formallabel{GlobalTheory.Cut.ShaBound.cob_of_bridge}
\end{lemma}

\begin{proof}
(i) Injectivity holds for a field homomorphism. For $y\in L_w$ there is $0\ne d\in\Zl$ with $dy\in\OO_w$, as $L_w$ is algebraic over $\Ql$;
writing $dy=e_w(n)$ with $n\in\OO$, the element $n/d\in K$ maps to $y$. An element of $K$ whose image lies in $\OO_w=e_w(\OO)$ therefore lies in
$\OO$, so $\phi_w$ is injective, and it is surjective by the first part. The compatibility with the actions holds because the action on
$L_w^k/\OO_w^k$ is through $e_w\circ\mathsf M\circ\iota_w$.
(ii) If $\phi_w(c(\iota_w(h)))=h\cdot x-x$ for all $h\in G_w$, put $m=\phi_w^{-1}(x)$ and use (i).
\end{proof}

\begin{lemma}[Unramified cut cocycles]\label{pf:linkN:cuttrivine}
Let $w\mid\ell$ and let $C$ be an $\OO$-module of maps $G_E\to W_1$ such that, for every $f\in C$, $\mathrm{loc}_w(f)$ is an unramified cut cocycle
for the trivial character $\mathsf M_1$. Then $\mathrm{res}_w(C)=0$.
\formallabel{GlobalTheory.CutTrivIne.main}
\end{lemma}

\begin{proof}
Let $f\in C$ and $\gamma\in I_w$, and let $g=\iota_w^{-1}(\gamma)\in I^{\mathrm{loc}}_w$. As $\mathsf M_1$ is trivial, $g$ acts trivially, so
$\phi_w(f(\gamma))=\mathrm{loc}_w(f)(g)=g\cdot x-x=0$, and $f(\gamma)=0$ by \cref{pf:linkN:bridge}(i).
\end{proof}

\begin{lemma}[Local conditions above $\ell$]\label{pf:linkN:countL}
Let $\mathsf M:G_E\to M_k(\OO)$, $e\in\N$, and let $C$ be an $\OO$-module of maps $G_E\to W_k$. Assume that for every place $w\mid\ell$,
$\mathrm{LB}_w(\mathsf M,e)$ holds and $\mathrm{loc}_w(f)\in\langle\mathrm{Cut}_w(\mathsf M)\rangle$ for all $f\in C$. Then there is $C'\subseteq C$ with
$\mathrm{loc}_w(f)\in F_w(\mathsf M)$ for all $f\in C'$ and $w\mid\ell$, and
\[
\len C\le\#\{w\mid\ell\}\cdot e\cdot\len(\OO/\ell\OO)+\len C'.
\]
\formallabel{GlobalTheory.T.length_global_cocycles_le_sum_places_over}
\formallabel{GlobalTheory.Cut.CountL.main}
\formallabel{GlobalTheory.Cut.CountL.length_le_inf_comap_add}
\formallabel{GlobalTheory.Cut.CountL.exists_le_forall_le}
\end{lemma}

\begin{proof}
For $w\mid\ell$ let $\kappa_w=\{f:\mathrm{loc}_w(f)\in F_w(\mathsf M)\}$ and $A_w=\langle\mathrm{Cut}_w(\mathsf M)\rangle$. We claim that every
submodule $C_0\subseteq C$ satisfies $\len C_0\le\len(C_0\cap\kappa_w)+e\len(\OO/\ell\OO)$. The map $C_0\to A_w/(A_w\cap F_w(\mathsf M))$,
$f\mapsto\mathrm{loc}_w(f)$, is $e_w$-semilinear with kernel $C_0\cap\kappa_w$. Since $e_w$ is an isomorphism, an injective $e_w$-semilinear map
carries strictly increasing chains of $\OO$-submodules to strictly increasing chains of $\OO_w$-submodules, so the length of $C_0/(C_0\cap\kappa_w)$
is at most $\len_{\OO_w}(A_w/(A_w\cap F_w(\mathsf M)))\le e\len(\OO_w/\ell\OO_w)=e\len(\OO/\ell\OO)$, by $\mathrm{LB}_w$ and the isomorphism $e_w$; the
claim follows from additivity of length. The set of places above $\ell$ is finite; applying the claim successively to
$C$, $C\cap\kappa_{w_1}$, $C\cap\kappa_{w_1}\cap\kappa_{w_2}$, \dots\ gives the lemma with $C'=C\cap\bigcap_w\kappa_w$.
\end{proof}

\begin{definition}[Selmer groups]\label{pf:linkN:seldef}
Let $\rho=(\rho_x)_{x\in X_\iota}$ be homomorphisms $G_E\to\GL_r(\OO)$, $\pi\in X_\iota$, and let $\mathsf M_{\rho,\pi}=\rho_\pi\chi_1^{-1}$, acting on
$W_r$. The Selmer group $\mathrm{Sel}(\iota,\rho,\pi)\subseteq H^1(G_E,W_r)$ consists of the classes whose localisation at every $w\mid\ell$ lies in
the image of $F_w(\mathsf M_{\rho,\pi})$, whose restriction to $I_v$ is trivial for every finite $v$ not above $\ell$ and not in
$S_\iota\cup Q_\iota$, with no condition at $S_\iota\cup Q_\iota$.
\formal{Formalism.MainNotions.N.Br.selPi}
\end{definition}

\begin{lemma}[Selmer bound]\label{pf:linkN:selclass}
Let $C\subseteq Z^1(\mathsf M_{\rho,\pi})$ be a submodule such that $\mathrm{loc}_w(f)\in F_w(\mathsf M_{\rho,\pi})$ for all $f\in C$ and $w\mid\ell$, and
$\mathrm{res}_v(C)=0$ for all finite $v$ not above $\ell$ and not in $S_\iota\cup Q_\iota$. Then
$\len C\le\len_{\mathrm{cob}}(C)+\len\mathrm{Sel}(\iota,\rho,\pi)$.
\formallabel{GlobalTheory.T.length_cocycles_le_coboundaries_add_selmer}
\formallabel{GlobalTheory.Cut.SelClass.main}
\formallabel{GlobalTheory.Cut.SelClass.length_le_cob_add}
\formallabel{GlobalTheory.Cut.SelClass.mk_mem_selG}
\end{lemma}

\begin{proof}
The class of every $f\in C$ lies in the Selmer group: the conditions above $\ell$ hold by assumption, and for $v$ not above $\ell$ and not in
$S_\iota\cup Q_\iota$ the restriction of $f$ to $I_v$ is $0$. The map $C\to\mathrm{Sel}(\iota,\rho,\pi)$, $f\mapsto[f]$, has kernel
$C\cap B^1(\mathsf M_{\rho,\pi})$, and additivity of length gives the claim.
\end{proof}

\begin{definition}[Tate--Shafarevich groups]\label{pf:linkN:shadef}
The Galois data of $D$ at $\ell$ provide the quotient $G_{E,S_\ell}$ of $G_E$ by the closed normal subgroup $N_{S_\ell}$ generated by the inertia
groups of all places of $\Omega$ above finite places of $E$ not above $\ell$, places of $\Omega$ above the places of $S_\ell$ (possibly different
from those fixed above), and the module $\Ql/\Zl(-e)$, on which $G_{E,S_\ell}$ acts through $\chi_{\mathrm{cyc}}^{-e}$; $\Sha^1(G_{E,S_\ell},\Ql/\Zl(-e))$
is the kernel of the restriction maps to the decomposition groups of these places, and we write $\mathrm{sha}(\ell,e)$ for its $\Zl$-length (in
$\N\cup\{\infty\}$).
\formal{Formalism.GaloisCohomology.Sha}
\formal{Formalism.GaloisCohomology.shaLen}
\end{definition}

\begin{lemma}[Cyclotomic blocks]\label{pf:linkN:shabound}
Let $2\le c\le t$ and let $C\subseteq Z^1(\mathsf M_c)$ be a submodule such that $\mathrm{loc}_w(f)\in B_w(\mathsf M_c)$ for all $f\in C$ and
$w\mid\ell$, and $\mathrm{res}_v(C)=0$ for every finite $v$ not above $\ell$. Then
\[
\len C\le\len_{\mathrm{cob}}(C)+\len(\OO/\ell\OO)\cdot\mathrm{sha}(\ell,c-1).
\]
\formallabel{GlobalTheory.T.length_cyclotomic_block_cocycles_le_sha}
\formallabel{GlobalTheory.Cut.ShaBound.main}
\formallabel{GlobalTheory.Cut.ShaBound.length_le}
\formallabel{GlobalTheory.Cut.ShaBound.length_quot}
\formallabel{GlobalTheory.Cut.ShaBound.exists_coord}
\formallabel{GlobalTheory.Cut.ShaBound.exists_toSha}
\formallabel{GlobalTheory.Cut.ShaBound.mk_mem_Sha}
\end{lemma}

\begin{proof}
\emph{Coordinates.} Let $d=[\OO:\Zl]$ and fix a $\Zl$-basis of $\OO$; it is a $\Ql$-basis of $K$, and the coordinates give a $\Zl$-linear
bijection $p=(p_1,\dots,p_d):K/\OO\to(\Ql/\Zl)^d$. By the formula for $\chi_c$, $\mathsf M_c(g)=\chi_{\mathrm{cyc}}(g)^{-(c-1)}\in\Zl^\times$ acts on $K/\OO$
as a $\Zl$-scalar, so $p_j(g\cdot x)=\chi_{\mathrm{cyc}}(g)^{-(c-1)}p_j(x)$. Hence for $f\in C$ each $a_j(f)=p_j\circ f$ is a continuous cocycle of $G_E$
with values in $\Ql/\Zl(-(c-1))$, and $a_j$ is $\Zl$-linear.

\emph{Local triviality.} Let $f\in C$ and $a=a_j(f)$. (1) $a$ vanishes on $I_v$ for $v\nmid\ell$. (2) For $w\mid\ell$, \cref{pf:linkN:bridge}(ii)
gives $m$ with $f(g)=g\cdot m-m$ on $D_w$, so $a(g)=g\cdot p_j(m)-p_j(m)$ on $D_w$. (3) $a$ vanishes on $N_{S_\ell}$: the set
$\{n\in N_{S_\ell}:a(n)=0\}$ is a subgroup (as $N_{S_\ell}$ acts trivially on $\Ql/\Zl(-(c-1))$, $a$ is additive on it), it is closed by
continuity of $a$, and it is normal, since the cocycle relation gives $a(gng^{-1})=g\cdot a(n)$ for $n\in N_{S_\ell}$ (as $n$ acts trivially
and $g\cdot a(g^{-1})=-a(g)$); it contains the inertia groups $I_v$ ($v\nmid\ell$) by (1), hence, being normal, all their conjugates, which are
the inertia groups of the other places of $\Omega$ above $v$. So it equals $N_{S_\ell}$. Consequently $a$ factors through a cocycle $\bar a$ of
$G_{E,S_\ell}$ (if $h^{-1}h'\in N_{S_\ell}$, then $a(h')=a(h)+h\cdot a(h^{-1}h')=a(h)$), and inflation is injective on cocycles. The decomposition groups of the places of $\Omega$ above
$w\in S_\ell$ used for $\Sha^1$ are conjugate to $D_w$ \cite[Ch.~I, Prop.~9.1]{Neukirch1999}, and the restriction of a cocycle to a conjugate
subgroup $gD_wg^{-1}$ is a coboundary if its restriction to $D_w$ is one (if $a(h)=h\cdot m-m$ on $D_w$, then
$a(ghg^{-1})=ghg^{-1}\cdot m'-m'$ with $m'=g\cdot m-a(g)$). Hence the class of $\bar a$ lies in $\Sha^1(G_{E,S_\ell},\Ql/\Zl(-(c-1)))$, and we
obtain $\Zl$-linear maps $\Psi_j:C\to\Sha^1$, $f\mapsto[\overline{a_j(f)}]$, such that $\Psi_j(f)=0$ implies that $a_j(f)$ is a coboundary.

\emph{Counting.} Let $\Psi=(\Psi_j)_j:C\to(\Sha^1)^d$. If $\Psi(f)=0$, choose $m_j$ with $a_j(f)(g)=g\cdot m_j-m_j$ and let $y=p^{-1}(m_1,\dots,m_d)$;
since $p$ is bijective and compatible with the actions, $f(g)=g\cdot y-y$, so $f\in B^1$. Put $B=C\cap B^1(\mathsf M_c)$. Then $C/B$ is a quotient of
$C/\ker\Psi$, which embeds into $(\Sha^1)^d$, so $\len_{\Zl}(C/B)\le d\cdot\mathrm{sha}(\ell,c-1)$. For an $\OO$-module $N$ of finite length,
$\len_{\Zl}N=f_\OO\len_\OO N$, where $f_\OO$ is the residue degree, and $d=e_\OO f_\OO$ with $e_\OO=\len_\OO(\OO/\ell\OO)$ (indeed
$\len_{\Zl}(\OO/\ell\OO)=d$, as $\OO/\ell\OO\cong\F_\ell^d$). Dividing by $f_\OO$ (the bound being trivial if $C/B$ has infinite length) gives
$\len_\OO(C/B)\le e_\OO\,\mathrm{sha}(\ell,c-1)$, and $\len C=\len B+\len(C/B)$.
\end{proof}

\begin{lemma}[Unramified homomorphisms]\label{pf:linkN:unramhom}
Assume global Artin reciprocity (as in \cref{pf:linkN:H2}). Let $C\subseteq Z^1(\mathsf M_1)$ be a submodule with $\mathrm{res}_v(C)=0$ for every finite
place $v$ of $E$. Then $\len C\le\len(\OO/\ell\OO)\cdot v_\ell(h_E)$.
\formallabel{GlobalTheory.UnramHom.length_le_classNumber}
\formallabel{GlobalTheory.UnramHom.length_le}
\formallabel{GlobalTheory.UnramHom.eq_one_of_class}
\formallabel{GlobalTheory.UnramHom.index_dvd_classNumber}
\end{lemma}

\begin{proof}
Since $\mathsf M_1$ is trivial, the elements of $C$ are continuous homomorphisms $G_E\to K/\OO$. Global Artin reciprocity, in the form of
\cref{pf:linkN:H2}, provides a continuous open surjection $\mathrm{art}$ from the id\`ele group $\mathbb I_E$ onto $G_E^{\mathrm{ab}}$, trivial on $E^\times$,
mapping the units $\OO_{E_v}^\times$ (placed at $v$) onto the image of $I_v$, for every finite place $v$. Let $\mathbb I^0\subseteq\mathbb I_E$ be the
kernel of the map to the ideal class group; its index divides $h_E$, and $\mathbb I^0\subseteq U\cdot E^\times$, where $U$ is the group of unit
id\`eles \cite[Ch.~VI, \S1]{Neukirch1999}. Let $H=\mathrm{art}(\mathbb I^0)\subseteq G_E^{\mathrm{ab}}$; as $\mathrm{art}$ is surjective, its index $n$
divides $h_E$.

\emph{Every $f\in C$ vanishes on the preimage of $H$.} Let $\psi=f\circ\mathrm{art}$, a continuous homomorphism to the discrete group $K/\OO$ with finite
image (the image of the compact group $G_E$). It kills $\OO_{E_v}^\times$ for every finite $v$, since $f$ vanishes on $I_v$. It takes the same
value at an id\`ele and at its finite part, by \lref{GlobalTheory.RingCls.map_finPartI} (for a totally complex field the archimedean
components are $N$-th powers for every $N$, and the image of $\psi$ is finite). A unit id\`ele with trivial archimedean components is the limit
of its partial products over finite sets of places (\lref{GlobalTheory.Mem.LocTriv.tendsto_part}), each killed by $\psi$, and $\ker\psi$ is
open, hence closed; so $\psi$ kills $U$. It kills $E^\times$. Hence $\psi$ kills $\mathbb I^0$.

\emph{Counting.} We show, by induction on $n$, that if $C$ is a module of homomorphisms $G_E\to K/\OO$ vanishing on the preimage of a subgroup
$H$ of $G_E^{\mathrm{ab}}$ of finite index $n$, then $\len C\le\len(\OO/\ell\OO)\,v_\ell(n)$. If $n=1$ then $C=0$. Otherwise choose a prime
$p\mid n$ and $a\in G_E$ whose image in $G_E^{\mathrm{ab}}/H$ has order $p$, and let $H'\supseteq H$ be the preimage of the subgroup generated by
the image of $a$, of index $n/p$. Evaluation $\mathrm{ev}:C\to K/\OO$, $f\mapsto f(a)$, has kernel vanishing on the preimage of $H'$ (if $b\equiv a^j$
modulo $H$, then $f(b)=jf(a)+f(a^{-j}b)=0$), so $\len\ker\mathrm{ev}\le\len(\OO/\ell\OO)v_\ell(n/p)$ by induction. Its image is killed by $p$,
since $a^p$ maps into $H$. If $p\ne\ell$, $p$ is a unit in $\OO$ and the image is $0$; if $p=\ell$, the image has length at most
$\len(\OO/\ell\OO)$ by \cref{pf:linkN:torsKO}. In both cases $\len C\le\len(\OO/\ell\OO)v_\ell(n)$. Finally $v_\ell(n)\le v_\ell(h_E)$ as $n\mid h_E$.
\end{proof}

% ---------------------------------------------------------------------------
\subsection{Remarks on the formal statements}\label{sec:pf:linkN:remarks}

\begin{remark}[Use of the hypothesis of part (b)]\label{pf:linkN:remhyp}
In the proof of \cref{pf:linkN:main}(b) the hypothesis that $p$ is linked for the bound at $2$ and $3$ is used through its auxiliary
slot $R_3$ (\cref{pf:linkN:linkedAt}), through its row property, which becomes part of the selection property of the conclusion
(\cref{pf:linkN:rowcalc}(iii)), and through the number $\eps_1$ and the threshold packages $K$ of its closedness statement
(\cref{pf:linkN:slot}). The bound at $2$ and $3$ itself for the families of $D3(3,\mathcal A)$ is not used: the bound needed in Step~3
of \cref{pf:linkN:famthm} is obtained for the cut family from \cref{lem:families:at23}, which holds for every family. By
\cref{lem:families:at23} the hypothesis of (b) holds for every parameter (\cref{rem:families:linkN}).
\end{remark}

\begin{remark}[Largeness]\label{pf:linkN:remlarge}
In \cref{pf:linkN:slot} the assumption that $\mathcal F$ is large for $(K,C_H)$ is not used: $\sfname{LinkN}$ holds for every
$\mathcal F\in D3(3,\mathcal A)\cap\sfname{FamP}(\eps)$ with $\eps\le\eps_1$.
\end{remark}

\begin{remark}[Conditional conclusions and abstract data]\label{pf:linkN:remcond}
$\sfname{LinkN}$ says nothing about members with $\linkNfin(D_{x,i})$ false or with $\sum_\ell\mathrm{vv}(D_{x,i},\ell)$ divergent. The
invariants $\eta_j(w)$, $\mathrm{sh}(\ell,e)$, $\#G$ and the sets $\linkNgood_\iota\subseteq Z_\iota$ are components of the datum, tied to
lattice indices and to lengths of Tate--Shafarevich groups only through the compatibility conditions \eqref{eq:pf:linkN:ties}; a
normalised length of a module of infinite length is $0$ by convention (\cref{pf:linkN:localinv}(b)), and \cref{pf:linkN:lengthle}
shows that the lengths entering $\linkNloc$ are finite. The main term $fmb\Lambda^2=fm\,e^{x^\eps}x^{2\eps}$ grows faster than any
power of $x$.
\end{remark}

\begin{remark}[Unused hypotheses]\label{pf:linkN:remunused}
The finiteness statement of \cref{pf:linkN:Hdec} is a hypothesis of \cref{pf:linkN:main} but is not used in its proof: it is passed
only to \cref{pf:linkN:tamevanish}, whose proof uses only duality and the Euler characteristic formula, and to \cref{pf:linkN:GPmem},
whose formal proof uses none of the three statements of \cref{pf:linkN:Hdec} (they are passed to the construction of the conditions
(GP1)--(GP12), which does not use them). One hypothesis of the formal version of \cref{pf:linkN:operator}, a tie between labels, is
not used either.
\end{remark}

\begin{remark}[Constants and names]\label{pf:linkN:remconst}
The constants produced are explicit: $C_T=40$ and $C=320$ in \cref{pf:linkN:extra}; $2000(|C|+1)^3$ in \cref{pf:linkN:sumthm}; and the
estimates of \cref{pf:linkN:pointbound} with the numerical constants $8$, $4$, $20$, $1000$, $1164000$ and $4\cdot10^6$. The formal
statement of \cref{pf:linkN:lenunit} is an inequality although equality holds. The names of several formal results (for instance
the one of \cref{pf:linkN:main}) do not describe their content, and the ``Theorem~T'' of \cref{pf:linkN:closinglocal} is an internal
name of the development, not a published theorem.
\end{remark}

\begin{remark}[The crystalline period bound]\label{pf:linkN:remthmT}
The lifting step \cref{pf:linkN:liftstep} is stated under the admitted local statements, but its proof does not use them; consequently
the hypotheses on the local Bloch--Kato theory, local duality and the Euler characteristic formula, and the comparison data, enter
\cref{pf:linkN:thmT} only through statements that are not used, and the conclusion of \cref{pf:linkN:thmT} does not involve the Hodge--Witt,
Fontaine and comparison data. The lifting theorem gives the exponent $7(Nr'+\eta(N+r'))$; the assumption $C\ge40$ is used for the
inequality $40(HN^2+N^3)<\ell$ required by it, and $7\le C$ for \cref{pf:linkN:liftmono}.
\end{remark}

\begin{remark}[Auxiliary statements]\label{pf:linkN:remaux}
In \cref{pf:linkN:hostsize}(ii) the formal statement is an inequality, although the proof gives equality. The hypothesis $1\in R'$ of
\cref{pf:linkN:lgsel} always holds. \Cref{pf:linkN:unramhom} assumes that the inertia restrictions vanish at every finite place, including the
places above $\ell$. In \cref{pf:linkN:countL} the same number $e$ is used at all places above $\ell$.
\end{remark}
